\documentclass[12pt, b5paper, oldfontcommands]{memoir}
% Margins
\usepackage[b5paper,total={6in, 9in}]{geometry}
\geometry{margin=2.4cm, top=2.5cm}

% Citations at the end of each chapter
\usepackage[sectionbib,square,comma]{natbib}
\usepackage{chapterbib}
\bibliographystyle{bib/genbibstyle}

% Required for LaTeXDraw
%\usepackage[usenames,dvipsnames]{pstricks}
%\usepackage{pstricks-add}
%\usepackage{epsfig}
%\usepackage{pst-grad} % For gradients
%\usepackage{pst-plot} % For axes

% Switched to tikz ecosystem
\usepackage{forest} % PGF/TikZ-based mechanism for drawing linguistic trees.
\usepackage{graphicx}
%\usepackage{svg} % For \includesvg
% Required for grammar trees
\usepackage{pst-node,pst-tree}

\usepackage[space]{grffile} % For spaces in paths
\usepackage{etoolbox} % For spaces in paths
\makeatletter % For spaces in paths
\patchcmd\Gread@eps{\@inputcheck#1 }{\@inputcheck"#1"\relax}{}{}
\makeatother



\usepackage{tgheros} % "TEX Gyre Heros", Helvetica, code: qhv
%\usepackage{tgbonum} % "TEX Gyre Bonum", ITCBookman, code: qbk
% Keep commented. Somehow ptm still loads, while uncommenting messes up Greek letters.
%\usepackage{mathptmx} % Times, code: ptm

\usepackage[utf8]{inputenc}
\usepackage[T1]{fontenc}
\usepackage[english]{babel}
\usepackage{anyfontsize}
\usepackage{microtype}
\usepackage{sfmath} % Sans Serif math fonts
\renewcommand{\rmdefault}{ptm} % Times
\renewcommand{\sfdefault}{qhv} % Helvetica
\renewcommand{\familydefault}{\rmdefault}

\usepackage[fleqn]{amsmath}
\usepackage{mathrsfs} % Math script
\usepackage{mathtools}
%\usepackage{textgreek} % Greek characters in text mode.
%\newcommand{\tl}{λ}
\newcommand{\tlb}[1]{$\boldsymbol{\lambda}$\,#1.\,}
\newcommand{\tl}{$\boldsymbol{\lambda}$}
\newcommand{\ta}{$\boldsymbol{\alpha}$}
\newcommand{\tb}{$\boldsymbol{\beta}$}
\newcommand{\te}{$\boldsymbol{\eta}$}
\newcommand{\td}{$\boldsymbol{\delta}$}
\newcommand{\tr}{$\boldsymbol{\rho}$}
\newcommand{\conversion}[1]{$\underset{\boldsymbol{#1}}{\leftrightarrow}$}
\newcommand{\reduction}[1]{$\underset{\boldsymbol{#1}}{\rightarrow}$}
\newcommand{\antireduction}[1]{$\underset{\boldsymbol{#1}}{\leftarrow}$}
%\DeclareUnicodeCharacter{03BB}{\lt}
%\usepackage{newunicodechar}
%\newunicodechar{λ}{\lt}

% For source code highlighting.
%\usepackage[no-math]{fontspec}
%\usepackage{minted}
%\usemintedstyle{mathematica}

%\usepackage{tabu} % Tables
%\usepackage{subcaption} % Subfigures (side-by-side)
%\usepackage{wrapfig}


% Get rid of those hideous arrows
\usepackage{chemarrow}
\renewcommand{\rightarrow}{\rarrowfill{1.2em}}
\renewcommand{\to}{\rarrowfill{1.2em}}
\renewcommand{\leftarrow}{\larrowfill{1em}}
\renewcommand{\leftrightarrow}{\larrowfill{1em}\hspace{-0.75em}\rarrowfill{1em}}
\newcommand{\xdownarrow}[1]{%
    {\left\downarrow\vbox to #1{}\right.\kern-\nulldelimiterspace}
}

% enumerations
\usepackage{paralist}
%\setlist{noitemsep,topsep=5pt,parsep=0pt,partopsep=5pt}

% Bibliography management
%\usepackage[backref=true,style=numeric,sorting=none]{biblatex}
%\bibliography{references}
%\usepackage[sectionbib,authoryear]{natbib}
%\usepackage{chapterbib}
\usepackage[hidelinks]{hyperref} % links (after biblatex)

%\usepackage{fullpage} %Comment this out before submission.
%Only used for DRAFTS to include the time of compilation as a versioning system
\usepackage{datetime}

\usepackage{arydshln} % Dashed line in table
\usepackage{dashrule} % \hdashrule

% Whitespace adjustments
\newcommand{\vs}{\vspace{.5\baselineskip}}
\newcommand{\pg}{\clearpage}

% We create a math mode for the Miranda code, because there is occasional need
% for mathematical operators.
%\usepackage{sansmath}
%\usepackage{mathastext}
%\DeclareMathVersion{sfmath}
%\DeclareSymbolFont{sfletters}{OT1}{qhv}{m}{up}
%\SetSymbolFont{letters}{sfmath}{OT1}{qhv}{m}{up}
%\DeclareSymbolFont{sfoperators}{OT1}{qhv}{m}{up}
%\SetSymbolFont{operators}{sfmath}{OT1}{qhv}{m}{up}
%\SetMathAlphabet\mathit{sfmath}{OT1}{qhv}{m}{sl}
%\SetMathAlphabet\mathrm{sfmath}{OT1}{qhv}{m}{n}

% This is the mlcoded environment, an environment in upright Helvetica.
\newenvironment{mlcoded}{
	% The following newline is necessary for some reason!

	\sffamily
	\footnotesize
	\spaceskip=2
	\fontdimen2
	\font
	\list{}{\rightmargin10pt\leftmargin10pt}\item[] % This line is how the quote environment is defined.
%	\noindent\par\vspace{0.5\baselineskip}
}{
	\endlist
%	\vspace{0.5\baselineskip}\noindent\\
}

\newenvironment{mlalign}{
	\begin{mlcoded}\setlength{\tabcolsep}{2pt}
%    \vspace{-.3cm}
	\begin{tabular}{rl}
}{
	\end{tabular}
%    \vspace{-.2cm}
	\end{mlcoded}
}

\newenvironment{letalign}{
    \begin{mlcoded}\setlength{\tabcolsep}{2pt}
        %    \vspace{-.3cm}
        \begin{tabular}{ll}
        }{
        \end{tabular}
        %    \vspace{-.2cm}
    \end{mlcoded}
}

\usepackage{environ}

\NewEnviron{mlnumbered}{%
    \begin{equation}
    \text{\hspace{-0.9em}\ml{\BODY}}%
    \end{equation}%
}{%
%what goes here?
}

% Inline Miranda code: "this is \ml{f x} applied to...."
%\newcommand{\ml}[1]{{\sansmath$#1$}}
\newcommand{\ml}[1]{{\sffamily\footnotesize #1}}

% For Roman Numerals
%\makeatletter
%\newcommand*{\rom}[1]{\expandafter\@slowromancap\romannumeral #1@}
%\makeatother

%%%%% Title page %%%%%
\makeatletter
\newcommand*{\maketitlepg}{{
		\null
		\let\cleardoublepage\clearpage
		\thispagestyle{empty}
		\leavevmode
		\normalfont
		\flushleft
		%	\parindent=0pt
		\vspace*{\drop}
		{\noindent\HUGE\fontfamily{qhv}\selectfont\uppercase \@title}\par
		\rule{\linewidth}{3pt}

		\vspace{2cm}
		{\LARGE\fontfamily{qhv}\selectfont\@author}\par
		%	\rule{\unitlength}{1.6pt}
		{\noindent\fontfamily{qhv}\selectfont\members\par\null}
		\vfill
		{\small\date{\today}}
		%	\vspace*{\drop}
		\cleardoublepage
}}
\makeatother

% This is how to do this in memoir:
\setsecnumdepth{subsubsection}
%\setcounter{secnumdepth}{4}
%\setcounter{tocdepth}{4}

\makeindex

%%%%% Other headings %%%%%

% Headings
\usepackage{fancyhdr}
\pagestyle{fancy}

\renewcommand{\chaptermark}[1]{\markboth{#1}{}}
\renewcommand{\sectionmark}[1]{\markright{#1}}
%\renewcommand{\subsectionmark}[1]{\markright{#1}}
\fancyhf{}
\fancyhead[LE,RO]{\footnotesize \thepage}
\fancyhead[LO]{\footnotesize \em Section \thesection\; \rightmark}
%\fancyhead[LO]{\fancyplain{}\slshape{\rightmark}}
\fancyhead[RE]{\footnotesize \em\chaptername \;\thechapter\; \leftmark}
%\fancyfoot[CE,CO]{\leftmark}
%\fancyfoot[LE,RO]{\thepage}

% For spelling out numbers: \numberstringnum{4} or \Numberstringnum{4}
\usepackage{fmtcount}
% For underlining the chapter number
\usepackage{calc}

\usepackage{caption}
\captionsetup{labelsep=quad,font={sf}}

% Section Formatting
\usepackage{titlesec}


% Dedication
\newcommand{\dedication}{To Dorothy}
% Dedication formatting, very similar to chapter and part formatting below
\newcommand{\dedicationstyle}[1]{{
		\normalfont\fontsize{20}{20}\selectfont \sffamily
		\makebox[2pt][l]{%\hspace-1.7cm
					\rule[-6pt]{\widthof{ttttttt\dedication}}{3pt}

		}
		%	\upshape\Large\sffamily
		\hspace{\widthof{ttttt}}
		\vspace{-5pt}
		#1
}}


% Chapter title formatting
\titleformat{\chapter}[display]{\normalfont\fontsize{30}{20}\selectfont\itshape \sffamily}{
   \makebox[2pt][l]{\hspace{-1.7cm}
	\rule[-6pt]{\widthof{ttttt\Numberstringnum{\thechapter}}}{3pt}
	}\Numberstringnum{\thechapter}
}{.7cm}
{
%	{\vspace{-1.1\baselineskip}\hrule height 3pt width 1.25in \relax}\vspace{\baselineskip}
	\upshape\Large\sffamily\MakeUppercase
}[]
\titlespacing*{\chapter}{ 1.7cm}{ 1cm}{1cm}  %[ right sep ]
\newcommand{\chapterbreak}{\clearpage\thispagestyle{empty}}

\newcommand{\chapterauthors}[1]{
{\upshape\large\sffamily
    \vspace{-1\baselineskip}
    \hspace{1.7cm}\noindent
    #1
}
}


% Part title formatting - identical to chapter formatting
\titleformat{\part}[display]{\normalfont\fontsize{30}{20}\selectfont\itshape \sffamily}{
	\makebox[2pt][l]{\hspace{-1.7cm}
		\rule[-6pt]{\widthof{tttttPart \thepart}}{3pt}
	}Part \thepart
}{.7cm}
{
	%	{\vspace{-1.1\baselineskip}\hrule height 3pt width 1.25in \relax}\vspace{\baselineskip}
	\upshape\LARGE\sffamily\MakeUppercase
}[]
\titlespacing*{\part}{ 1.7cm}{ 1cm}{3cm}  %[ right sep ]
\titleclass{\part}{top}
\newcommand{\partbreak}{\clearpage\thispagestyle{empty}}
% Section heading spacing
% \titleformat{<command>}[<shape>]{<format>}{<label>}{<sec>}{<before-code>}[<after-code>]
\titleformat{\section}{\bf\sffamily}{\thesection}{8pt}{}
\titleformat{\subsection}{\sffamily}{\thesubsection}{8pt}{}
%\titleformat{\subsubsection}{\large}{\textsf{\thesubsubsection}}{8pt}{\textsf}

% Removes page numbers on first part/chapter/title pages. Since it clobbers the plain page
% style, we can't use that style anywhere.
\makeatletter
\let\ps@plain\ps@empty
\makeatother

\title{{The Implementation\\ of Functional\\ \vspace{10pt}Programming Languages}}

%\subtitle{ }

\author{Simon L. Peyton Jones}

\newcommand{\members}{\leavevmode Department of Computer Science\\ University College London\par
	\vspace*{\baselineskip}
	\textit{with chapters by}\\
	Philip Wadler, Programming Research Group, Oxford\\
	Peter Hancock, Metier Management Systems Ltd\\
	David Turner, University of Kent, Canterbury
}

\newlength{\drop}
\drop=0.1\textheight


\newenvironment{numbered}{
	\vs
	\begin{compactenum}[(i)]
}{
	\end{compactenum}
	\vs
}

\newenvironment{references}{\begin{description}\itemsep -5pt \small
}{\end{description}}

% Boxed asides.
\usepackage{float}
\newcommand{\boxedfigure}[2]{
	\begin{figure}[H]
		\centering

		{%
			\setlength{\fboxrule}{1pt}%
			\setlength{\fboxsep}{10pt}%
			\fbox{%
				\begin{minipage}{0.9\textwidth}
					\small
					\setlength{\parindent}{10pt}
					\setlength{\parskip}{0mm plus 0mm minus 0mm}
					#1
			\end{minipage}%
			}%
		}%

		\caption{\textsf #2}
	\end{figure}
}


\newcommand{\plainbox}[1]{
    {\centering

        {%
            \vspace{0.5\baselineskip}
            \setlength{\fboxrule}{1pt}%
            \setlength{\fboxsep}{10pt}%
            \fbox{%
                \begin{minipage}{0.9\textwidth}
                    \small
                    \setlength{\parindent}{10pt}
                    \setlength{\parskip}{0mm plus 0mm minus 0mm}
                    #1
                \end{minipage}%
            }%
            \vspace{0.5\baselineskip}
        }%

    }
}


\newcommand{\simpletitledbox}[2]{%
	\plainbox{%
		{\centering

			#1

		}
		\vspace{0.4\baselineskip}

		\noindent #2
	}%
}%


\newcommand{\titledbox}[3]{%
\plainbox{%
{\centering

#1

}
\vspace{0.4\baselineskip}

\noindent\textit{#2}
\vspace{0.4\baselineskip}

\noindent #3
}%
}%

\newcommand{\hdashsep}{\vbox{\noindent\hdashrule[0pt][x]{\fill}{0.5pt}{5pt 2pt}}}

\newcommand{\theorembox}[2]{\titledbox{THEOREM}{\noindent #1}{\noindent #2}}
\newcommand{\definitionbox}[2]{\titledbox{DEFINITION}{\noindent #1}{\noindent #2}}

\newenvironment{pcorollary}{\begin{quote}\textit{Corollary.\;}}{\end{quote}}
\newenvironment{pproof}{\begin{quote}\textit{Proof.\;}}{\end{quote}}

\def\bracketadjust{\hspace{-1.2pt}}
\newcommand{\doublebracket}[1]{%
    \ml{\mbox{\bfseries[\bracketadjust[\ }}%
    \ml{#1}%
    \ml{\mbox{\bfseries\ ]\bracketadjust]}}%
}%

\newcommand{\metafn}[1]{\ml{\bfseries #1}}
\newcommand{\metafnbb}[2]{\metafn{#1}\doublebracket{#2}}


% Infix box operator
\newcommand{\fatbar}{{\sffamily [\hspace{-1pt}]}}

% The `::=` operator
\newcommand{\typedecl}{\ml{$::=$}}
\newcommand{\hastype}{\ml{$::$}}
% The equivalence operator with spacing for use in ml mode
\newcommand{\equivalent}{\quad $\equiv$\quad}

\newenvironment{parsetree}[1]{\begin{center}\begin{forest}[#1]}{\end{forest}\end{center}}

\begin{document}

\frontmatter

\maketitlepg

% Make Dedication
\begin{flushleft}
	\thispagestyle{empty}
	\vspace*{\drop}
%	\huge
%	\dedication
%	\dedicationstyle{\dedication}
\end{flushleft}
\clearpage

\tableofcontents
%\listoffigures
%\listoftables

\setlength{\parindent}{10pt}
\setlength{\parskip}{0mm plus 0mm minus 0mm}

\chapter[Preface]{}

% The Preface "chapter" is typeset as if the word preface were the number instead of the title,
% except it is in all caps as well.
\vspace{-3.5cm}
\makebox[2pt][l]{
{\normalfont\fontsize{20}{20}\selectfont\sffamily
\hspace{0.75cm}
\makebox[2pt][l]{ \hspace{-1.5cm}
	\rule[-6pt]{\widthof{mmPREFACE}}{3pt}
}PREFACE}
}

\vspace{1.5cm}

\noindent This book is about implementing functional programming languages using
graph reduction.

Functional languages have become the focus of much active research in
recent years
%\citep{backus_can_1978} \citep{peyton-jones_directions_1984}
[Backus, 1978] [Peyton Jones, 1984], but their acceptance has
been delayed by the inefficiency of their available implementations when
compared with more conventional languages.

This situation has changed recently, with the advent of rather fast implementations of functional languages such as Cardelli's ML [Cardelli, 1983],
Fairbairn's Ponder [Fairbairn, 1982], and the Chalmers Lazy ML compiler
[Johnsson, 1984]. These implementations rival the speed of compilers for
more conventional languages.

There appear to be two main approaches to the efficient implementation of
functional languages. The first is an environment-based scheme, exemplified
by Cardelli's ML implementation, which derives from the experience of the
Lisp community. The other is graph reduction, a much newer technique first
invented by Wadsworth [Wadsworth, 1971], and on which the Ponder and
Lazy ML implementations are founded. Despite the radical differences in
beginnings, the most sophisticated examples of each approach show remark-
able similarities.

The techniques of graph reduction are to be found scattered amongst the
proceedings of various conferences and workshops, and it is one purpose of
this book to collect some of this work together. The book is intended to have
two main applications:
\begin{compactenum}[(i)]
\item As a course text for part of an undergraduate or postgraduate course on
the implementation of functional languages.
\item As a handbook for those attempting to write a functional language
implementation based on graph reduction.
\end{compactenum}
The material is presented in a fairly informal tutorial fashion, the idea being to
convey the framework and some of the intuitions that will render the original
sources more accessible.

Chapters 5 and 7 were written by Philip Wadler, of the Programming
Research Group, Oxford, and Chapter 4 was written in collaboration with
him. Chapters 8 and 9 were written by Peter Hancock, of Metier Management
Systems Ltd, and currently at the Programming Research Group, Oxford. I
gratefully acknowledge their patience in writing and rewriting their drafts.

I am extremely grateful to a number of other people who have made
significant technical contributions to the book. David Turner's help was
invaluable, in offering comments on the parts of the book that relate to
Miranda. The Appendix which gives an introduction to Miranda, was written
by him. (Miranda is a trademark of Research Software Limited.)

Much of the information about Miranda in the first part of the book is based
on a prerelease version of the Miranda system, and I am grateful to Research
Software Limited for the permission to include this material.

Simon Finn, of the University of Stirling, made a number of penetrating
observations about the treatment of pattern-matching, which resulted in
significant improvements in Chapters 4 and 6.

Several long discussions with Thomas Johnsson of Chalmers University,
Goteborg, radically changed the shape of the G-machine chapters, and
Chapter 21 was written in collaboration with him. John Fairbairn and Stuart
Wray of Cambridge also made important contributions in this area.

Paul Hudak and Robert Keller helped me in writing the last section of
Chapter 24.

Many other people have given me welcome help and encouragement, and
have helped enormously by reading the text and making constructive
suggestions. These include Lennart Augustsson, Philip Boswell, Geoff Burn,
Nigel Chapman, Chris Clack, Pierre-Louis Curien, Dan Friedman, Kevin
Hammond, Peter Hancock, Chris Hankin, John Hughes, Thomas Johnsson,
Dick Kieburtz, Phil Molyneux, Benedict Nixon, Martin Packer, Ellen Poon,
Jon Salkild, William Stoye, David Turner, Phil Wadler, John Washbrook and
Russel Winder. I am particularly grateful to John Washbrook for his unfailing
support.

\vspace{-\baselineskip}
\begin{flushright}
Simon L. Peyton Jones\\
University College London\\
London WC1E 6BT\\
simonpj@uk.ac.ucl.cs
\end{flushright}

%\pg
\vspace{-2\baselineskip}

%\bibliography{bib/SPLJ}

\section*{References}

\begin{references}
\item Backus, J. 1978. Can programming be liberated from the von Neumann style? A
functional style and its algebra of programs. \textit{Communications of the ACM}. Vol.~21, no.~8, pp. 613-41.
\item Cardelli, L. 1983. The functional abstract machine. \textit{Polymorphism}. Vol.~1,~no.~1.
\item Fairbairn, J. 1982. Ponder and its type system. \textit{Technical Report 31}. Computer Lab.,
Cambridge. November.
\item Johnsson, T. 1984. Efficient compilation of lazy evaluation. In \textit{Proceedings of the ACM
Conference on Compiler Construction, Montreal}, pp. 58-69. June.
\item Peyton Jones, S.L. 1984. Directions in functional programming research. In SERC
\textit{Distributed Computing Systems}, pp. 220-49. Duce (editor). Peter Peregrinus.
\item Wadsworth, C.P. 1971. Semantics and pragmatics of the lambda calculus, Chapter 4.
PhD thesis, Oxford.

\end{references}


\mainmatter


\chapter{Introduction}

This book is about implementing functional programming languages using
\textit{lazy graph reduction}, and it divides into three parts.

The first part describes how to translate a high-level functional language
into an intermediate language, called the lambda calculus, including detailed
coverage of pattern-matching and type-checking. The second part begins with
a simple implementation of the lambda calculus, based on graph reduction,
and then develops a number of refinements and alternatives, such as super-
combinators, full laziness and SK combinators. Finally, the third part
describes the G-machine, a sophisticated implementation of graph reduction,
which provides a dramatic increase in performance over the implementations
described earlier.

One of the agreed advantages of functional languages is their semantic
simplicity. This simplicity has considerable payoffs in the book. Over and
over again we are able to make semi-formal arguments for the correctness of
the compilation algorithms, and the whole book has a distinctly mathematical
flavor - an unusual feature in a book about implementations.

Most of the material to be presented has appeared in the published
literature in some form (though some has not), but mainly in the form of
conference proceedings and isolated papers. References to this work appear
at the end of each chapter.

\section{Assumptions}
This book is about implementations, not languages, so we shall make no
attempt to extol the virtues of functional languages or the functional
programming style. Instead we shall assume that the reader is familiar with
functional programming; those without this familiarity may find it heavy
going. A brief introduction to functional programming may be found in
Darlington [1984], while Henderson [1980] and Glaser et al. [1984] give more
substantial treatments. Another useful text is Abelson and Sussman [1985]
which describes Scheme, an almost-functional dialect of Lisp.

An encouraging consensus seems to be emerging in the basic features of
high-level functional programming languages, exemplified by languages such
as SASL [Turner, 1976], ML [Gordon et al., 1979], KRC [Turner, 1982],
Hope [Burstall et al., 1980], Ponder [Fairbairn, 1985], LML [Augustsson,
1984], Miranda [Turner, 1985] and Orwell [Wadler, 1985]. However, for the
sake of definiteness, we use the language Miranda as a concrete example
throughout the book (When used as the name of a programming language,
`Miranda' is a trademark of Research Software Limited.) A brief intro-
duction to Miranda may be found in the appendix, but no serious attempt is
made to give a tutorial about functional programming in general, or Miranda
in particular. For those familiar with functional programming, however, no
difficulties should arise.

Generally speaking, all the material of the book should apply to the other
functional languages mentioned, with only syntactic changes. The only
exception to this is that we concern ourselves almost exclusively with the
implementation of languages with \textit{non-strict semantics} (such as SASL, KRC,
Ponder, LML, Miranda and Orwell). The advantages and disadvantages of
this are discussed in Chapter 11, but it seems that graph reduction is probably
less attractive than the environment-based approach for the implementation
of languages with strict semantics; hence the focus on non-strict languages.
However, some functional languages are strict (ML and Hope, for example),
and while much of the book is still relevant to strict languages, some of the
material would need to be interpreted with care.

Generally speaking, all the material of the book should apply to the other
functional languages mentioned, with only syntactic changes. The only
exception to this is that we concern ourselves almost exclusively with the
implementation of languages with \textit{non-strict semantics} (such as SASL, KRC,
Ponder, LML, Miranda and Orwell). The advantages and disadvantages of
this are discussed in Chapter 11, but it seems that graph reduction is probably
less attractive than the environment-based approach for the implementation
of languages with strict semantics; hence the focus on non-strict languages.
However, some functional languages are strict (ML and Hope, for example),
and while much of the book is still relevant to strict languages, some of the
material would need to be interpreted with care.

The emphasis throughout is on an informal approach, aimed at developing
understanding rather than at formal rigor. It would be an interesting task to
rewrite the book in a formal way, giving watertight proofs of correctness at
each stage.

\section{Part I: Compiling High-level Functional Languages}
It has been widely observed that most functional languages are quite similar to
each other, and differ more in their syntax than their semantics. In order to
simplify our thinking about implementations, the first part of this book shows
how to translate a high-level functional program into an \textit{intermediate language}
which has a very simple syntax and semantics. Then, in the second and third
parts of the book, we will show how to implement this intermediate language
using graph reduction. Proceeding in this way allows us to describe graph
reduction in considerable detail, but in a way that is not specific to any
particular high-level language.

The intermediate language into which we will translate the high-level
functional program is the notation of the lambda calculus (Figure~\ref{fig:implfunctprog}). The
lambda calculus is an extremely well-studied language, and we give an introduction
to it in Chapter 2.


\boxedfigure{
\centering

\framebox{\
    \begin{minipage}{6cm}
        \centering

        High-level language program\
    \end{minipage}
}\\
\hspace{4cm}$\Bigg\downarrow$ \quad
\begin{minipage}{4cm}
    {Part I}
\end{minipage}\\
\framebox{\
    \begin{minipage}{6cm}
        \centering

        Program expressed in lambda notation
    \end{minipage}
    \ }\\
\hspace{4cm}$\Bigg\downarrow$  \quad
\begin{minipage}{4cm}
    {Parts II and III}
\end{minipage}\\
\framebox{\
    \begin{minipage}{6cm}
        \centering

        Concrete implementation
    \end{minipage}
    \ }\\

}{Implementing a functional program}



The lambda calculus is not only simple, it is also sufficiently expressive to
allow us to translate any high-level functional language into it. However,
translating some high-level language constructs into the lambda notation is
less straightforward than it at first appears, and the rest of Part I is concerned
with this translation.

Part I is organized as follows. First of all, in Chapter 3, we define a language
which is a superset of the lambda calculus, which we call the enriched lambda
calculus. The extra constructs provided by the enriched lambda calculus are
specifically designed to allow a straightforward translation of a Miranda
program into an expression in the enriched lambda calculus, and Chapter 3
shows how to perform this translation for simple Miranda programs.

After a brief introduction to pattern-matching, Chapter 4 then extends the
translation algorithm to cover more complex Miranda programs, and gives a
formal semantics for pattern-matching. Subsequently, Chapter 7 rounds out
the picture, by showing how Miranda's ZF expressions can also be translated
in the same way. (Various advanced features of Miranda are not covered,
such as algebraic types with laws, abstract data types, and modules.)

Much of the rest of Part I concerns the transformation of enriched lambda
calculus expressions into the ordinary lambda calculus subset, a process which
is quite independent of Miranda. This language-independence was one of the
reasons for defining the enriched lambda calculus language in the first place.
Chapter 5 shows how expressions involving pattern-matching constructs may
be transformed to use case-expressions, with a considerable gain in efficiency.
Then Chapter 6 shows how all the constructs of the enriched lambda calculus,
including case-expressions, may be transformed into the ordinary lambda
calculus.

Part I concludes with Chapter 8 which discusses type-checking in general,
and Chapter 9 in which a type-checker is constructed in Miranda.

\section{Part II: Graph Reduction}

The rest of the book describes how the lambda calculus may be implemented
using a technique called \textit{graph reduction}. It is largely independent of the later
chapters in Part I, Chapters 2--4 being the essential prerequisites.

As a foretaste of things to come, we offer the following brief introduction to
graph reduction. Suppose that the function \ml{f} is defined (in Miranda) like this:
\begin{mlcoded}
f\, x = (x + 1) $*$ (x -- 1)
\end{mlcoded}
This definition specifies that \ml{f} is a function of a single argument \ml{x}, which
computes `\ml{(x + 1) $*$ (x -- 1)}'. Now suppose that we are required to evaluate
\begin{mlcoded}
f\, 4
\end{mlcoded}
that is, the function \ml{f} applied to \ml{4}. We can think of the program like this:
\begin{center}
    \begin{forest}
        [\ml{@}
            [\ml{f}]
            [\ml{4}]
        ]
    \end{forest}
\end{center}
where the \ml{@} stands for function application. Applying \ml{f} to \ml{4} gives
\begin{center}
    \begin{forest}
        [\ml{$*$}
            [$+$
                [\ml{4}]
                [\ml{1}]
            ]
            [$-$
                [\ml{4}]
                [\ml{1}]
            ]
        ]
    \end{forest}
\end{center}
(Note: in the main text we will use a slightly different representation for
applications of \ml{$*$}, \ml{$+$} and \ml{$-$}, but this fact is not significant here.) We may now
execute the addition and the subtraction (in either order), giving
\begin{center}
    \begin{forest}
    [\ml{$*$}
        [\ml{5}]
        [\ml{3}]
    ]
\end{forest}
\end{center}
Finally we can execute the multiplication, to give the result
\begin{center}
\ml{15}
\end{center}

From this simple example we can see that:
\begin{numbered}
\item \textit{Executing} a functional program consists of \textit{evaluating} an expression.
\item A functional program has a natural representation as a \textit{tree} (or, more
generally, a \textit{graph}).
\item Evaluation proceeds by means of a sequence of simple steps, called
\textit{reductions}. Each reduction performs a local transformation of the graph
(hence the term \textit{graph reduction}).
\item Reductions may safely take place in a variety of orders, or indeed in
parallel, since they cannot interfere with each other.
\item Evaluation is complete when there are no further reducible expressions.
\end{numbered}
Graph reduction gives an appealingly simple and elegant model for the
execution of a functional program, and one that is radically different from the
execution model of a conventional imperative language.

We begin in Chapter 10 by discussing the representation of a functional
program as a graph. The next two chapters form a pair which discusses first the
question of deciding which reduction to perform next (Chapter 11), and then
the act of performing the reduction (Chapter 12).

Chapters 13 and 14 introduce the powerful technique of \textit{supercombinators},
which is the key to the remainder of the book. This is followed in Chapter 15
with a discussion of \textit{full laziness}, an aspect of lazy evaluation; this chapter can
be omitted on first reading since later material does not depend on it.

Chapter 16 then presents \textit{SK combinators}, an alternative implementation
technique to supercombinators. Hence, this chapter can be understood
independently of Chapters 13--15. Thereafter, however, we concentrate on
supercombinator-based implementations.

Part II concludes with a chapter on \textit{garbage collection}.

\section{Part Ill: Advanced Graph Reduction}

It may seem at first that graph reduction is inherently less efficient than more
conventional execution models, at least for conventional von Neumann
machines. The bulk of Part III is devoted to an extended discussion of the
G-machine, which shows how graph reduction can be compiled to a form that
is suitable for \textit{direct execution} by ordinary sequential computers.

In view of the radical difference between graph reduction on the one hand,
and the linear sequence of instructions executed by conventional machines on
the other, this may seem a somewhat surprising achievement. This (fairly
recent) development is responsible for a dramatic improvement in the speed
of functional language implementations.

Chapters 18 and 19 introduce the main concepts of the G-machine, while
Chapters 20 and 21 are devoted entirely to optimizations of the approach.
The book concludes with three chapters that fill in some gaps, and offer
some pointers to the future.

Chapter 22 introduces \textit{strictness analysis}, a compile-time program analysis
method which has been the subject of much recent work, and which is crucial
to many of the optimizations of the G-machine.

Perhaps the major shortcoming of functional programming languages,
from the point of view of the programmer, is the difficulty of estimating the
space and time complexity of the program. This question is intimately bound
up with the implementation, and we discuss the matter in Chapter 23.
graph reduction.

Finally, the book concludes with a chapter on parallel implementations of
graph reduction.

\pg

\section*{References}

\begin{references}
	\item Abelson, H., and Sussman, G.J. 1985. \textit{Structure and Interpretation of Computer
	Programs}. MIT Press.
	\item  Augustsson, L. 1984. A compiler for lazy ML. \textit{Proceedings of the ACM Symposium on
	Lisp and Functional Programming, Austin}. August, pp. 218--27.
	\item Burstall, R.M., MacQueen, D.B., and Sanella, D.T. 1980. \textit{Hope: an experimental
	applicative language}. CSR-62--80. Department of Computer Science, University of
	Edinburgh. May.
	\item  Darlington, J. 1984. Functional programming. In \textit{Distributed Computing}. Duce
(Editor). Academic Press.
	\item  Fairbairn, J. 1985. Design and implementation of a simple typed language based on the
lambda calculus. PhD thesis, \textit{Technical Report 75}. University of Cambridge. May.
	\item  Glaser, H., Hankin, C., and Till, D. 1984. \textit{Principles of Functional Programming}.
Prentice-Hall.
	\item  Gordon, M.J., Milner, A.J., and Wadsworth, C.P. 1979. \textit{Edinburgh LCF}. LNCS 78.
Springer Verlag.
	\item  Henderson, P. 1980. \textit{Functional Programming}. Prentice-Hall.
	\item  Turner, D.A. 1976. \textit{The SASL language manual}. University of St Andrews. December.
	\item  Turner, D.A. 1982. Recursion equations as a programming language. In \textit{Functional
Programming and Its Applications}, Darlington et al. (editors), pp. 1-28. Cam-
bridge University Press.
	\item  Turner, D.A. 1985. Miranda -- a non-strict functional language with polymorphic
types. In \textit{Conference on Functional Programming Languages and Computer
Architecture, Nancy}, pp. 1-16. Jouannaud (editor), LNCS 201. Springer Verlag.
	\item  Wadler, P. 1985. \textit{Introduction to Orwell}. Programming Research Group, University of
Oxford.
\end{references}


\part[Compiling High-Level Functional Languages]{Compiling High-Level\\ Functional Languages}

\chapter{The Lambda Calculus}

This chapter introduces the lambda calculus, a simple language which will be
used throughout the rest of the book as a bridge between high-level functional
languages and their low-level implementations. The reasons for introducing
the lambda calculus as an intermediate language are:
\begin{numbered}
\item It is a simple language, with only a few, syntactic constructs, and simple
semantics. These properties make it a good basis for a discussion of
implementations, because an implementation of the lambda calculus only
has to support a few constructs, and the simple semantics allows us to
reason about the correctness of the implementation.
\item It is an expressive language, which is sufficiently powerful to express all
functional programs (and indeed, all computable functions). This means
that if we have an implementation of the lambda calculus, we can
implement any other functional language by translating it into the lambda
calculus.
\end{numbered}
In this chapter we focus on the syntax and semantics of the lambda calculus
itself, before turning our attention to high-level functional languages in the
next chapter.

\section{The Syntax of the Lambda Calculus}

Here is a simple expression in the lambda calculus:
\begin{mlcoded}
(+ 4 5)
\end{mlcoded}
All function applications in the lambda calculus are written in \textit{prefix} form, so,
for example, the function + precedes its arguments 4 and 5. A slightly more
complex example, showing the (quite conventional) use of brackets, is
\begin{mlcoded}
(+ ($*$ 5 6) (+ 8 3))
\end{mlcoded}
In both examples, the outermost brackets are redundant, but have been
added for clarity (see Section 2.1.2).

From the implementation viewpoint, a functional program should be
thought of as an \textit{expression}, which is `executed' by \textit{evaluating} it. Evaluation
proceeds by repeatedly selecting a \textit{reducible expression} (or \textit{redex}) and
reducing it. In our last example there are two redexes: ($*$ 5 6) and (+ 8 3).
The whole expression \ml{(+ ($*$ 5 6) ($*$ 8 3))} is not a redex, since a \ml{+} needs to be
applied to two \textit{numbers} before it is reducible. Arbitrarily choosing the first
redex for reduction, we write
\begin{mlcoded}
(+ ($*$ 5 6) ($*$ 8 3)) $\to$ (+ 30 ($*$ 8 3))
\end{mlcoded}
where the \ml{$\to$} is pronounced `reduces to'. Now there is only one redex, \ml{(* 8 3)},
which gives
\begin{mlcoded}
(+ 30 ($*$ 8 3)) $\to$ (+ 30 24)
\end{mlcoded}
This reduction creates a new redex, which we now reduce
\begin{mlcoded}
(+ 30 24) $\to$ 54
\end{mlcoded}
\indent When there are several redexes we have a choice of which one to reduce
first. This issue will be addressed later in this chapter.

\subsection{Function Application and Currying}
In the lambda calculus, function application is so important that it is denoted
by simple juxtaposition; thus we write
\begin{mlcoded}
f x
\end{mlcoded}
to denote `the function \ml{f} applied to the argument \ml{x}'. How should we express
the application of a function to several arguments? We could use a new
notation, like \ml{(f (x,y))}, but instead we use a simple and rather ingenious
alternative. To express `the sum of 3 and 4' we write
\begin{mlcoded}
	((+ 3) 4)
\end{mlcoded}
The expression \ml{(+ 3)} denotes the function that adds \ml{3} to its argument. Thus
the whole expression means `the function \ml{+} applied to the argument \ml{3}, the
result of which is a function applied to \ml{4}'. (In common with all functional
programming languages, the lambda calculus allows a function to return a
function as its result.)

This device allows us to think of all functions as having a \textit{single argument
only}. It was introduced by Schonfinkel [1924] and extensively used by Curry
[Curry and Feys, 1958]; as a result it is known as \textit{currying}.

\subsection{Use of Brackets}

In mathematics it is conventional to omit redundant brackets to avoid
cluttering up expressions. For example, we might omit brackets from the
expression
\begin{mlcoded}
	(ab) + ((2c)/d)
\end{mlcoded}
to give
\begin{mlcoded}
	ab + 2c/d
\end{mlcoded}
The second expression is easier to read than the first, but there is a danger that
it may be ambiguous. It is rendered unambiguous by establishing conventions
about the precedence of the various functions (for example, multiplication
binds more tightly than addition)..

Sometimes brackets cannot be omitted, as in the expression:
\begin{mlcoded}
	(b + c)/a
\end{mlcoded}
Similar conventions are useful when writing down expressions in the
lambda calculus. Consider the expression:
\begin{mlcoded}
	((+ 3) 2)
\end{mlcoded}
By establishing the convention that function application associates to the left,
we can write the expression more simply as:
\begin{mlcoded}
	(+ 3 2)
\end{mlcoded}
or even
\begin{mlcoded}
	+ 3 2
\end{mlcoded}

We performed some such abbreviations in the examples given earlier. As a
more complicated example, the expression:
\begin{mlcoded}
	((f ((+ 4) 3)) (g x))
\end{mlcoded}
is fully bracketed and unambiguous. Following our convention, we may omit
redundant brackets to make the expression easier to read, giving:
\begin{mlcoded}
	f (+ 4 3) (g x)
\end{mlcoded}
No further brackets can be omitted. Extra brackets may, of course, be
inserted freely without changing the meaning of the expression; for example
\begin{mlcoded}
	(f (+ 4 3) (g x))
\end{mlcoded}
is the same expression again.

\subsection{Built-in Functions and Constants}
In its purest form the lambda calculus does not have built-in functions such as
\ml{+}, but our intentions are practical and so we extend the pure lambda calculus
with a suitable collection of such built-in functions.

These include arithmetic functions (such as \ml{+}, \ml{--}, $*$, \ml{/}\,) and constants (\ml{0}, \ml{1},
\ldots), logical functions (such as \ml{AND}, \ml{OR}, \ml{NOT}) and constants (\ml{TRUE},
\ml{FALSE}), and character constants (\ml{`a'}, \ml{`b'}, \ldots). For example
\begin{mlalign}
	-- 5 4 &$\to$ 1\\
AND TRUE FALSE &$\to$ FALSE
\end{mlalign}
We also include a conditional function, \ml{IF}, whose behavior is described by the
reduction rules:
\begin{mlcoded}
IF TRUE E\textsubscript{t} E\textsubscript{f} $\to$ E\textsubscript{t}\\
IF FALSE E\textsubscript{t} E\textsubscript{f} $\to$ E\textsubscript{f}
\end{mlcoded}

We will initially introduce data constructors into the lambda calculus by
using the built-in functions \ml{CONS} (short for \ml{CONSTRUCT}\,), \ml{HEAD} and \ml{TAIL}
(which behave exactly like the Lisp functions \ml{CONS}, \ml{CAR} and \ml{CDR}\,). The
constructor \ml{CONS} builds a compound object which can be taken apart with
\ml{HEAD} and \ml{TAIL}. We may describe their operation by the following rules:
\begin{mlalign}
HEAD (CONS a b) &$\to$ a\\
TAIL (CONS a b) &$\to$ b
\end{mlalign}
We also include \ml{NIL}, the empty list, as a constant. The data constructors will
be discussed at greater length in Chapter 4.

The exact choice of built-in functions is, of course, somewhat arbitrary, and
further ones will be added as the need arises.

\subsection{Lambda Abstractions}

The only functions introduced so far have been the built-in functions (such as
\ml{+} and \ml{CONS}\,). However, the lambda calculus provides a construct, called a
lambda abstraction, to denote new (non-built-in) functions. A lambda
abstraction is a particular sort of expression which denotes a function. Here is
an example of a lambda abstraction:
\begin{mlcoded}
	(\tlb{x}+ x 1)
\end{mlcoded}
The \tl says `here comes a function', and is immediately followed by a variable,
\ml{x} in this case; then comes a . followed by the \textit{body} of the function, \ml{(+ x 1)} in
this case. The variable is called the \textit{formal parameter}, and we say that the \tl
binds it. You can think of it like this:
\begin{quote}\setlength{\tabcolsep}{3pt}
\begin{tabular}{llllll}
	\ml{(\tl}	&\ml{x}  &\ml{.}  &\ml{+}  &\ml{x} &\ml{1}\,)  \\
	\;\;$\uparrow$	&$\uparrow$  &$\uparrow$  &$\uparrow$  &$\uparrow$ &$\uparrow$ \\
	That function of &\ml{x}  & which  &adds  &\ml{x} to &\ml{1} 
\end{tabular}
\end{quote}
A lambda abstraction \textit{always} consists of all the four parts mentioned: the \tl,
the formal parameter, the \ml{.} and the body.

A lambda abstraction is rather similar to a function definition in a
conventional language, such as C:
\begin{mlcoded}
	Inc( x )\\
	int x;\\
	\{ return( x + 1 ); \}
\end{mlcoded}
The formal parameter of the lambda abstraction corresponds to the formal
parameter of the function, and the body of the abstraction is an expression
rather than a sequence of commands. However, functions in conventional
languages must have a name (such as \ml{Inc}), whereas lambda abstractions are
`anonymous' functions.

The body of a lambda abstraction extends \textit{as far to the right as possible}, so
that in the expression
\begin{mlcoded}
	(\tlb{x}+ x 1) 4
\end{mlcoded}
the body of the \ml{\tl{}x} abstraction is \ml{(+ x 1)}, not just \ml{+}. As usual, we may add
extra brackets to clarify, thus
\begin{mlcoded}
	(\tlb{x}(+ x 1)) 4
\end{mlcoded}
When a lambda abstraction appears in isolation we may write it without any
brackets:
\begin{mlcoded}
	\tlb{x}+ x 1
\end{mlcoded}

\subsection{Summary}
We define a \textit{lambda expression} to be an expression in the lambda calculus, and
Figure 2. 1 summarizes the forms which a lambda expression may take. Notice
that a lambda \textit{abstraction} is not the same as a lambda \textit{expression}; in fact the
former is a particular instance of the latter.

<exp> :: = <constant>
Built-in constants
<variable>
Variable names
<exp> <exp>
A <variable>.<exp>
Applications
Lambda abstractions
This is the abstract syntax of lambda expressions. In order to write down
such an expression in concrete form we use brackets to disambiguate its
structure (see Section 2.1.2).
We will use lower-case letters for variables (e.g. x, f), and upper-case
letters to stand for whole lambda expressions (e.g. M, E).
The choice of constants is rather arbitrary; we assume integers and
booleans (e.g. 4, TRUE), together with built-in functions to manipulate
them (e.g. AND, IF, +). We also assume built-in list-processing functions
(e.g. CONS, HEAD).
Figure 2.1 Syntax of a lambda expression (in BNF)
\chapter[Translating a high-level functional language into the lambda calculus][Translating a high-level functional language into the lambda calculus]{Translating a high-level\\ functional language into the\\ lambda calculus}

In the next few chapters we will describe how to translate a high-level functional language into the lambda calculus. We can regard this translation in two ways:

\begin{numbered}
    \item As a description of the semantics of the language, giving the meaning of each of its constructs in terms of lambda expressions, whose meaning is well understood. This is precisely the approach taken by \textit{denotational semantics} [Gordon 1979].
    \item As a step in the implementation of the high-level language, by expressing all its constructs in terms of the lambda notation.
\end{numbered}

For the sake of definiteness we use a subset of the language Miranda [Turner 1985], but the techniques apply to any functional language. An introduction to Miranda can be found in the Appendix.
\vfill

\plainbox{
    \footnotesize
    {\centering

    \textbf{Disclaimer:}

    }
    \noindent In this book Miranda is used as an example of a modern functional programming language, to illustrate various points about the implementation of functional programming languages in general. This book is not intended to be a source of reference for the definition of Miranda. Note that:
    \begin{numbered}
        \item Miranda has a number of features, both major and minor, which are not discussed here at all.
        \item The material about Miranda in this book was based on a prerelease version of the Miranda system and may therefore be inaccurate by the time it is published.
    \end{numbered}

    The Miranda functional programming system is a product of Research Software Limited, and a full description of the language and its programming environment is in preparation by them.
}

\section{The Overall Structure of the Translation Process}
Miranda is a powerful, high-level functional language, providing a rich set of programming constructs. The purpose of the next few chapters is to demonstrate how some of these constructs can be translated into the lambda calculus. Specifically, we will discuss structured data types, pattern-matching, conditional equations and ZF expressions. Miranda includes a number of other constructs, such as abstract data types and structured data types with laws, which we will not study in this book.

Even so, the translation we describe is a substantial task, and we begin by outlining the structure of the translation process. It might be possible to translate a program directly from Miranda into the lambda calculus, but this would be an extremely complicated translation, so we will take a more step-by-step approach. In order to do this, it is convenient to regard much of the translation as a process of successively \textit{transforming} one program into another, until finally the result is a program in the lambda notation. (We are here using `translation' to suggest a process which takes a program in one language and produces a program in another, while a `transformation' produces a program in the \textit{same} language.)

Two ways of organizing the translation then suggest themselves:

\begin{numbered}
    \item We could perform most of the translation by successive transformations of one Miranda program into another, each transformation performing a simplification step. We would complete the process by translating the resulting (simple) Miranda program into the lambda calculus. The idea is that the earlier transformations would have done all the hard work, so the final step should consist of little more than a change of syntax.

    \item Alternatively, we could begin the translation by performing a simple syntactic translation of the Miranda program into an enriched version of the lambda calculus. This enriched lambda calculus would include the ordinary lambda calculus as a subset, but would also include extra constructs, chosen so that the first step consists of little more than a change of syntax. Then we could do most of the hard work by successively transforming the expression into simpler and simpler forms, until it becomes an ordinary lambda expression, free from any of the extra constructs.
\end{numbered}

Initially, the first method looks more attractive than the second, because it does not require us to define a new language (the enriched lambda calculus). However, we choose to follow the second course of action for the following reasons:

\begin{numbered}
    \item Miranda is designed to be a language for programmers, not compilers, and it lacks certain features that are desirable for a transformation-based compiler. (The particular features lacking are lambda abstractions and the ability to qualify any expression with local definitions. This is not a criticism of Miranda -- it just has a different purpose.)

    \item To a much greater extent than is the case for imperative languages, functional languages are largely syntactic variations of one another, with relatively few semantic differences. Using the second method allows the transformations we present to be applied easily to other languages, by altering only the translation of the high-level language into the enriched lambda calculus.
\end{numbered}

Figure 3.1 depicts the overall plan of action. We will use the term \textit{ordinary lambda calculus} to refer to the language described in Chapter 2, and \textit{enriched lambda calculus} to refer to the language introduced here.

The enriched lambda calculus is simply the ordinary lambda calculus augmented with extra constructs, chosen to allow an easy translation from Miranda. For each construct we will:

\begin{numbered}
    \item say what it looks like (give its syntax);
    \item say what it means (give its semantics).
\end{numbered}

\boxedfigure{
    \centering\sffamily\footnotesize

    \begin{minipage}{10cm}
        \centering

        \framebox{\ Miranda program\ }\\
        \hspace{4cm}$\Bigg\downarrow$ \begin{minipage}{4cm}
            {A simple translation\\ (specific to Miranda)}
        \end{minipage}\\
        \framebox{\
            \begin{minipage}{4.5cm}
                {Expression in the\\ enriched lambda calculus}
            \end{minipage}
            \ }\\
        \hspace{4cm}$\Bigg\downarrow$ \begin{minipage}{4cm}
            {Multiple transformations\\ (independent of Miranda)}
        \end{minipage}\\
        \framebox{\
            \begin{minipage}{4.5cm}
                {Expression in the\\ ordinary lambda calculus}
            \end{minipage}
            \ }\\

    \end{minipage}

}{Translation of Miranda into the lambda calculus}


The semantics for each construct can be given by providing a simple transformation which shows how to express that construct in terms of the ordinary lambda calculus. Then we could, in principle, translate from Miranda into the ordinary lambda calculus by first translating into the enriched lambda calculus, and then using the semantics of each construct repeatedly to transform the expression into an ordinary lambda expression.


While this method generates correct results, far greater efficiency is attainable by using more complicated transformations, but we can always confirm their correctness by reference to the inefficient version.

\section{The Enriched Lambda Calculus}
The enriched lambda calculus is a superset of the ordinary lambda calculus, so that any expression in the ordinary lambda calculus is also an expression in the enriched lambda calculus. The syntax for function application, lambda
abstractions, constants and built-in functions therefore remains exactly as described in Chapter 2. Likewise, all functions are written in prefix form, and the same conventions hold concerning brackets.

The only difference from the ordinary lambda calculus is the provision of four extra constructs. They are:

\begin{numbered}
    \item \ml{let}-expressions and \ml{letrec}-expressions;
    \item pattern-matching lambda abstractions;
    \item the infix operator \fatbar;
    \item case-expressions.
\end{numbered}

Of these, we will only describe the first here. The other three all concern pattern-matching, and cannot be defined before the discussion of pattern-matching itself. This is given in Chapter 4, and the remaining three constructs are defined there.

Figure 3.2 summarizes the syntax of the enriched lambda calculus for future reference.

\subsection{Simple let-expressions}
One of the main constructs in any functional language is the \textit{definition}, whereby a name is bound to a value. This mechanism is provided in the enriched lambda calculus, using \ml{let}-expressions and \ml{letrec}-expressions.

We begin by defining \textit{simple} \ml{let}-expressions. They are called `simple' by contrast with \textit{pattern-matching} \ml{let}-expressions, which we deal with later. A simple \ml{let}-expression has the following syntax:
\begin{mlcoded}
    let v = B in E
\end{mlcoded}

\boxedfigure{
\hspace{-2em}
\begin{tabular}{llll}
    <exp> & ::= & <constant> & Constants \\
          & | & <variable>   & Variables \\
          & | & <exp> <exp>  & Applications \\
          & | & \tl{} <variable> . <exp>   & Lambda abstractions \\
          & | & \ml{let\phantom{rec}} <pattern> = <exp> \ml{in} <exp> & \ml{Let}-expressions \\
          & | & \ml{letrec} <pattern> = <exp>  & \ml{Letrec}-expressions \\
          &   & \phantom{\ml{letrec}} $\cdots$ & \\
          &   & \phantom{\ml{letrec}} <pattern> = <exp>  &  \\
          &   & \ml{in} <exp> & \\
          & | & <exp> \fatbar{} <exp>  & Fat bar \\
          & | & \ml{case} <variable> \ml{of} & \ml{Case}-expression \\
          &   & <pattern>\quad $\Rightarrow$\quad <exp> & \\
          &   & $\cdots$ & \\
          &   & <pattern>\quad $\Rightarrow$\quad <exp> & \\
\end{tabular}

\hspace{-2em}
\begin{tabular}{lllll}
    <pattern> & ::= & <constant> &  & Constant patterns\\
              & |   & <variable> &  & Variable patterns\\
              & |   & <constructor> & <pattern> ,\qquad\qquad  & Constructor patterns\\
              &     &  & $\cdots$  & \\
              &     &  & <pattern>  & \\
\end{tabular}

}{Syntax of enriched lambda expressions}


\noindent where the \ml{v} is a variable, and \ml{B} and \ml{E} are expressions in the (enriched) lambda notation.

It introduces a definition for a variable \ml{v}, which binds \ml{v} to \ml{B} in \ml{E}. The definition is in scope with \ml{E} but not \ml{B}. We say that the `\ml{v = B}' is the \textit{definition of the} \ml{let}, the \ml{v} is the variable \textit{bound by the} \ml{let}, and the \ml{B} is the \textit{definition body}.

For example, consider the following \ml{let}-expression:
\begin{mlcoded}
    let x $=$ 3 in ($*$ x x)
\end{mlcoded}
Intuitively, the value of this expression is found by substituting \ml{3} for \ml{x} in the body \ml{($*$ x x)}, and then evaluating the body, giving the result \ml{9}:
\begin{mlalign}
    & let x $=$ 3 in ($*$ x x)\\
    $\rightarrow$ & $*$ 3 3\\
    $\rightarrow$ & 9
\end{mlalign}

A \ml{let}-expression is an expression like any other, and can be used in the same way as any other expression. For example,
\begin{mlalign}
    & $+$ 1 (let x $=$ 3 in ($*$ x x))\\
    $\rightarrow$ & $+$ 1 ($*$ 3 3)\\
    $\rightarrow$ & $+$ 1 9\\
    $\rightarrow$ & 10
\end{mlalign}

For the same reason, \ml{let}-expresions can be nested:
\begin{mlalign}
    & let x $=$ 3 in (let y $=$ 4 in ($*$ x y))\\
    $\rightarrow$ & let y $=$ 4 in ($*$ 3 y)\\
    $\rightarrow$ & $*$ 3 4\\
    $\rightarrow$ & 12
\end{mlalign}

As a matter of convenience, we also allow ourselves to write multiple definitions in the same let; thus:
\begin{letalign}
    let & x $=$ 3\\
        & y $=$ 4\\
    in  & $*$ x y
\end{letalign}

This expression means precisely the same as the previous one. We define a \ml{let}-expresion with several definitions to mean the same as the nested set of \ml{let}-expresions which defines the same variables in the same order, one per \ml{let}-expresion. (Syntactically, it would have been possible to specify that multiple definitions are separated with semicolons, but layout will suffice for our purposes.)

Earlier in this section we developed an informal reduction rule for \ml{let}-expresions. This involved \textit{substitution} and is very reminiscent of the \tb{}-reduction rule, which also uses substitution. For example, to evaluate
\begin{mlcoded}
    (\tlb{x}$*$ x x) 3
\end{mlcoded}

we substitute \ml{3} for \ml{x} in the body \ml{($*$ x x)}, and then evaluate the body. Generalizing this idea, we can now define the semantics of a simple \ml{let}-expresion as follows:
\begin{mlcoded}
    (let v $=$ B in E) $=$ ((\tlb{v}E) B)
\end{mlcoded}

(We use the symbol $=$ to denote the equivalence of two expressions.) That is all that is needed to define its semantics! By repeated application of this equivalence, we could eliminate all simple \ml{let}-expresions from an expression, in favor of lambda abstractions.

\subsection{Simple \ml{letrec}-expresions}
The syntax of a \textit{simple} \ml{letrec}-expresion is similar to that of a simple \ml{let}-expresion:
\begin{letalign}
    letrec & v$_1 =$ E$_1$ \\
           & v$_2 =$ E$_2$ \\
           & $\cdots$ \\
           & v$_n =$ E$_n$ \\
    in & \\
    & E
\end{letalign}

where the \ml{v$_i$} are variables, and \ml{E}, \ml{E$_1$}, ..., \ml{E$_n$} are expressions in the (enriched) lambda notation. We will sometimes abbreviate `\ml{letrec}-expresion' to `letrec' (and `\ml{let}-expresion' to `\ml{let}'), where no ambiguity arises.

The term `\ml{letrec}' is short for `let recursively', and it introduces possibly recursive bindings for a number of variables \ml{v$_i$}. The difference between \ml{let}s and \ml{letrec}s is that the \ml{v$_i$} are in scope in the \ml{E$_i$} (as well as \ml{E}) of a \ml{letrec}. To take an example, the expression
\noindent
\begin{letalign}
    letrec & factorial $=$ \tlb{n}IF ($=$ n 0) 1 ($*$ n (factorial ($-$ n 1)))\\
    in & factorial 4
\end{letalign}
defines a recursive function \ml{factorial}, and applies it to the argument \ml{4}. The value of the expression is thus \ml{24}.

Like \ml{let}-expresions, \ml{letrec}-expresions can appear embedded anywhere in an expression. Unlike \ml{let}-expresions, however, it is essential to allow multiple definitions in a \ml{letrec}-expresion, so as to permit mutual recursion. This is demonstrated by the following example:
\begin{letalign}
    letrec & f $= \ldots$ f $\ldots$ g $\ldots$ \\
           & g $= \ldots$ f $\ldots$ \\
    in $\ldots$ &
\end{letalign}

Here, \ml{f} refers to itself and \ml{g}, and \ml{g} refers to \ml{f}. This cannot be transformed into a nested pair of \ml{letrec}s, because then either \ml{g} would not be in scope in the body of \ml{f}, or vice versa.

It is easy to provide a semantics for a \ml{letrecs} with only a single definition, using the \ml{Y} operator developed in Section 2.4. In particular,
\begin{mlcoded}
    (letrec v $=$ B in E) $=$ (let v $=$ Y (\tlb{v}B) in E)
\end{mlcoded}
The use of \ml{Y} renders the definition non-recursive, so we can then use a \ml{let}-expresion, whose semantics has already been defined.

The case of multiple definitions requires the use of pattern-matching, and so is postponed until Chapter 6.

\subsection{Pattern-matching let$-$ and \ml{letrec}-expresions}
We will also allow \textit{patterns}, as well as variables, to appear on the left-hand side of definitions in \ml{let}s and \ml{letrec}s. We have not yet defined what a pattern is, so we postpone the topic until Chapter 6. However, a variable is just a simple form of pattern, so simple \ml{let(rec)}-expressions are just simple forms of pattern-matching \ml{let(rec)}-expressions.

\subsection{Let(rec)s versus Lambda Abstractions}
So far we have regarded the ordinary lambda calculus as the target language, into which we will transform the program, and \ml{let(rec)}-expressions as intermediate embellishments. However, there are strong efficiency reasons for including \textit{simple} \ml{let(rec)}-expressions in the target language, rather than transforming them into the ordinary lambda calculus.

Specifically, the transformation of a \ml{let}-expresion
\begin{mlcoded}
    let v $=$ B in E
\end{mlcoded}
into the application of a lambda abstraction
\begin{mlcoded}
    (\tlb{v}E) B
\end{mlcoded}
is using a sledgehammer (lambda abstraction) to crack a nut (\ml{let}-expresions). The lambda abstraction \ml{(\tlb{v}.E)} could be applied to many arguments, but it is in fact only ever applied to one, namely \ml{B}. The generality of lambda abstraction is not required, and the special case (that of application to a unique argument) can be exploited by the more sophisticated compilers described later in this book.

This issue manifests itself in a number of ways:
\begin{numbered}
    \item Miranda is a polymorphically typed language, and in Chapter 8 we give an algorithm for type-checking programs. Unfortunately, it is not possible to type-check the program once it has been transformed into the ordinary lambda calculus, but the addition of simple \ml{let(rec)}-expressions is sufficient to solve the problem.
    \item In all implementations except the very simplest, \ml{let}-expressions can be evaluated very much more efficiently than the corresponding application of a lambda abstraction. This applies to all the implementations described from Chapter 14 onwards.
    \item A related problem is that the transformation of \ml{letrec}-expressions into the ordinary lambda calculus compels us to use \ml{Y} to express recursion. The resulting expression is not an efficient implementation, and a more
    sophisticated compiler may wish to handle recursion in a different way
    (see Chapter 14). Keeping the recursion explicit using \ml{letrec} allows scope
    for these optimizations.
\end{numbered}

To summarize, all our implementations, except the very simplest, will require the program to be transformed into the ordinary lambda calculus augmented with simple \ml{let(rec)}-expressions. This approach makes a dramatic contribution to the efficiency of the resulting implementations. On the other hand, little seems to be gained by augmenting the language still further.

\section{Translating Miranda into the Enriched Lambda Calculus}
A program consists of a set of definitions, together with an expression to be evaluated. To keep these two components of the program separate we will use a box, like this:

\begin{center}
\setlength{\tabcolsep}{18pt}
\renewcommand{\arraystretch}{1.5}
\begin{tabular}{|c|}
    \hline
    Set of definitions \\
    \hdashline
    Expression to be evaluated \\
    \hline
\end{tabular}
\end{center}

\noindent For example, we could compute twice the square of 5 with the following Miranda program:
\begin{center}
\setlength{\tabcolsep}{18pt}
\renewcommand{\arraystretch}{1.5}
    \begin{tabular}{|c|}
        \hline
        \ml{square n $=$ n*n } \\
        \hdashline
        \ml{2 $*$ (square 5)} \\
        \hline
    \end{tabular}
\end{center}

\noindent(Note: Miranda is an interactive language, and defines a `program' to be a set of definitions, while the `expression to be evaluated' is typed by the user. For the rest of this book, however, we will use `program' to mean `a set of definitions together with an expression to be evaluated'.) Proceeding informally, we can translate this Miranda program into the enriched lambda calculus quite easily, to produce the expression
\begin{letalign}
    let &square $=$ \tlb{n}$*$ n n\\
    in &($*$ 2 (square 5))
\end{letalign}

We now introduce some notation to help describe the translation process. Consider the translation of the Miranda expression \ml{(2 $*$ (square 5))} into the lambda expression \ml{($*$ 2 (square 5))}. We may regard this translation process as a \textit{function} \metafn{TE}, which takes the Miranda expression as its input, and produces the lambda expression as its output. We write the translation like this:
\begin{mlcoded}
    \metafnbb{TE}{2 $*$ (square 5)} $\equiv *$ 2 (square 5)
\end{mlcoded}
The double square brackets \doublebracket{} are used to enclose the Miranda expression,
to emphasize that the argument to \metafn{TE} is a \textit{syntactic} object. This convention was used in Chapter 2, but the difference on this occasion is that the result of the translation is a syntactic object also, and we use \ml{$\equiv$} rather than \ml{$=$} to remind us of this fact. We call \metafn{TE} a \textit{translation scheme}.

We also need another translation scheme \metafn{TD}, which translates Miranda definitions into definitions suitable for a \ml{letrec}. For example,

\begin{mlcoded}
    \metafnbb{TD}{square n $=$ n*n} $\equiv$ square $=$ \tlb{n}$*$ n n
\end{mlcoded}

Here we see another reason for using \ml{$\equiv$} when writing translation schemes: it avoids confusion with \ml{$=$} symbols in the program being translated. We can now generalize the translation scheme as follows. Given the Miranda program

\begin{center}
    \setlength{\tabcolsep}{18pt}
    \renewcommand{\arraystretch}{1.5}
    \begin{tabular}{|c|}
        \hline
        \ml{\strut Definition$_1$} \\
        $\vdots$\\
        \ml{\strut Definition$_n$} \\
        \hdashline
        \ml{\strut Expression} \\
        \hline
    \end{tabular}
\end{center}
we generate the following (enriched) lambda expression:
\begin{mlcoded}
    letrec \\
    \phantom{ww}\metafnbb{TD}{Definition$_1$} \\
    \phantom{ww}$\cdots$ \\
    \phantom{ww}\metafnbb{TD}{Definition$_n$} \\
    in\\
    \phantom{ww}\metafnbb{TE}{Expression}
\end{mlcoded}
In the previous example we used a \ml{let} instead of a \ml{letrec}, but Miranda definitions are all potentially recursive, so we must use a \ml{letrec} in general (later work will optimize this--Section 6.2.8).

What we have now done is to reduce the translation problem to one of defining the two translation schemes \metafn{TD} and \metafn{TE}. We will define them for simple cases in the succeeding two sections, and then lay out the plan of the next few chapters, which will extend them to cover more complicated cases.

For the moment, we completely avoid the question of declarations of new types and type-checking. The former will be introduced in Chapter 4 and the latter in Chapter 8.

\subsection{The \metafn{TE} Translation Scheme}

The translation scheme \metafn{TE} is a function, which takes a Miranda expression as its argument, and produces an equivalent lambda expression as its result, thus:

\begin{center}
    \framebox{\strut\ Miranda expression\ } $\xrightarrow{\qquad\text{\normalsize \metafn{TE}}\qquad}$ \framebox{\strut\ Lambda expression\ }
\end{center}

\noindent We will describe \metafn{TE} by case analysis, giving a rule for each possible form of a
Miranda expression.

\chapter[Structured Types and the Semantics of Pattern-Matching][Structured Types and the Semantics of Pattern-Matching]{Structured Types and the\\Semantics of\\Pattern-Matching}
\chapterauthors{Simon L. Peyton Jones and Philip Wadler}
\vspace{3cm}


\noindent This chapter concerns structured types, a powerful and general mechanism
for defining data types, provided by several functional languages, including
Miranda, ML and Hope. Intimately associated with structured types is a
notational device known as pattern-matching, which is used by such
languages for defining functions.

Section 4.1 gives a general introduction to structured types and pattern-matching. Section 4.2 begins with a more in-depth look at pattern-matching
and conditional equations, and then introduces two new constructs in the
enriched lambda calculus, \fatbar{} and pattern-matching lambda abstractions. Using
these constructs, we then show how to translate a general Miranda function
definition into the enriched lambda calculus. Section 4.3 is devoted to
providing a precise semantics for pattern-matching lambda abstractions.

We conclude in Section 4.4 by defining case-expressions, the last new
construct of the enriched lambda calculus. This clears the way for Chapter 5,
which will show how to transform pattern-matching lambda abstractions into
case-expressions, thus giving a considerable gain in efficiency.

What in this chapter are called `structured types' are called `algebraic types'
in Miranda, and `free data types' by some others [Burstall and Goguen, 1982].

\section{Introduction to Structured Types}

Suppose that we wish to define binary trees with leaves that are numbers. In the notation of Miranda, this could be done by declaring a structured type tree as follows:
\begin{mlcoded}
    tree \typedecl{} \ml{LEAF} num | \ml{BRANCH} tree tree
\end{mlcoded}
(The symbol `\typedecl{}' identifies this as a type declaration.) This might be read as follows: `a tree is either a \ml{LEAF}, which contains a \ml{num}, or a \ml{BRANCH}, which contains a tree and a tree'. Here \ml{LEAF} and \ml{BRANCH} are called \textit{constructors} of the type. Miranda requires that constructors (and only constructors) begin with an upper-case letter, but we will always write them entirely in upper case. \ml{LEAF} has one \textit{field}, of type \ml{num}, and \ml{BRANCH} has two, both of type tree. The number of fields associated with a constructor is called its \textit{arity}; thus \ml{LEAF} has arity 1 and \ml{BRANCH} has arity 2.

Constructors can be used as functions, to create values of type \ml{tree}. For example, the following equation
\begin{mlcoded}
    tree1 $=$ BRANCH (BRANCH (LEAF 1) (LEAF 2)) (LEAF 3)
\end{mlcoded}
defines \ml{tree1} to be a \ml{tree}. Informally, this tree might be drawn as:
\begin{mlcoded}
\begin{center}
    \begin{forest}
        [. [. [1] [2]] [3]]
    \end{forest}
\end{center}
\end{mlcoded}

Constructors can also appear on the left-hand side of an equation, as in the following Miranda function definition:
\begin{mlcoded}
    \begin{tabular}{ll}
    reflect (LEAF n) &$=$ LEAF n \\
    reflect (BRANCH t1 t2) &$=$ BRANCH (reflect t2) (reflect t1)
    \end{tabular}
\end{mlcoded}
For example, \ml{(reflect tree1)} returns
\begin{mlcoded}
    BRANCH (LEAF 3) (BRANCH (LEAF 2) (LEAF 1))
\end{mlcoded}

A definition with patterns on the left-hand side, such as that of \ml{reflect}, is said to use \textit{pattern-matching} to perform \textit{case analysis}. For example, in evaluating \ml{(reflect t)} there are two cases to choose from: \ml{t} matches the pattern \ml{(LEAF n)}, or \ml{t} matches the pattern \ml{(BRANCH t1 t2)}. If, say, \ml{t} is \ml{(LEAF 1)} then the first case is chosen, with \ml{n} bound to \ml{1}. Much more will be said about pattern-matching later.

An important difference in the treatment of structured types in Miranda from that in ML or Hope, is that in Miranda constructor functions are lazy; that is, they do not evaluate their arguments. The components of a structured object are evaluated only when (and if) they are subsequently extracted and used, not when the object is built.

\subsection{Type Variables}

Type declarations may also contain type variables. For example, the definition of the type tree above may be rewritten to allow trees with leaves of any type:
\begin{mlcoded}
    tree $*$ \hastype$=$ \ml{LEAF} $*$ | \ml{BRANCH} (tree $*$) (tree $*$)
\end{mlcoded}
Here $*$ is called a \textit{generic} (or \textit{schematic}) \textit{type variable}. The declaration could be read as follows: `a tree of $*$ is either a \ml{LEAF}, which contains a $*$, or a \ml{BRANCH} which contains a tree of $*$ and a tree of $*$, for any type $*$'.

Leaves of any particular tree must all contain values of the same type, but different trees may have leaves of different types. Examples of trees and their types are
\begin{mlcoded}
    \begin{tabular}{ll}

    \ml{BRANCH} (\ml{LEAF} 1) (\ml{LEAF} 2)
    &\hastype tree num\\
    \ml{BRANCH} (\ml{LEAF} `a') (\ml{LEAF} `b')
    &\hastype tree char

    \end{tabular}
\end{mlcoded}
(The symbol \hastype is pronounced `has type'.) Here, `tree' is called a \textit{type-forming operator}, since it takes a type (such as \ml{num} or \ml{char}) as an `argument' and produces a type (respectively, \ml{(tree num)} or \ml{(tree char)}).

The repeated use of $*$ on the right-hand side of the type declaration specifies that the two branches of a tree must be of uniform type. For example,
\begin{mlcoded}
    \ml{BRANCH} (\ml{LEAF} 1) (\ml{LEAF} 'a')
\end{mlcoded}
is not legal, since it has leaves of mixed type. More will be said about types and type variables in Chapter 8.

\subsection{Special Cases}

This section shows how three `built-in' types, namely lists, tuples and enumerated types, can be regarded as instances of general structured types.

\subsubsection{Lists}

Miranda has a special syntax to denote lists, but lists are just an instance of a general structured type. Lists could be defined as follows:

\begin{mlcoded}
    list $*$ \hastype$=$ \ml{NIL} | \ml{CONS} $*$ (list $*$)
\end{mlcoded}

This type declaration defines the two new constructors \ml{NIL} and \ml{CONS}. Miranda's built-in syntax for lists could then be translated to use \ml{NIL} and \ml{CONS}, as follows:

\begin{quote}
    \ml{[ ]} is translated to \ml{NIL}\\
    \ml{(x:xs)} is translated to \ml{(CONS x xs)}.\\
    \ml{[x,y,z]} is a Miranda abbreviation for \ml{(x:y:z:[ ])} and hence is translated to
    \ml{(CONS x (CONS y (CONS z NIL)))}\\
    \ml{[$*$]} is translated to (list \ml{$*$})
\end{quote}

\begin{figure}[H]
\centering

{%
    \setlength{\fboxrule}{1pt}%
    \setlength{\fboxsep}{10pt}%
    \fbox{%
        \begin{minipage}{\textwidth}
            \footnotesize
\begin{mlcoded}
    \begin{tabular}{lll}
        \metafnbb{TE}{:}  &$\equiv$ &CONS \\
        \metafnbb{TE}{[ ]}  &$\equiv$ &NIL \\
        \metafnbb{TE}{[E$_1$, E$_2$, $\ldots$, E$_n$]}  &$\equiv$ &CONS \metafnbb{TE}{E$_1$} \metafnbb{TE}{[E$_2$, $\ldots$, E$_n$]}\\
        \metafnbb{TE}{(E$_1$, E$_2$)}  &$\equiv$ &PAIR \metafnbb{TE}{E$_1$} \metafnbb{TE}{E$_2$} \\
        \metafnbb{TE}{(E$_1$, E$_2$, E$_3$)}  &$\equiv$ &TRIPLE \metafnbb{TE}{E$_1$} \metafnbb{TE}{E$_2$} \metafnbb{TE}{E$_3$}\\
        {\normalfont\normalsize and so on} & & \\
        \metafnbb{TE}{True}  &$\equiv$ &TRUE \\
        \metafnbb{TE}{False}  &$\equiv$ &FALSE \\
    \end{tabular}
\end{mlcoded}
        \end{minipage}%
    }%
}%

\caption{\textsf Modifications to the TE scheme for lists, tuples and booleans}
\end{figure}

\noindent(Note: the last example is different from the others, because it describes a type-expression rather than a value-expression.)

We can conveniently perform this translation when translating from Miranda into the enriched lambda calculus; Figure 4.1 gives the required equations.

Notice that the elements of a list of type \ml{(list $*$)} must all be of type \ml{$*$}, but the number of elements in a list is not determined by its type. Thus \ml{(CONS} 2 \ml{NIL)} and \ml{(CONS 3 (CONS 6 NIL)} are both of type \ml{(list num)}, though they are of different lengths.

\subsubsection{Tuples}

Miranda also provides special syntax to denote tuples, and these also can be defined using a structured type. Tuples could be defined as follows:

\begin{mlcoded}
    \footnotesize
    \begin{tabular}{llll}
    pair        & $*\ **$           &\typedecl{} PAIR        & $*\ **$ \\
    triple      & $*\ **\ ***$       &\typedecl{} TRIPLE      & $*\ **\ ***$ \\
    quadruple   & $*\ **\ ***\ ***$   &\typedecl{} QUADRUPLE   & $*\ **\ ***\ ***$
    \end{tabular}
\end{mlcoded}

Notice the difference between `\ml{pair}' and \ml{PAIR}: the former is a type-forming operator, used only in type-expressions, while the latter is the constructor function of the type, used only in value-expressions.

As with lists, Miranda's special syntax can be translated as follows:
\begin{quote}
    \ml{(x,y)} is translated to \ml{(PAIR x y)}\\
    \ml{(x,y,z)} is translated to \ml{(TRIPLE x y z)}\\
    and so on.\\
    \ml{($*$, $**$)} is translated to (pair $*$ $*$ $*$)\\
    \ml{($*$, $**$, $***$)} is translated to \ml{(triple $*$ $**$ $***$)}
\end{quote}
Figure 4.1 gives the required equations.

Notice that a tuple may contain elements of mixed type; for example
\begin{mlcoded}
\begin{tabular}{ll}
    (3, TRUE) &\hastype{} PAIR num bool\\
    (`a', (3, 2)) &\hastype{} PAIR char (PAIR num num)
\end{tabular}
\end{mlcoded}
However, the type of a tuple completely determines the number and the types
of its fields. For example, a \ml{pair} always contains exactly two fields, a \ml{triple}
contains exactly three fields, and so on.

\subsubsection{Enumerated types}

The type declaration

\begin{mlcoded}
    color \hastype$=$ \ml{VERMILLION} | \ml{PUCE} | \ml{LAVENDER}
\end{mlcoded}
in which each constructor has zero fields, is just like an enumerated type in Pascal. Thus, we can define the type of boolean values:
\begin{mlcoded}
    bool \hastype$=$ \ml{TRUE} | \ml{FALSE}
\end{mlcoded}

The usual functions on booleans can then be defined using pattern-matching; for example:
\begin{mlcoded}
    \begin{tabular}{ll}
    if \ml{TRUE} &e1 e2 $=$ e1\\
    if \ml{FALSE} &e1 e2 $=$ e2
    \end{tabular}
\end{mlcoded}

Miranda uses the names `\ml{True}' and `\ml{False}' for its built-in truth-values.

\subsubsection{Summary}

Since it is easy to translate `built-in' types like lists and tuples into equivalent structured types, then any implementation of a functional language that handles structured types will also handle these `built-in' types for free. This can greatly simplify an implementation. Instead of implementing several type mechanisms, one for lists, one for tuples, one for enumerated types, and so on, we need only implement a single mechanism for structured types, and translate other types into structured types. Figure 4.1 gives the required equations.

\subsection{General Structured Types}

In general, the form of a structured type definition is:

\begin{mlcoded}
    T \hastype$=$ c$_1$ T$_{1,1}$ $\ldots$ T$_{1,r_1}$\\
    \quad | $\ldots$\\
    \quad | c$_n$ T$_{n,1}$ $\ldots$ T$_{n,r_n}$\\
\end{mlcoded}
where the \ml{T$_{i,j}$} are types and the \ml{c$_i$} are constructors of arity $r_i$. In the `\ml{tree}' example above, \ml{T} was \ml{(tree $*$)}, \ml{c$_1$} was \ml{LEAF}, \ml{T$_{1,1}$} was \ml{num}, \ml{c$_2$} was \ml{BRANCH}, \ml{T$_{2,1}$} was \ml{(tree $*$)}, and \ml{T$_{2,2}$} was \ml{(tree $*$)}.

Readers familiar with the mathematical operations for constructing types will recognize that the general type above can be written as the sum (that is, discriminated union):
\begin{mlcoded}
T $=$ T$_1$ $+ \ldots +$ T$_n$
\end{mlcoded}
where each \ml{T$_i$}, for $i$ from 1 to $n$, can be written as a product:
\begin{mlcoded}
T$_i =$ T$_{i,1} \times$ T$_{i,2} \times \cdots \times$ T$_{i,r_i}$
\end{mlcoded}
In other words, a structured type is a \textit{sum-of-products}.

When $n=1$ we say that the type is a \textit{product type}; the types \ml{(pair $*\ **$)}, \ml{(triple $*\ **\ ***$)}$\ldots$ are all product types. When $n>1$ we say that the type is a \textit{sum type}, since it is the sum of more than one domain; the types \ml{(tree $*$)}, \ml{(list $*$)}, \ml{color} and \ml{bool} are all sum types. Thus a product type has exactly one constructor, and a sum type has two or more constructors.

We will often wish to distinguish between the constructors of product types and sum types. Just as we use the names \ml{c$_i$} to stand for constructors of all types, we will use the name \ml{t} to stand for the constructor of a product type, and the names \ml{s} and \ml{s$_i$} to stand for the constructors of a sum type (\ml{t} suggests `tuple' and \ml{s} suggests `sum').

(Note: we use lower-case letters to stand for constructors, to avoid confusion with the constructors themselves, which are written in upper case. Similarly, we use upper-case letters to stand for types, which are themselves written in lower case $-$ see Section 4.1.)

(\textit{Important:} at the time when this chapter was first written the semantics of Miranda provisionally specified that a structured type with only one constructor was a product type, as above. However, an alternative view is that a structured type with only one constructor should behave as a sum type with one component in the sum, and that product types (tuples) be treated as an independent construct. It now seems likely that Research Software Limited will follow this latter course in their definition of Miranda. As a consequence some of the statements made in this chapter about the semantics of structured types in Miranda may be incorrect. We draw the reader's attention to the caveat on page 37.)

\subsection{History}

As mentioned, structured types are a combination of sum types and product types, which have a long history in mathematics.

Landin's Iswim, one of the earliest functional languages, was described using a stylized form of English for defining structured types [Landin, 1966]. Burstall introduced a more formal notation for defining such types in NPL [Burstall, 1977]. Hope and ML have type systems based on separate sum and product types, whereas Miranda and Orwell have type systems based on sum-of-product types.

Iswim also contained a simple form of pattern-matching, where one could write definitions such as
\begin{mlcoded}
    addPair (x,y) $=$ x $+$ y
\end{mlcoded}
However, the important idea of using pattern-matching for case analysis appears to have been developed independently by Burstall and Turner. Pattern-matching appeared in NPL and SASL, and was used to good effect in proofs by structural induction [Burstall, 1969] and program transformation
[Burstall and Darlington, 1977]. It was incorporated into many later
languages such as Hope, KRC, ML, Miranda and Orwell.

\section{Translating Miranda into the Enriched Lambda Calculus}

We must now demonstrate how to translate Miranda function definitions involving pattern-matching into the enriched lambda calculus. In the process of doing so we will introduce \textit{pattern-matching lambda abstractions} and the \fatbar{} operator, two of the constructs in the enriched lambda calculus whose explanation was postponed.

\subsection{Introduction to Pattern-matching}
We begin this section by illustrating some further aspects of pattern-matching, which have to be handled by an implementation. (Not all the illustrations should be taken as examples of good programming style. Some are expressly chosen to demonstrate all the possible nasty things that can happen!)  Recall the definition of \ml{reflect}:
\begin{mlcoded}
    \begin{tabular}{ll}
    reflect (LEAF n) &$=$ LEAF n\\
    reflect (BRANCH t1 12) &$=$ BRANCH (reflect 12) (reflect t1)
    \end{tabular}
\end{mlcoded}

The terms \ml{(LEAF n)} and \ml{(BRANCH t1 12)} occurring on the left-hand side of these equations are called \textit{patterns}. When \ml{reflect} is applied to an argument, the argument is first evaluated to see whether it \textit{matches} the pattern \ml{(LEAF n)} or \ml{(BRANCH t1 12)}. It will certainly match one or the other, because the type-checker ensures that \ml{reflect} is only applied to objects of type \ml{(tree $*$)}, for some type \ml{$*$}. For example, if \ml{reflect} is applied to an expression which evaluates to \ml{(BRANCH E$_1$ E$_2$)}, the second equation is selected, with \ml{t1} bound to \ml{E$_1$} and \ml{12} bound to \ml{E$_2$}.

In the preceding example, the order in which the equations were written was immaterial, but this is not always the case. Consider the Miranda function definition
\begin{mlcoded}
    factorial 0 $=$ 1\\
    factorial n $=$ n $*$ factorial (n-1)
\end{mlcoded}

The order of the equations in this definition is significant. In the evaluation of \ml{(factorial x)}, there are two cases to choose from: either \ml{x} matches \ml{0} (that is, \ml{x} evaluates to \ml{0}), so the first equation is chosen, or it does not, so the second case is chosen with \ml{n} bound to \ml{x}. The equations are tried out one at a time, from top to bottom. If they had been written in the other order then the first equation would always match. In this situation we say that the patterns \textit{overlap}. (As we shall see in Chapter 5, there are good reasons to avoid writing overlapping patterns, but occasionally they prove useful.)

Another point, illustrated by the first \ml{factorial} equation, is that a pattern may consist of a literal constant, such as a number or character.

As another example, consider the Miranda function definition
\begin{mlcoded}
    \begin{tabular}{ll}
    lastElt (x:[]) &$=$ x\\
    lastElt (x:xs) &$=$ lastElt xs
    \end{tabular}
\end{mlcoded}
The function call \ml{(lastElt xs)} returns the last element of the list \ml{xs}. Again, the order of the equations is significant, since the patterns overlap. Furthermore, the first pattern is an example of a \textit{nested} pattern, in which the pattern \ml{[\,]} is nested inside the pattern \ml{(x:[\,])}. Finally, the equations are not exhaustive, since neither pattern matches the argument \ml{[\,]}. If \ml{lastElt} is applied to \ml{[\,]} some sort of error should be reported.

Pattern-matching can apply to several arguments, as the following Miranda definition shows:
\begin{mlcoded}
    \begin{tabular}{lll}
    xor False &y &$=$ y\\
    xor True &False &$=$ True\\
    xor True &True &$=$ False
    \end{tabular}
\end{mlcoded}

Another feature of Miranda that is closely connected with pattern-matching is \textit{conditional equations}, which control the selection of \textit{alternatives} by the use of \textit{guards}. We could, for example, rewrite the \ml{factorial} function in the following way:
\begin{mlalign}
    factorial n & $=$ 1,\qquad n=0\\
              n & $=$ n $*$ factorial (n-1)
\end{mlalign}
A single left-hand side governs several alternatives, which together constitute the right-hand side. In this case there is only one guard, namely the boolean-valued expression `\ml{n=0}', which appears following a comma. Guards are evaluated one at a time, beginning at the top, and when a guard evaluates to \ml{True}, the corresponding alternative expression is selected. The guard may be omitted in the final right-hand side, giving an `otherwise' case (equivalent to a guard of \ml{True}).

The \ml{factorial} example shows, incidentally, that a constant appearing in a pattern can easily be eliminated by replacing it with a variable and adding a guard to the equation instead.

Conditional equations interact with pattern-matching, as demonstrated in the next example. The function \ml{funnyLastElt} returns the last element of its argument list, except that if a negative element is encountered then it is returned instead:
\begin{mlcoded}
    \begin{tabular}{ll}
        funnyLastElt (x:xs) &$=$ x,\qquad x<0\\
        funnyLastElt (x:[]) &$=$ x\\
        funnyLastElt (x:xs) &$=$ funnyLastElt xs
    \end{tabular}
\end{mlcoded}
Pattern-matching proceeds, as usual, from top to bottom; when a left-hand side matches the argument, the guarded alternative(s) are tried, from top to bottom. If none of the guards is \ml{True}, then pattern-matching continues, Starting with the next equation. Applying \ml{funnyLastElt} to the list \ml{[1,2]} would cause this behavior, since the first equation would match, but the guard fails, so the second and then third equations are tried.

Finally, variables may be repeated on the left-hand side of an equation. For example, the function \ml{noDups} eliminates adjacent duplicate elements in a list:
\begin{mlcoded}
    \begin{tabular}{ll}
    noDups [\,] &$=$ [\,]\\
    noDups [x] &$=$ [x]\\
    noDups (x:x:xs) &$=$ noDups (x:xs)\\
    noDups (x:y:ys) &$=$ x : noDups (y:ys)
    \end{tabular}
\end{mlcoded}
The third equation matches only if the first two elements of the argument list are equal; the repeated use of \ml{x} on the left-hand side implies the equality condition.

We may summarize the features that the implementation must support as follows:
\begin{numbered}
    \item overlapping patterns;
    \item constant patterns;
    \item nested patterns;
    \item multiple arguments;
    \item non-exhaustive sets of equations;
    \item conditional equations;
    \item repeated variables.
\end{numbered}
Given these complications, it is unwise to rely on a purely intuitive understanding of what a function definition using pattern-matching means. The rest of this section and the next is therefore devoted to providing a formal semantics of pattern-matching.

\subsection{Patterns}
First of all, we will need a precise definition of patterns.

\definitionbox{{\normalfont A pattern \ml{p} is}}{
\begin{tabular}{ll}
    either & a variable \ml{v}\\
    or     & a constant \ml{k}, such as a number, a character, a boolean and so on,\\
    or &a constructor pattern, of the form \ml{(c p$_1$ \ldots p$_r$)} where \ml{c} is a\\
    & constructor of arity \ml{r}, and p$_1,\, \ldots,\, $p$_r$ are themselves patterns.
\end{tabular}
\vspace{0.25\baselineskip}

All of the variables in a pattern should be distinct.

A pattern of the form \ml{(s p$_1$ \ldots p$_r$)}, where \ml{s} is a sum constructor, is called a \textit{sum-constructor pattern}, or \textit{sum pattern}. A pattern of the form \ml{(t p$_1$ \ldots p$_r$)}, where \ml{t} is a product constructor, is called a \textit{product-constructor pattern}, or \textit{product pattern}.
\vspace{0.25\baselineskip}

\noindent Note: according to this definition, patterns may not contain repeated variables, although Miranda allows them to do so. This point is discussed in Section 4.2.7.
}

Here are some examples of patterns:

\begin{tabular}{ll}
\ml{x} & \\
\ml{LEAF n} & \\
\ml{BRANCH (LEAF n) t} & \\
\ml{CONS x xs} & (written \ml{(x:xs)} in Miranda) \\
\ml{CONS x (CONS 3 NIL)} & (written \ml{[x,3]} in Miranda) \\
\ml{PAIR x 4} &(written \ml{(x,4)} in Miranda)
\end{tabular}

\noindent The term \ml{(PAIR z z)} is not a pattern, because it contains a repeated variable. The term \ml{(CONS x)} is not a pattern, because the \ml{CONS} does not have enough arguments.

Miranda allows patterns with repeated variables, like \ml{(PAIR z z)} but the patterns defined here do not. This is discussed in Section 4.2.7.

A constructor pattern is \textit{simple} if it has the form \ml{(c v$_1$ $\ldots$ v$_r$)}, where \ml{v$_1$, $\ldots$, v$_r$} are distinct variables. If a constructor pattern is not simple it is \textit{nested}.

\subsection{Introducing Pattern-matching Lambda Abstractions}

Up to now we have translated function definitions into the lambda calculus using the following rule:
\begin{mlcoded}
	\metafnbb{TD}{f v$_1$ ... v$_n$ $=$ E}   $\equiv$ f $=$ \tl{}v1\ldots\tlb{v$_n$}\metafnbb{TE}{E}
\end{mlcoded}
where \ml{v$_1$, $\ldots$, v$_n$} are variables. Temporarily restricting our attention to functions of a single variable, we could derive the less general rule
\begin{mlcoded}
	\metafnbb{TD}{f v $=$ E} $\equiv$ f $=$ \tlb{v}\metafnbb{TE}{E}
\end{mlcoded}
By analogy, given the function definition
\begin{mlcoded}
	\ml{f p $=$ E}
\end{mlcoded}
(where \ml{p} is a pattern), it seems plausible to translate it using the rule
\begin{mlcoded}
	\metafnbb{TD}{f p $=$ E} $\equiv$ f $=$ \tlb{p}\metafnbb{TE}{E}
\end{mlcoded}

This is not quite right yet, because we must remember to translate the pattern, so that Miranda's list notation is translated into uses of \ml{CONS} and \ml{NIL} (and likewise for tuples and booleans). Fortunately, the syntax of patterns is a subset of that of expressions, so we can use the \metafn{TE} scheme.
\begin{mlcoded}
	\metafnbb{TD}{f p $=$ E} $\equiv$ f $=$ \tlb{\metafnbb{TE}{p}}\metafnbb{TE}{E}
\end{mlcoded}
For example, consider the Miranda function definition for \ml{fst}:
\begin{mlcoded}
	fst (x, y) $=$ x
\end{mlcoded}
Using the rule above gives:
\begin{mlcoded}
	\metafnbb{TD}{fst (x, y) $=$ x} $\equiv$ fst $=$ \tlb{(PAIR x y)}x
\end{mlcoded}
This introduces a new sort of lambda abstraction, a \textit{pattern-matching lambda abstraction}, which has the form \ml{(\tlb{p}E)} where \ml{p} is a pattern. This leaves us with two questions:
\begin{numbered}
	\item How can we translate a general Miranda function definition into pattern-matching lambda abstractions?
	\item What, exactly, does \ml{(\tlb{p}E)} mean?
\end{numbered}
We discuss the first in the remainder of this section, leaving the second for the next section.

\subsection{Multiple Equations and Failure}

Consider first a Miranda function definition of the form
\begin{mlcoded}
	f p$_1$ $=$ E$_1$\\
	f p$_2$ $=$ E$_2$\\
	$\cdots$\\
	f p$_n$ $=$ E$_n$
\end{mlcoded}
Intuitively, we expect the semantics to be `try the first equation, and if that fails try the second, and so on'. This introduces the idea that a pattern-match might \textit{fail}. Such failure does not necessarily indicate an error, since there might be a subsequent equation which would match. Hence, we introduce a new built-in value \ml{FAIL}, which is returned when a pattern-match fails.

With the aid of this idea, we can translate the definition of \ml{f} into the following enriched lambda calculus expression:

\begin{mlalign}
	f $=$  \tlb{x}&(  (( \tlb{p$_1'$}E$_1'$) x)\\
	& \fatbar{} (( \tlb{p$_2'$}E$_2'$)  x)\\
	& $\cdots$ \\
	& \fatbar{} (( \tlb{p$_n'$}E$_n'$)  x)\\
	& \fatbar{} ERROR )
\end{mlalign}

where \ml{x} is a new variable name that does not occur free in any \ml{E$_i$}, the expressions \ml{E$_i'$} are the result of translating the \ml{E$_i$}, and the patterns \ml{p$_i'$} are the result of translating the \ml{p$_i'$}. The new definition of \ml{f} can be read `try to apply \ml{(\tlb{p$_1'$}E$_1'$)} to \ml{x}, and if that succeeds return its result; otherwise try \ml{(\tlb{p$_2'$}E$_2'$)}, and so on; if they all fail, return \ml{ERROR}'.

Here \ml{ERROR} is meant to be a special value whose evaluation indicates an error, an event which should never occur.

The function \ml{\fatbar{}} is an infix function, whose behavior is described by the semantic equations:
\begin{mlcoded}
	\begin{tabular}{ll}
		a &\fatbar{} b $=$ a \qquad {\normalfont if} a $\neq \bot$ {\normalfont and} a $\neq$ FAIL\\
	FAIL & \fatbar{} b $=$ b  \\
	$\bot$ & \fatbar{} b $= \bot$
	\end{tabular}
\end{mlcoded}

Operationally, \fatbar{} evaluates its left argument; if the evaluation terminates and yields something other than \ml{FAIL}, then \fatbar{} returns that value (first rule); if it evaluates to \ml{FAIL}, \fatbar{} returns its right argument (second rule); if the evaluation of the left argument fails to terminate, then so does the application of \fatbar{} (third rule).

It is easy to verify that \fatbar{} is an \textit{associative} operator, and has \textit{identity} \ml{FAIL}. Its associativity means that we may write expressions such as \ml{(E$_1$ \fatbar{} E$_2$ \fatbar{} E$_3$)} without ambiguity. It is extremely convenient to write \fatbar{} between its operands (that is, infix) but, since all functions are written prefix in the lambda calculus, we are forced to dignify \fatbar{} by making it one of the new constructs of the enriched lambda calculus. The sole reason for doing so is notational.

As an example of the suggested translation in action, recall the definition of the \ml{reflect} function:
\begin{mlcoded} \footnotesize
    \begin{tabular}{ll}
    reflect (LEAF n) &$=$ LEAF n \\
    reflect (BRANCH t1 t2) &$=$ BRANCH (reflect t2) (reflect t1)
    \end{tabular}
\end{mlcoded}
This would be translated to:
\begin{mlalign}\footnotesize
    \tlb{reflect} $=$ &\tlb{t}( ((\tlb{(LEAF n)}LEAF n) t) \\
    &\fatbar{} ((\tlb{(BRANCH t1 t2)}BRANCH (reflect t2) (reflect t1)) t)\\
    &\fatbar{} ERROR )
\end{mlalign}

In this case, of course, \ml{ERROR} can never be returned, since one of the previous pattern-matches will succeed. This is not always the case, as the following example shows. Consider the Miranda definition of \ml{hd}, which extracts the first element of a list:
\begin{mlcoded}
    hd (x:xs) $=$ x
\end{mlcoded}
It would be translated to
\begin{mlcoded}
    hd $=$ \tlb{xs'}(((\tlb{(CONS x xs)}x) xs') \fatbar{} \ml{ERROR})
\end{mlcoded}
If \ml{hd} is applied to \ml{NIL}, then \ml{ERROR} will be the result. (We have used \ml{xs'} as the formal parameter of the lambda abstraction, to avoid confusion with the \ml{xs} in the pattern. Technically, however, there would be no problem with using \ml{xs}, or any other variable, since \ml{hd} has no free variables.)

\subsection{Multiple Arguments}

Functions with multiple arguments are easily handled. As we recalled earlier, the basic approach is to translate a function of several arguments using the rule
\begin{mlcoded}
    \metafnbb{TD}{f v$_1$ $\ldots$ v$_n$ $=$ E} $\equiv$ f $=$ \tlb{v$_1$}$\cdots$\tlb{v$_n$}\metafnbb{TE}{E}
\end{mlcoded}
Combining this with the approach of the previous section suggests that we should translate the definition
\begin{mlcoded}
    \ml{f p$_1$ p$_2$ $\ldots$ p$_m$ $=$ E}
\end{mlcoded}
where \ml{p$_1$, $\ldots$, p$_m$} are patterns, into
\begin{mlcoded}
    f $=$ \tl{}v$_1\ldots$ \tlb{v$_m$}(((\tl{}p$_{1}'\ldots$\tlb{p$_{m}'$}E$'$) v$_1 \ldots$ v$_m$) \fatbar{} ERROR)
\end{mlcoded}
where \ml{v$_1$, $\ldots$, v$_m$} are new variables that do not occur free in \ml{E}, the \ml{p$_{i}'$} are the results of translating the \ml{p$_i$}, and \ml{E$'$} is the result of translating \ml{E}. The only new complication is that we must specify what happens in case of failure. Suppose \ml{f} is applied to \ml{m} arguments, and the first pattern-match fails:
\begin{mlcoded}
    (\tl{}p$_{1}' \ldots$\tlb{p$_{m}'$}E$'$) E$_1$ E$_2 \ldots$ E$_m \rightarrow$ FAIL E$_2 \ldots$ E$_m$
\end{mlcoded}

Then we want the whole expression to fail, so we need to add a reduction rule for \ml{FAIL}:
\begin{mlcoded}
    FAIL E $\rightarrow$ FAIL
\end{mlcoded}
Now we can continue reduction.
\begin{mlcoded}
    \ml{FAIL} E$_2$ E$_3\ldots$ E$_m$ $\rightarrow$ \ml{FAIL} E$_3\ldots$ E$_m$ $\rightarrow \cdots \rightarrow$ FAIL
\end{mlcoded}
The translation is readily extended for the case when \ml{f} is defined by several equations. To see an example of this in action, consider the definition of \ml{xor} given above:
\begin{mlcoded}
    \begin{tabular}{lll}
    xor False &y &$=$ y\\
    xor True &False &$=$ True\\
    xor True &True &$=$ False
    \end{tabular}
\end{mlcoded}
Combining the rules of this section and the last allows us to transform this to
\rightline{\footnotesize(Notice that the arguments are matched from left to right)}
\vspace{-\baselineskip}
\begin{mlalign}
    xor $=$ \tlb{x}\tlb{y}&( ((\tlb{FALSE}\tlb{y}y) x y) \\
    & \fatbar{} ((\tlb{TRUE}\tlb{FALSE}TRUE) x y) \\
    & \fatbar{} ((\tlb{TRUE}\tlb{TRUE}FALSE) x y) \\
    & \fatbar{} ERROR)
\end{mlalign}


\subsection{Conditional Equations}

Next, we describe how to translate conditional equations into the enriched lambda calculus. Consider the following Miranda definition:
\begin{mlcoded}
    \begin{tabular}{lll}
        gcd a b &$=$ gcd (a $-$ b) b, &a $>$ b\\
        &$=$ gcd a (b $-$ a), &a $<$ b\\
        &$=$ a, &a $=$ b
    \end{tabular}
\end{mlcoded}
It is easy to see that the right-hand side of this definition could be translated to
\begin{mlcoded}
    (IF ($>$ a b) (gcd ($-$ a b) b)\\
    (IF ($<$ a b) (gcd a ($-$ b a))\\
    (IF ($=$ a b) a FAIL)))
\end{mlcoded}
Notice that if all the guards fail, then \ml{FAIL} is returned by the nested \ml{IF} expression. (In the case of \ml{gcd} this can never occur, and a very clever compiler might be able to discover this fact and optimize the last \ml{IF}.) In a more complicated definition, the failure of all the guards would cause the next equation to be tried (see example below).

Regarding all of an equation after the first \ml{$=$} sign as a `right-hand side', we can now give a new translation scheme, \metafn{TR}, which translates right-hand sides:

\plainbox{
    {\centering

    \metafnbb{TR}{rhs} translates the right-hand side of a definition

    }

    \metafn{TR}\ml{\hspace{-2em}
        \begin{minipage}{0.3\textwidth}
        \[
            \left[\hspace{-4pt}\left[
            \begin{array}{rl}
                & \text{A}_1, \text{G}_1\\
                =& \text{A}_2, \text{G}_2\\
                & \cdots \\
                =& \text{A}_n, \text{G}_n
            \end{array}
            \right]\hspace{-4pt}\right]
            \equiv
        \]
        \end{minipage}
    }
    \begin{minipage}{0.6\textwidth}
        \vs
      \begin{mlcoded}
          (IF \metafnbb{TE}{G$_1$} \metafnbb{TE}{a$_1$} \\
          (IF \metafnbb{TE}{G$_2$} \metafnbb{TE}{a$_2$} \\
          (IF \metafnbb{TE}{G$_n$} \metafnbb{TE}{a$_n$} FAIL)$\cdots$ ))
      \end{mlcoded}
    \end{minipage}\vs

    where \ml{A$_i$} is an expression and \ml{G$_i$} is a boolean-valued expression.
}


Now we can use \metafn{TR} instead of \metafn{TE} to translate the right-hand sides of function definitions. As an example, recall the definition of \ml{funnyLastElt}:
\begin{mlcoded}
    \begin{tabular}{ll}
    funnyLastElt (x:xs) &$=$ x, x < 0\\
    funnyLastElt (x:[]) &$=$ x\\
    funnyLastElt (x:xs) &$=$ funnyLastElt xs
    \end{tabular}
\end{mlcoded}
We can now translate it to
\begin{mlalign}
    funnyLastElt $=$ \tlb{v}&( ((\tlb{CONS x xs} . IF (< x 0) x FAIL) v) \\
    &\fatbar{} ((\tlb{CONS x NIL}x) v)  \\
    &\fatbar{} ((\tlb{CONS x xs}funnyLastElt xs) v)  \\
    &\fatbar{} ERROR)
\end{mlalign}
If the first equation matches, but the guard fails, then the \ml{IF} returns \ml{FAIL}, and the next equation is tried.

In Miranda, the final guard \ml{G$_n$} may be omitted, which is equivalent to giving a final guard of \ml{True}. In this case, the innermost \ml{IF} is of the form
\begin{mlcoded}
    IF TRUE E$_1$ FAIL
\end{mlcoded}
which can be optimized to
\begin{mlcoded}
    E$_1$
\end{mlcoded}
For example, the definition of \ml{factorial}
\begin{mlalign}
    factorial n &$=$ 1, \qquad n $=$ 0\\
    &$=$ n $*$ factorial (n-1)
\end{mlalign}
would be translated to
\begin{mlalign}
    factorial $=$ \tlb{v}&( ((\tlb{n}IF ($=$ n 0) 1 ($*$ n (factorial ($-$ n 1)))) v)\\
     &\fatbar{} ERROR)
\end{mlalign}
This can be simplified further, since the pattern-match cannot fail, and this special case will be spotted by the transformations of Chapter 5.

\subsection{Repeated Variables}

It appears at first that it is easy to use a conditional equation to eliminate repeated variables, by introducing a new variable name to replace one of the occurrences of the repeated variable, and adding an appropriate equality condition. For example, we could rewrite the definition of \ml{noDups} (given in Section 4.2.1) thus:

\begin{mlcoded}
\begin{tabular}{ll}
    noDups [\,] &$=$ [\,]\\
    noDups [x] &$=$ [x]\\
    noDups (x:y:ys) &$=$ noDups (y:ys), x $=$ y\\
    noDups (x:y:ys) &$=$ x : noDups (y:ys)
\end{tabular}
\end{mlcoded}

(The last two equations could now be combined into a conditional equation with two alternatives.) Unfortunately, this approach occasionally conflicts with the left-to-right rule originally given for pattern-matching. For example, given the following definition:

\begin{mlcoded}
\begin{tabular}{ll}
    nasty x x True &$=$ 1\\
    nasty x y z &$=$ 2
\end{tabular}
\end{mlcoded}
consider the evaluation of
\begin{mlcoded}
    nasty bottom 3 False
\end{mlcoded}
where the evaluation of \ml{bottom} fails to terminate (for example, \ml{bottom} could be defined by the degenerate equation: \ml{bottom $=$ bottom}). We might expect that the evaluation \ml{(nasty bottom 3 False)} would not terminate, since we will try to evaluate \ml{bottom} in order to compare it with \ml{3}. However, suppose we transformed the definition of \ml{nasty} to use a conditional equation:
\begin{mlcoded}
\begin{tabular}{ll}
    nasty$'$ x y True &$=$ 1,\qquad x $=$ y\\
    nasty$'$ x y z &$=$ 2
\end{tabular}
\end{mlcoded}
Now, if we evaluate \ml{(nasty$'$ bottom 3 False)}, \ml{bottom} will match \ml{x} and \ml{3} will match \ml{y}, but the match of \ml{True} against \ml{False} will fail, so the second equation will be tried, and deliver the answer 2. Hence, \ml{nasty} and \ml{nasty$'$} behave differently, and the transformation is invalid. (Note: \ml{nasty} and \ml{nasty$'$} also behave differently for expressions such as \ml{(nasty 1 2 bottom)}.)

There is a further complication raised by repeated variables. Consider the function \ml{multi}:
\begin{mlcoded}
\begin{tabular}{ll}
        multi p q q p $=$ 1\\
    multi p q r s $=$ 2
\end{tabular}
\end{mlcoded}

Should we compare the first and fourth arguments, and then compare the second and third arguments, or the other way around? The order of comparison is important, because it affects termination; consider \ml{(multi bottom 2 3 4)}.

This section has shown that repeated variables in a pattern are not as straightforward as at first appeared (the examples were suggested by Simon Finn of the University of Stirling). To simplify the rest of this chapter we will therefore sidestep these complications, by restricting our attention to a subset of Miranda which does not allow repeated variables in a pattern. We lose no expressive power thereby, though we do lose some notational convenience.

\subsection{Where-clauses}
Miranda allows the right-hand side of a definition to be qualified with a \ml{where}-clause. For example,
\begin{mlalign}
    sumsq x y $=$ xsq $+$ ysq&\\
    where\qquad{} &\\
    xsq &$=$ x $*$ x\\
    ysq &$=$ y $*$ y
\end{mlalign}
It is intuitively clear that this could be translated to
\begin{mlalign}
    sumsq $=$ \tlb{x}\tlb{y}(let &xsq $=$ $*$ x x\\
    &ysq $=$ $*$ y y\\
    in &\\
    & ($+$ xsq ysq))
\end{mlalign}
where we use a \ml{let}-expresion instead of a \ml{where}-clause. In general, the definitions in a \ml{where}-clause may be mutually recursive, so we have to use a \ml{letrec}-expresion instead. This will be optimized in Section 6.2.8.

Finally, the scope of a where-clause may include a set of alternatives and guards in a conditional equation:
\begin{mlcoded}
    \begin{tabular}{rll}
        gcd a b $=$ &gcd diff b, &a $>$ b\\
        $=$ &gcd a ($-$diff), &a $<$ b\\
        $=$ &a, & a $=$ b\\
        &where & \\
        & \qquad diff $=$ a $-$ b &
    \end{tabular}
\end{mlcoded}

% Figure 4.2 has been relocated to be with 4.3 and 4.4.

The scope of the definition of \ml{diff} includes all the alternatives and guards.
Figure 4.2 gives the final \metafn{TR} translation scheme, which translates right-hand sides, using a \ml{letrec} to translate a \ml{where}-clause.

\subsection{Patterns on the Left-hand Side of Definitions}

So far we have only described how to translate \textit{function} definitions, but Miranda also allows a \textit{pattern} to appear on the left-hand side of a definition. For example, consider the following Miranda definition:
\begin{mlalign}
    addPair w $=$ x $+$ y & \\
    where &(x, y) $=$ w
\end{mlalign}
The product pattern \ml{(x, y)} appears on the left-hand side of the definition in the \ml{where}-clause. It implies that \ml{w} evaluates to a pair, and it binds the names \ml{x} and \ml{y} to the components of \ml{w}.

As mentioned in Section 3.2.3, we also allow general patterns to appear on the left-hand side of definitions in a \ml{let(rec)}. This extension allows us to make a simple translation of \ml{addPair} to
\begin{mlcoded}
    addPair $=$ \tlb{w}(letrec (PAIR x y) $=$ w in ($+$ x y))
\end{mlcoded}

The hard work of dealing with patterns on the left-hand side of definitions is now carried out by transforming this \ml{letrec} into the ordinary lambda calculus, which is described in Section 6.2. The modification required to \metafn{TD} is very simple:
\begin{mlcoded}
    \metafnbb{TD}{ p $=$ R } $\equiv$ \metafnbb{TE}{p} $=$ \metafnbb{TR}{R}
\end{mlcoded}
where \ml{p} is a pattern and \ml{R} is a right-hand side.

\subsection{Summary}
We have now completed the development of the translation of a significant subset of Miranda into the enriched lambda calculus. The final translation schemes, summarized in Figures 4.2, 4.3 and 4.4, look rather forbidding, but this is because of their generality rather than their complexity.

\section{The Semantics of Pattern-matching Lambda Abstractions}
Having described how to translate from Miranda into a language involving pattern-matching lambda abstractions, we now give the semantics of pattern-matching lambda abstractions of the form \ml{( \tlb{p}E)}.

We will do so by devoting a subsection to each form of the pattern, \ml{p}: variable, constant, sum-constructor, and product-constructor.


\boxedfigure{
    {\centering

        \metafnbb{TR}{rhs} translates the right-hand side of a definition

    }

    \hspace{-1em}\metafn{TR}\ml{\hspace{-2em}
        \begin{minipage}{0.3\textwidth}
            \[
            \left[\hspace{-4pt}\left[
            \begin{array}{rl}
                & \text{A}_1, \text{G}_1\\
                =& \cdots \\
                =& \text{A}_n, \text{G}_n\\
                & \hspace{-1em}\text{\normalfont{where}}\\
                & \text{D}_1\\
                & \cdots \\
                & \text{D}_m\\
            \end{array}
            \right]\hspace{-4pt}\right]
            \equiv
            \]
        \end{minipage}
    }
    \begin{minipage}{0.6\textwidth}
        \vs
        \begin{mlcoded}
            \begin{tabular}{ll}
                letrec & \metafnbb{TD}{D$_1$}\\
                & $\cdots$ \\
                & \metafnbb{TD}{D$_m$}\\
                in & \\
                & (IF \metafnbb{TE}{G$_1$} \metafnbb{TE}{a$_1$} \\
                & (IF \metafnbb{TE}{G$_2$} \metafnbb{TE}{a$_2$} \\
                & (IF \metafnbb{TE}{G$_n$} \metafnbb{TE}{a$_n$} FAIL)$\cdots$ ))
            \end{tabular}
        \end{mlcoded}
    \end{minipage}\vs

    \noindent If \ml{G$_n$} is absent, or \ml{True}, then the final \ml{IF}-expression should be replaced by \metafnbb{TE}{A$_n$}.\vs

    \begin{tabular}{rl}
        where & \ml{A$_i$} is an expression\\
        & \ml{G$_i$} is a boolean-valued expression\\
        & \ml{D$_i$} is a definition
    \end{tabular}
}{The final \ml{TR} translation scheme}



\boxedfigure{
\begin{center}
    \metafnbb{TE}{Exp} translates the expression \ml{Exp}
\end{center}
\vspace{-0.5\baselineskip}
\begin{tabular}{ll}
    \metafnbb{TE}{k} &$\equiv$ \ml{k} \quad\quad (assumes no name-changing) \\
    \metafnbb{TE}{v} &$\equiv$ \ml{v} \\
    \metafnbb{TE}{ E$_1$ E$_2$ } &$\equiv$ \metafnbb{TE}{ E$_1$ } \metafnbb{TE}{ E$_2$ } \\
    \metafnbb{TE}{ E$_1$ infix E$_2$ } &$\equiv$ \metafnbb{TE}{ infix } \metafnbb{TE}{ E$_1$ } \metafnbb{TE}{ E$_2$ } \\
    \metafnbb{TE}{ E$_1$ \$v E$_2$ } &$\equiv$ \metafnbb{TE}{ v } \metafnbb{TE}{ E$_1$ } \metafnbb{TE}{ E$_2$ } \\
    \metafnbb{TE}{:}  &$\equiv$ CONS \\
    \metafnbb{TE}{[ ]}  &$\equiv$ NIL \\
    \metafnbb{TE}{[E$_1$, E$_2$, $\ldots$, E$_n$]}  &$\equiv$ CONS \metafnbb{TE}{E$_1$} \metafnbb{TE}{[E$_2$, $\ldots$, E$_n$]}\\
    \metafnbb{TE}{(E$_1$, E$_2$)}  &$\equiv$ PAIR \metafnbb{TE}{E$_1$} \metafnbb{TE}{E$_2$} \\
    \metafnbb{TE}{(E$_1$, E$_2$, E$_3$)}  &$\equiv$ TRIPLE \metafnbb{TE}{E$_1$} \metafnbb{TE}{E$_2$} \metafnbb{TE}{E$_3$}\\
    {\normalfont\normalsize and so on} & \\
    \metafnbb{TE}{True}  &$\equiv$ TRUE \\
    \metafnbb{TE}{False}  &$\equiv$ FALSE \\
\end{tabular}\\
\vs
\begin{tabular}{rll}
    \hspace{1cm} where & \ml{k} & is a literal constant or built-in operator \\
    & \ml{v} & is a variable \\
    & \ml{E$_1$} & is an expression \\
    & \ml{infix} & is an infix operator
\end{tabular}


}{The final \metafn{TE} translation scheme}


\begin{figure}[H]
\centering

{%
    \setlength{\fboxrule}{1pt}%
    \setlength{\fboxsep}{10pt}%
    \fbox{%
        \begin{minipage}{\textwidth}
            \small
            \setlength{\parindent}{0pt}
            \setlength{\parskip}{0mm plus 0mm minus 0mm}
\begin{center}
    \metafnbb{TD}{Def} translates the definition \ml{Def}
\end{center}
\begin{mlcoded}
    \metafnbb{TD}{p $=$ R} $\equiv$ \metafnbb{TE}{p} $=$ \metafnbb{TR}{R} \\
    \hspace{-1em}\metafn{TD}\hspace{-0.5em}
    \begin{minipage}{0.3\textwidth}
        \[
        \left[\hspace{-4pt}\left[
        \begin{array}{rl}
            \text{f p}_{1,1}  \ldots \text{p}_{1,m} &= \text{R}_1\\
            \vdots \qquad \vdots & \vdots \\
            \text{f p}_{n,1}  \ldots \text{p}_{n,m} &= \text{R}_n\\
        \end{array}
        \right]\hspace{-4pt}\right]
        \]
    \end{minipage}\vs

    \begin{minipage}{0.9\textwidth}
        \begin{mlalign}
            \hspace{-1em}$\equiv$ f $=$ (\tl{}v$_1$\ldots\tlb{v$_m$}&( ((\tl{}\metafnbb{TE}{p$_{1,1}$}\ldots\tlb{\metafnbb{TE}{p$_{1,m}$}}\metafnbb{TR}{R$_1$}) v$_1$ \ldots v$_m$)\\
            & \fatbar{} $\ldots$\\
            & \fatbar{} ((\tl{}\metafnbb{TE}{p$_{n,1}$}\ldots\tlb{\metafnbb{TE}{p$_{n,m}$}}\metafnbb{TR}{R$_n$}) v$_1$ \ldots v$_m$)\\
            & \fatbar{} ERROR))
        \end{mlalign}
    \end{minipage}
\end{mlcoded}

\begin{tabular}{rll}
    \hspace{1cm} where & \ml{f} & is a variable \\
    & \ml{v$_i$} & is a variable not free in any \ml{R$_j$}\\
    & \ml{p$_{i,j}$} & is a pattern\\
    & \ml{R} & is a right-hand side\\
    & \ml{R$_i$} & is a right-hand side
\end{tabular}
        \end{minipage}%
    }%
}%

\caption{\textsf The final \metafn{TD} translation scheme}
\end{figure}



\subsection{The Semantics of Variable Patterns}
If the pattern \ml{p} is a variable \ml{v}, then the pattern-matching lambda abstraction \ml{(\tlb{p}.E)} is just an ordinary lambda abstraction \ml{(\tlb{v}E)}, whose semantics have already been discussed in Section 2.5.

\subsection{The Semantics of Constant Patterns}
To describe the semantics of constant patterns we must specify the value of

\begin{mlcoded}
    \metafnbb{Eval}{\tlb{k}E}
\end{mlcoded}

where \ml{k} is a constant. Its value is certainly a function, so we can specify it by giving the value of

\begin{mlcoded}
    \metafnbb{Eval}{\tlb{k}E} a
\end{mlcoded}

for any argument \ml{a}. There are three possibilities: either \ml{a} is the same as \ml{k}, or it is \ml{1}, or it is something else. This leads to the following semantic equations:

\begin{mlcoded}
    \begin{tabular}{lll}
        \metafnbb{Eval}{\tlb{k}E} a &$=$ \metafnbb{Eval}{E} & \text{if} a $=$ \metafnbb{Eval}{k} \\
    \metafnbb{Eval}{\tlb{k}E} a &$=$ \text{FAIL} & \text{if} a $\neq$ \metafnbb{Eval}{k} \text{and} a $\neq$ $\bot$ \\
    \metafnbb{Eval}{\tlb{k}E} $\bot$ &$=$ $\bot$
    \end{tabular}
\end{mlcoded}

The first equation says that if \ml{(\tlb{k}E)} is applied to something that evaluates to \ml{k}, then the result comes from evaluating \ml{E}. The second equation says that the result is \ml{FAIL} if the argument evaluates to anything else, and the third equation specifies that, if the evaluation of the argument fails to terminate, then so does the whole application. As usual, these semantic equations specify reduction rules by implication. Thus, for example

\begin{mlalign}
    (\tlb{1}$+$ 3 4) 1 &$\rightarrow +$ 3 4 \\
    (\tlb{1}$+$ 3 4) 2 &$\rightarrow$ \text{FAIL}
\end{mlalign}

It is also possible to regard constants as sum-constructors of arity zero, as outlined in Section 4.1.2.3, in which case the rules of this section become a special case of those of the next.

\subsection{The Semantics of Sum-constructor Patterns}
Next, we consider the case of constructor patterns, of the form \ml{(s p$_1$ $\ldots$ p$_r$)}. Initially we will only consider sum patterns, since product patterns turn out to require special treatment. Here are the semantic rules for such patterns:

\begin{mlcoded}
    \footnotesize
    \begin{tabular}{ll}
        \metafnbb{Eval}{\tlb{(s p$_1$ $\ldots$ p$_r$)}E} (s a$_1$ $\ldots$ a$_r$) &$=$ \metafnbb{Eval}{\tlb{p$_1$ $\ldots$ p$_r$}E} a$_1$ $\ldots$ a$_r$ \\
    \metafnbb{Eval}{\tlb{(s p$_1$ $\ldots$ p$_r$)}E} (s$'$ a$_1$ $\ldots$ a$_r$) &$=$ \text{FAIL}  \text{if} \quad s $\neq$ s$'$ \\
    \metafnbb{Eval}{\tlb{(s p$_1$ $\ldots$ p$_r$)}E} $\bot$ &$=$ $\bot$
    \end{tabular}
\end{mlcoded}

Operationally, the rules work as follows. To apply \ml{(\tlb{(s p$_1$ $\ldots$ p$_r$)}E)} to an argument \ml{A} we first evaluate \ml{A} to find out what sort of object it is. This implies that if the evaluation of \ml{A} does not terminate then neither does the application in question (third rule). (Note: to `evaluate \ml{A}' we only evaluate it to constructor form; we do not evaluate its components. They will be evaluated lazily, only if they are extracted and used. This is what it means for constructors to be lazy.)

If \ml{A} evaluates to an object built with a constructor other than \ml{s}, then the pattern-match fails (second rule). To see how this rule works, consider an
application of the lambda abstraction \ml{(\tlb{(BRANCH t1 t2)}BRANCH t2 t1)} to \ml{(LEAF 0)}:
\begin{mlalign}
    (\tlb{(BRANCH t1 t2)}BRANCH t2 t1) (LEAF 0) &$\rightarrow$ \text{FAIL}
\end{mlalign}
The application returns \ml{FAIL} because the constructor in the pattern is different from that of the argument.

Finally, if \ml{A} was built with the same constructor as the pattern, then the first rule applies. To see how this rule works, consider an application of the same abstraction to a \ml{BRANCH}:
\begin{mlcoded}
    (\tlb{(BRANCH t1 t2)}BRANCH t2 t1) (BRANCH (LEAF 0) (LEAF 1)) \\
    $\rightarrow$ (\tlb{t1}\tlb{t2}BRANCH t2 t1) (LEAF 0) (LEAF 1) \\
    $\rightarrow$ (\tlb{t2}BRANCH t2 (LEAF 0)) (LEAF 1) \\
    $\rightarrow$ 1 \\
    $\rightarrow$ BRANCH (LEAF 1) (LEAF 0)
\end{mlcoded}
In this case the match succeeds, and \ml{t1} and \ml{t2} are bound to the components of the branch with the ordinary $\beta$-reduction rule.

Notice that for constructors of arity zero \ml{(r=0)} the three rules correspond exactly to those of the previous section. For example, using the first case of the \ml{xor} function gives:
\begin{mlalign}
    (\tlb{FALSE}\tlb{y}y) FALSE TRUE &$\rightarrow$ (\tlb{y}y) TRUE \\
    &$\rightarrow$ TRUE
\end{mlalign}
Finally, notice that the rules deal correctly with nested patterns. Consider, for example, the following application of the first case of the function \ml{lastElt} to \ml{(CONS 4 (CONS 3 NIL))}:
\begin{mlcoded}
    (\tlb{(CONS x NIL)}x) (CONS 4 (CONS 3 NIL))\\
    $\rightarrow$ (\tlb{x}\tlb{NIL}x) 4 (CONS 3 NIL) (\text{first rule}) \\
    $\rightarrow$ (\tlb{NIL}4) (CONS 3 NIL) (first rule)\\
    $\rightarrow$ \text{FAIL} \quad (\text{second rule})
\end{mlcoded}
Here, the outer pattern matches but the inner one does not, so the whole expression returns \ml{FAIL}.

\subsection{The Semantics of Product-constructor Patterns}
Finally we consider the semantics of matching product patterns. This is an area in which a rather subtle issue surfaces.

Consider the Miranda functions
\begin{mlcoded}
    \begin{tabular}{ll}
        zeroAny x &$=$ 0\\
        zeroList [\,] &$=$ 0\\
        zeroPair (x,y) &$=$ 0
    \end{tabular}
\end{mlcoded}
The function \ml{zeroAny} takes a single argument and returns \ml{0}. Miranda's lazy semantics clearly means that the argument is not evaluated, so that \ml{0} is
returned even if the evaluation of the argument is very expensive or non-terminating:
\begin{mlcoded}
    \metafnbb{Eval}{zeroAny} $\bot$ $=$ 0
\end{mlcoded}
We say that \ml{zeroAny} is \textit{lazy} since it does not evaluate its argument.

The semantics of the function \ml{zeroList} has already been described by the preceding sections. It specifies that \ml{zeroList} evaluates its argument, and checks whether it is \ml{[\,]}. If it is, then \ml{zeroList} returns \ml{0}, otherwise it returns \ml{ERROR}. We say that \ml{zeroList} is \textit{strict} since it does evaluate its argument:
\begin{mlalign}
    \metafnbb{Eval}{zeroAny} $\bot$ $=$ $\bot$
\end{mlalign}
Should the \ml{zeroPair} function be lazy or strict? Since the argument is a tuple there is no point in evaluating it to check that it really is a tuple, as was required in the case of \ml{zeroList}, because the check would always succeed (assuming that the program is type-checked). It would be more in the spirit of a lazy language to specify that
\begin{mlcoded}
    \metafnbb{Eval}{zeroAny} $\bot$ $=$ 0
\end{mlcoded}
and the Miranda language specifies this choice. We call this \textit{lazy product-matching}. On the other hand, an alternative choice would be to specify that
\begin{mlcoded}
    \metafnbb{Eval}{zeroAny} $\bot$ $=$ $\bot$
\end{mlcoded}
and we call this \textit{strict product-matching}.

Notice that there is no `right' or `wrong' answer; it is simply a question of making a clear choice of semantics for product-matching. The only `wrong' approach is not to notice that there is a choice to be made (and hence to risk making different choices in different parts of the implementation, with unpredictable results).

Nevertheless, we contend that there are persuasive arguments in favor of the lazy approach. We discuss this issue in the next section, while in the rest of this section we concentrate on the semantics of lazy product-matching.

We may describe lazy product-matching by the following semantic rule:

\begin{mlalign}
\metafnbb{Eval}{\tlb{(t p$_1$ $\ldots$ p$_r$)}E} a
$=$ \metafnbb{Eval}{\tl{}p$_1$ $\ldots$ \tlb{p$_r$}E} &(SEL-t-1 a)\\
& \qquad$\ldots$ \\
&(SEL-t-r \,a)
\end{mlalign}

Here \ml{SEL-t-i} is a built-in function which selects the $i^\text{th}$ field from a structured object built with constructor \ml{t}. It may be described by the following semantic equations:

\begin{mlcoded}
\begin{tabular}{ll}
    {SEL}-t-i (t a$_1$ $\ldots$ a$_i$ $\ldots$ a$_r$) &$=$ a$_i$ \\
    {SEL}-t-i $\bot$ &$=$ $\bot$
\end{tabular}
\end{mlcoded}
Suppose that \ml{(\tlb{p}E)}, where \ml{p} is a product pattern, is applied to an expression \ml{A}. The rule for lazy product-matching postpones the evaluation of the argument \ml{A} by binding the names for the components to applications of \ml{SEL-t-i} to \ml{A}, rather than evaluating \ml{A} and extracting its components directly. If none of the components of \ml{A} is evaluated, then \ml{A} will not be evaluated either, which is the effect we wanted to achieve.
Let us see how this works on \ml{zeroPair}:
\begin{mlcoded}
zeroPair = \tlb{(PAIR x y)}0
\end{mlcoded}
Hence,
\begin{mlalign}
  &\metafnbb{Eval}{zeroPair} $\bot$\\
 = &\metafnbb{Eval}{\tlb{(PAIR x y)}0} $\bot$\\
 = &\metafnbb{Eval}{\tlb{x}\tlb{y}0} (SEL-PAIR-1 $\bot$) (SEL-PAIR-2 $\bot$)\\
 = &\metafnbb{Eval}{\tlb{y}0} (SEL-PAIR-2 $\bot$)\\
 = &0
\end{mlalign}
as required.

\subsection{A Defense of Lazy Product-matching}

Consider the Miranda function \ml{firsts}, which takes a list of numbers, and
returns a pair consisting of the first odd and first even elements of the list:
\begin{mlcoded}
\begin{tabular}{ll}
firsts [\,] &= (0,0)\\
firsts (x:xs) &= combine x (firsts xs)
\end{tabular}

\begin{tabular}{llll}
combine x (od, ev) &= (x, ev),\qquad &odd &x\\
  &= (od, x), &even &x
\end{tabular}
\end{mlcoded}
Suppose that we were to use strict product-matching, so that when evaluating
an application \ml{(combine A$_1$ A$_2$)} we would first evaluate \ml{A$_2$}. Now consider
evaluating \ml{(firsts [1..])}, where \ml{[1..]} is the infinite list of integers starting at \ml{1}:
\begin{mlalign}
firsts [1..] &$\rightarrow$ combine 1 (firsts [2..])\\
 &$\rightarrow$ combine 1 (combine 2 (firsts [3..]))
\end{mlalign}
and so on.
The evaluation of \ml{(firsts [1..])} will never terminate. This is hardly satisfactory, because it is clear that the value of \ml{(firsts [1..])} should be \ml{(1,2)}.

All is well, however, if we use lazy product-matching. Then, in effect, the
evaluation goes like this:
\begin{mlalign}
firsts [1..] &$\rightarrow$ combine 1 (firsts [2..])\\
 &$\rightarrow$ (1, SEL-PAIR-2 (firsts [2..]))\\
 &$\rightarrow$ (1, SEL-PAIR-2 (combine 2 (firsts [3..])))\\
 &$\rightarrow$ (1, SEL-PAIR-2 (SEL-PAIR-1 (firsts [3..]), 2))\\
 &$\rightarrow$ (1, 2)
\end{mlalign}
Under lazy product-matching, combine does not evaluate its second
argument. Instead it binds \ml{od} to \ml{(SEL-PAIR-1 A)} and \ml{ev} to \ml{(SEL-PAIR-2 A)},
where \ml{A} is the argument.

We conclude that lazy product-matching gives significant benefits to the
programmer. The effect is quite subtle: strict product-matching caused the
entire argument list to be scanned even though all the operations on lists are
lazy. One purpose of this section is to point out that it is easy for a subtle
difference in evaluation strategy (strict versus lazy product-matching) to
cause a gross difference in the operational behavior of the program (scanning
the whole of an infinite list versus looking at the first element only). The
example is derived from a paper by Wadler [1985].

A further reason for advocating lazy product-matching is that it allows us to
describe mutual recursion correctly. For an explanation of this point, see
Section 6.2.6.

There is another interesting mathematical way of looking at the differences
between strict and lazy product-matching. In domain theory there is more
than one way of forming the product of two domains $A$ and $B$, that vary in their
treatment of $\bot$. The \textit{ordinary product}, $A\times B$, is defined like this:
\[
A \times B = \{ (a,b) \mid a\in A \text{ and } b\in B\}
\]
All the elements of this domain are pairs, and the bottom element of $A\times B$ is $(\bot, \bot)$.

The \textit{lifted product}, $(A \times B)_\bot$, is defined like this:
\[
(A \times B)_\bot = (A \times B) \cup \{\bot\}
\]
In this product the element $\bot$ is distinct from $(\bot, \bot)$. This corresponds closely to
our operational ideas of how tuples (or any other data structure) are formed: $\bot$
stands for a non-terminating computation, while $(\bot, \bot)$ is a pair, both of whose
elements are non-terminating computations.

The key insight is that lazy product-matching corresponds to ordinary
product, and strict product-matching corresponds to lifted product. To
implement the ordinary product domain $(A \times B)$ we have to make $(\bot, \bot)$
indistinguishable from non-termination. Since they clearly differ operation-
ally, the only way to conceal their differences is to use values in an ordinary
product domain in a way that makes them indistinguishable. This is precisely
what the lazy product-matching rule does:
\begin{mlalign}
&\metafnbb{Eval}{\tlb{(PAIR p$_1$ p$_2$)}E} $\bot$ \\
= &\metafnbb{Eval}{\tlb{p$_1$}\tlb{p$_2$}E} (SEL-PAIR-1 $\bot$) (SEL-PAIR-2 $\bot$)\\
= &\metafnbb{Eval}{\tlb{p$_1$}\tlb{p$_2$}E} $\bot$ $\bot$\\
\\
&\metafnbb{Eval}{\tlb{(PAIR p$_1$ p$_2$)}E} (PAIR $\bot$ $\bot$) \\
= &\metafnbb{Eval}{\tlb{p$_1$}\tlb{p$_2$}E} (SEL-PAIR-1 (PAIR $\bot$ $\bot$)) (SEL-PAIR-2 (PAIR $\bot$ $\bot$))\\
= &\metafnbb{Eval}{\tlb{p$_1$}\tlb{p$_2$}E} $\bot$ $\bot$
\end{mlalign}
In other words, the abstraction \ml{(\tlb{PAIR p$_1$ p$_2$}E)} is indifferent to whether its
argument is $\bot$ or $(\bot, \bot)$; it returns the same result in either case. So lazy
product-matching can be regarded as a way of implementing ordinary product
domains $(A \times B)$ by using the values in the lifted product domain $(A \times B)_\bot$ in
such a way that $(\bot, \bot)$ is indistinguishable from $\bot$.

Finally, it is worth noting that the use of lazy product-matching carries an
implementation cost. Consider a function \ml{addPair}, which adds together the
elements of a pair:
\begin{mlcoded}
addPair = \tlb{PAIR x y}+ x y
\end{mlcoded}

Now, using lazy product-matching, the reduction of \ml{(addPair (PAIR 3 4))} goes as follows:
\begin{mlalign}
    & addPair (PAIR 3 4) \\
    $=$ & (\tlb{(PAIR x y)}+ x y) (PAIR 3 4) \\
    $\rightarrow$ & (\tlb{x}\tlb{y}+ x y) (SEL-PAIR-1 (PAIR 3 4)) (SEL-PAIR-2 (PAIR 3 4)) \\
    $\rightarrow$ & (\tlb{y}+ (SEL-PAIR-1 (PAIR 3 4)) y) (SEL-PAIR-2 (PAIR 3 4)) \\
    $\rightarrow$ & + (SEL-PAIR-1 (PAIR 3 4)) (SEL-PAIR-2 (PAIR 3 4)) \\
    $\rightarrow$ & + 3 (SEL-PAIR-2 (PAIR 3 4)) \\
    $\rightarrow$ & + 3 4\\
    $\rightarrow$ & 7
\end{mlalign}

This takes one reduction to apply the \ml{addPair} lambda abstraction, and then two further reductions (subsequently) to reduce the two applications of \ml{SEL-PAIR}. Contrast this with the effect of using strict product-matching:
\begin{mlalign}
    & addPair (PAIR 3 4) \\
    $=$ & (\tlb{(PAIR x y)}+ x y) (PAIR 3 4) \\
    $\rightarrow$ & (\tlb{x}\tlb{y}+ x y) 3 4 \\
    $\rightarrow$ & (\tlb{y}+ 3 y) 4 \\
    $\rightarrow$ & + 3 4 \\
    $\rightarrow$ & 7
\end{mlalign}
This uses fewer reductions, since the application of the \ml{addPair} lambda abstraction also takes the argument apart. Furthermore, it uses less store since no temporary applications of \ml{SEL-PAIR} are constructed. This suggests that we should use strict product-matching instead of lazy product-matching wherever this does not affect the semantics.

In the case of \ml{addPair}, it is clear that the argument will certainly be evaluated in the end, so it would do no harm to evaluate it at the time of function application (that is, to use strict product-matching). In general, whenever a function is strict in an argument (see Section 2.5.4) it is safe to use strict product-matching for that argument. The process of working out which functions are strict is called \textit{strictness analysis}, and is discussed in Chapter 22.

\subsection{Summary}
This section has examined the semantics of pattern-matching in some detail, because much confusion has surrounded this area in the past. Figure 4.5 summarizes the results of the section. The distinction between strict and lazy product-matching, and the use of \ml{\fatbar} and \ml{FAIL}, are both first described in Turner's thesis [Turner, 1981], but the present formulation based on structured types is due to the authors.

\section{Introducing \ml{case}-expressions}
The transformations in the last section produce remarkably inefficient programs! The main reason for this is that pattern-matches are attempted, testing for \ml{FAIL} each time, as each equation in the function definition is tried in turn.

Frequently, however, a single test would suffice to select the appropriate equation. For example, recall again the \ml{reflect} function:
\begin{mlcoded}
    \begin{tabular}{ll}
    reflect (LEAF n) &$=$ LEAF n \\
    reflect (BRANCH t1 t2) &$=$ BRANCH (reflect t2) (reflect t1)
    \end{tabular}
\end{mlcoded}
To apply \ml{reflect}, it would suffice to test the argument, and select the first or second right-hand side according to whether it was a \ml{LEAF} or a \ml{BRANCH}.

In this section, therefore, we introduce \ml{case}-expressions, a convenient construct for describing a particularly simple form of pattern-matching which has this single-test property. Chapter 5 will then demonstrate how to translate Miranda function definitions into \ml{case}-expressions, and Chapter 6 will show how \ml{case}-expressions can be transformed into the ordinary lambda calculus. The net effect will be a significant improvement in the efficiency of the resulting program.

\ml{Case}-expressions are a notation for describing a simple form of pattern-matching. To begin with an example, we may translate the definition of \ml{reflect}, using a \ml{case}-expression, in the following way:
\begin{mlcoded}
    \begin{tabular}{llll}
    reflect $=$ \tlb{t} &case t of \\
    &LEAF n  &$\Rightarrow$  &LEAF n \\
    &BRANCH t1 t2 &$\Rightarrow$ &BRANCH (reflect t2) (reflect t1)
    \end{tabular}
\end{mlcoded}
The important points about a \ml{case}-expression are that the patterns are \textit{simple} (that is, not nested) and \textit{exhaustive} (that is, they cover all constructors of the type). This makes them particularly simple to implement.

The general form of a \ml{case}-expression is
\begin{mlcoded}
    \begin{tabular}{llll}
    case &v of \\
    &c$_1$ v$_{1, 1} \ldots $ v$_{1, r_1}$ &$\Rightarrow$ &E$_1$\\
    &$\ldots$ \\
    &c$_n$ v$_{n, 1} \ldots $ v$_{n, r_n}$ &$\Rightarrow$ &E$_n$\\
    \end{tabular}
\end{mlcoded}
where \ml{v} is a variable, \ml{E$_1\ldots$ E$_n$} are expressions, the \ml{v$_{i, j}$} are distinct variables, and the \ml{c$_1, \ldots,$ c$_n$} are a \textit{complete family} of constructors from a structured type declaration. The syntax of \ml{case}-expressions was defined in Figure 3.2.

Operationally, to evaluate this \ml{case}-expression, \ml{v} is first evaluated. Then, according to what constructor \ml{v} was built with, the appropriate \ml{E$_i$} is selected and evaluated, with the \ml{v$_{i, j}$} bound to the components of \ml{v}.

Formally, the construct is defined to be equivalent to
\begin{mlalign}
 & ((\tlb{c$_1$ v$_{1, 1} \ldots $ v$_{1, r_1}$}E$_1$) v) \\
 \fatbar{} & $\ldots$\\
 \fatbar{} & ((\tlb{c$_n$ v$_{n, 1} \ldots $ v$_{n, r_n}$}E$_n$) v)
\end{mlalign}
but a \ml{case}-expression is far more readable!

Intuitively, \ml{case}-expressions correspond to a multiway jump, whereas the equivalent expression using \ml{\fatbar} corresponds to a sequential `\ml{if$\ldots$then$\ldots$elseif$\ldots$}'


\begin{figure}[H]
    \centering

    {%
        \setlength{\fboxrule}{1pt}%
        \setlength{\fboxsep}{10pt}%
        \fbox{%
            \begin{minipage}{\textwidth}
                \small
                \setlength{\parindent}{10pt}
                \setlength{\parskip}{0mm plus 0mm minus 0mm}
                \setlength{\tabcolsep}{0.25em}
The semantic equations of \ml{(\tlb{p}E)} are:
\begin{mlcoded}
    \hspace{-2em}
    \begin{tabular}{lll}
        \metafnbb{Eval}{\tlb{k}E} & a $=$ \metafnbb{Eval}{E} & \normalfont{\qquad if} a $=$ \metafnbb{Eval}{k}\\
        \metafnbb{Eval}{\tlb{k}E} & a $=$ FAIL & \normalfont{\qquad if} a $\neq$ \metafnbb{Eval}{k} \normalfont{and} a $\neq \bot$  \\
        \metafnbb{Eval}{\tlb{k}E} & $\bot = \bot$
    \end{tabular}
\end{mlcoded}

\begin{mlcoded}
    \hspace{-2em}
    \begin{tabular}{ll}
        \metafnbb{Eval}{\tlb{(s p$_1$ $\ldots$ p$_{r_s}$)}E} (s a$_1$ $\ldots$ a$_{r_s}$ )
        &$=$ \metafnbb{Eval}{\tl{}p$_1$ $\ldots$ \tlb{p$_{r_s}$}E} a$_1$\!$\ldots$\! a$_{r_s}$\\
        \metafnbb{Eval}{\tlb{(s p$_1$ $\ldots$ p$_{r_s}$)}E} (s' a$_1$ $\ldots$ a$_{r_s}$ )
        &$=$ FAIL \quad \normalfont{if} s $\neq$ s' \\
        \metafnbb{Eval}{\tlb{(s p$_1$ $\ldots$ p$_{r_s}$)}E} $\bot$ & $= \bot$
    \end{tabular}
\end{mlcoded}

\begin{mlalign}
    \hspace{-1.5em}
    \metafnbb{Eval}{\tlb{(t p$_1$ $\ldots$ p$_r$)}E} a
    $=$ \metafnbb{Eval}{\tl{}p$_1$ $\ldots$ \tlb{p$_r$}E} &(SEL-t-1 a)\\
    & \qquad$\ldots$ \\
    &(SEL-t-r \,a)
\end{mlalign}

\hspace{-1.5em}
\begin{tabular}{rl}
    where & \ml{k} is a constant\\
    & \ml{s} is a sum constructor of arity \ml{r$s$}\\
    & \ml{t} is a product constructor of arity \ml{r$_t$}\\
    & \ml{P$_i$} is a pattern\\
    & \ml{E} is an expression\\
    & \ml{a$_i$}, \ml{a} are values
\end{tabular}
\vs

The \ml{SEL-t-i} functions are defined as follows:
\begin{mlcoded}
\begin{tabular}{ll}
SEL-t-i (t a$_1$ $\ldots$ a$_j$ $\ldots$ a$_r$) &$=$ a$_i$\\
SEL-t-i $\bot$ &$= \bot$
\end{tabular}
\end{mlcoded}
where \ml{t} is a product constructor of arity \ml{r}.
\vs

The \fatbar{} operator is defined as follows:
\begin{mlcoded}
\begin{tabular}{lll}
a & \fatbar{} b = a \qquad \qquad & \normalfont{if} a $\neq \bot$ \normalfont{and} a $\neq$ FAIL\\
FAIL & \fatbar{} b = b  \\
$\bot$ & \fatbar{} b = $\bot$
\end{tabular}
\end{mlcoded}

            \end{minipage}%
        }%
    }%

    \caption{\textsf Semantics of pattern-matching lambda abstractions \\and \fatbar{}}
\end{figure}

\section{Summary}
Structured data types have proved more complicated than at first appeared! We have discussed the background and semantics of pattern-matching, showing how to translate a Miranda function definition involving pattern-matching into the enriched lambda calculus. This required us to define two new constructs, pattern-matching lambda abstractions and the \fatbar{} operator, whose semantics we then defined. To clear the way for a more efficient translation, we then introduced \ml{case}-expressions, describing their semantics in terms of a transformation into the constructs previously described.

The next two chapters complete the pattern-matching story. Chapter 5 gives a more efficient translation of Miranda function definitions into \ml{case}-expressions, and Chapter 6 shows how to transform the new constructs into the ordinary lambda calculus.


\section*{References}

\begin{references}
    \item Burstall, R.M. 1969. Proving properties of programs by structural induction. \textit{The Computer Journal}. Vol. 12, No. 1, pp. 41-8.
    \item Burstall, R.M. 1977. Design considerations for a functional programming language. In \textit{Proceedings Infotech State of the Art Conference}, Copenhagen, pp. 54-7.
    \item Burstall, R.M., and Darlington, J. 1977. A transformation system for developing recursive programs. \textit{Journal of the ACM}. Vol. 24, No. 1, pp. 44-67.
    \item Burstall, R.M., and Goguen, J.A. 1982. \textit{Algebras, Theories, and Freeness: An Introduction for Computer Scientists}. Report CSR-101-82, Dept of Computer Science, University of Edinburgh. February.
    \item Landin, P.J. 1966. The next 700 programming languages. \textit{Communications of the ACM}. Vol. 9, No. 3, pp. 157-64.
    \item Turner, D.A., 1981. Aspects of the implementation of programming languages. D.Phil. thesis, University of Oxford. February.
    \item Wadler, P. 1985. \textit{A Splitting Headache -- and Its Cure}. Programming Research Group, Oxford. January.
\end{references}
\chapter[Efficient Compilation of Pattern-Matching][Efficient Compilation of Pattern-Matching]{Efficient Compilation of\\Pattern-Matching}
\chapterauthors{Philip Wadler}
\vspace{3cm}

 This chapter shows how to compile function definitions with pattern-matching into \ml{case}-expressions that can be efficiently evaluated. Previously, pattern-matching has been formally defined, and we have seen some examples of function definitions with pattern-matching.

\section{Introduction and Examples}
We begin by reviewing two examples.

The first example shows pattern-matching on more than one pattern. The function call \ml{(mappairs f xs ys)} applies the function \ml{f} to corresponding pairs from the lists \ml{xs} and \ml{ys}.
\begin{letalign}
    \ml{mappairs f [\,] ys} &$=$ \ml{[\,]} \\
    \ml{mappairs f (x:xs) [\,]} &$=$ \ml{[\,]} \\
    \ml{mappairs f (x:xs) (y:ys)} &$=$ \ml{f x y : mappairs f xs ys}
\end{letalign}
For example, \ml{(mappairs ($+$) [1,2] [3,4])} returns \ml{[4,6]}. The definition given here specifies that if the argument lists are not the same length, then the result will be as long as the shorter of the two lists. For example, \ml{(mappairs ($+$) [1,2] [3,4,5])} also returns \ml{[4,6]}.

The simplest way to think of pattern-matching is as trying to match each equation in turn. Within each equation, patterns are matched from left to right. For example, evaluating \ml{(mappairs ($+$) [1,2] [3,4])} first matches \ml{($+$)} against \ml{f} in the first equation, which succeeds, and then matches \ml{[1,2]} against
\ml{[\,]}, which fails. Then the second equation is tried. Matching \ml{($+$)} against \ml{f} and \ml{[1,2]} against \ml{(x:xs)} both succeed, but matching \ml{[3,4]} against \ml{[\,]} fails. Finally, matching in the third equation succeeds, binding \ml{f} to \ml{($+$)}, \ml{x} to \ml{1}, \ml{xs} to \ml{[2]}, \ml{y} to \ml{3}, and \ml{ys} to \ml{[4]}. This corresponds exactly to the way pattern-matching was defined in Chapter 4.

Performing pattern-matching in this way can require a lot of work. The example above had to examine the list \ml{[1,2]} three times and the list \ml{[3,4]} twice. It seems clear that it should be possible to evaluate this function application in a more efficient manner that examines each list only once, but still gives the result prescribed by the semantics. This can be done by transforming the above definition into an equivalent one using \ml{case}-expressions:

\begin{mlcoded}
    \setlength{\tabcolsep}{2pt}
    mappairs\\
    \begin{tabular}{lllll}
    $=$ \tlb{f} &\tlb{xs$'$} \tlb{ys$'$} & & & \\
    & \ml{case xs$'$ of} & & & \\
    & \ml{NIL} &$\Rightarrow$ &\ml{NIL} & \\
    & \ml{CONS x xs} &$\Rightarrow$ &\ml{case ys$'$ of} & \\
    & & &\ml{NIL} &$\Rightarrow$ \ml{NIL} \\
    & & &\ml{CONS y ys} &$\Rightarrow$ \ml{CONS (f x y)\! (mappairs\! f\! xs\! ys)}
    \end{tabular}
\end{mlcoded}
(\ml{Case}-expressions were introduced in Section 4.4.) This chapter describes an algorithm that can automatically translate the first definition into the second. This algorithm is called the pattern-matching compiler.

The second example shows pattern-matching on a nested pattern. The function call \ml{(nodups xs)} removes adjacent duplicate elements from a list \ml{xs}. It can be defined as follows:

\begin{mlcoded}
\begin{tabular}{lll}
    nodups [\,] &$=$ [\,] &\\
    nodups [x] &$=$ [x] & \\
    nodups (y:x:xs) &$=$ \ml{nodups (x:xs)}, &\ml{y $=$ x} \\
     &$=$ y : \ml{nodups (x:xs)}, &\normalfont{otherwise}
\end{tabular}
\end{mlcoded}

(As you would expect, the guard `otherwise' applies if no other guard does. See Appendix) For example, \ml{(nodups [3,3,1,2,2,2,3])} returns \ml{[3,1,2,3]}. Note that the naming need not be consistent: \ml{x} stands for the first element of the list equation in the second equation, and for the second element of the list in the third.

Again, one can apply this definition by matching each equation in turn. For example, evaluation of \ml{(nodups [1,2,3])} will first try to match \ml{[1,2,3]} against \ml{[\,]}, which fails. Next, it will try to match \ml{[1,2,3]} against \ml{[x]}, which also fails. Finally, it will succeed in matching \ml{[1,2,3]} against \ml{(y:x:xs)}, binding \ml{y} to \ml{1}, \ml{x} to \ml{2}
and \ml{xs} to \ml{[3]}. Again, this corresponds exactly to the semantics in Chapter 4.

As before, this is not very efficient. The list \ml{[1,2,3]} is examined three times,
and the sublist \ml{[2,3]} is examined twice (once in the second equation, where it
fails to match \ml{[\,]}, and once in the third equation, where it succeeds in matching
\ml{(x:xs)}). The pattern-matching compiler can transform this into a form that
examines the list and the sublist only once:
\begin{mlcoded}
    \setlength{\tabcolsep}{2pt}
    \begin{tabular}{llllll}
    \tlb{nodups}&&&&& \\
    $=$ \tlb{xs$''$} &case xs$''$ & of &&&\\
    &NIL     &  &$\Rightarrow$ NIL && \\
    &CONS x$'$ &xs$'$  &$\Rightarrow$ && \\
    & &case &xs$'$ of && \\
    & &NIL & &$\Rightarrow$ & CONS x$'$ NIL \\
    & &CONS &x xs & $\Rightarrow$ & IF ($=$ x$'$ x)\\
    &&&&& \quad (nodups (CONS x xs)) \\
    &&&&& \quad (CONS x$'$ (nodups (CONS x xs))
    \end{tabular}
\end{mlcoded}
(Here \ml{x$'$} is the variable that was called \ml{x} in the second equation and \ml{y} in the third.)

The two kinds of pattern-matching, nested patterns and multiple patterns, are closely related to one another. The pattern-matching compiler discussed below works uniformly for both.

In the examples above, the patterns on the left-hand sides of the equations do not overlap. Many people would rewrite the first definition in the form:
\begin{letalign}
    mappairs$'$ f [\,] ys &$=$ \ml{NIL} \\
    mappairs$'$ f xs [\,] &$=$ \ml{NIL} \\
    mappairs$'$ f (x:xs) (y:ys) &$=$ f x y : mappairs$'$ f xs ys
\end{letalign}
In this case, the patterns overlap because both the first and the second equation match against \ml{(mappairs$'$ f [\,] [\,])}.

One reason for preferring \ml{mappairs$'$} to \ml{mappairs} is that it is considered to be more efficient. Indeed, if the simplest implementation of pattern-matching is used, matching each equation in turn, then it is slightly less work to match against \ml{xs} than to match against \ml{(x:xs)}. However, as we shall see, this definition may actually be \textit{less} efficient when the pattern-matching compiler is used. Some other problems with definitions like \ml{mappairs$'$} will be discussed in Section 5.5.

The remainder of this chapter is organized as follows. Section 5.2 explains the pattern-matching compiler algorithm. Section 5.3 presents a Miranda program that implements the algorithm. Section 5.4 describes some optimizations to the pattern-matching compiler. Section 5.5 discusses a restricted class of definitions, called uniform definitions, which have useful properties.

Credit for the first published description of a pattern-matching compiler goes to Augustsson, who used it in the LML compiler [Augustsson, 1985]. Techniques similar to Augustsson's have been discovered independently by several researchers, including the authors of the Hope compiler [Burstall \textit{et al.}, 1980]. The material presented here is derived partly from Augustsson's paper and partly from original work by the author (Wadler).

It is also possible to derive the pattern-matching compiler from its specification using program transformation techniques; see Barrett and Wadler [1986].

\section{The Pattern-matching Compiler Algorithm}
A Miranda function definition of the form
\begin{mlcoded}
    f p$_{1,1}\ldots$ p$_{1,n} =$ E$_1$\\
    $\ldots$\\
    f p$_{m,1}\ldots$ p$_{m,n} =$ E$_m$
\end{mlcoded}
can be translated into the enriched lambda calculus definition
\begin{mlalign}
    f $=$ \tl{}u$_1 \ldots$ \tlb{u$_n$}&((\tlb{p$_{1,1'}$} $\ldots$\tlb{p$_{1,n'}$} E$_{1'}$) u$_1 \ldots$ u$_n$)\\
    \fatbar{} &$\ldots$\\
    \fatbar{} &((\tlb{p$_{m,1'}$} $\ldots$\tlb{p$_{m,n'}$} E$_{m'}$) u$_1 \ldots$ u$_n$)\\
    \fatbar{} &ERROR
\end{mlalign}
where the \ml{u$_i$} are new variables which do not occur free in any \ml{E$_i$}, and the \ml{E$_i'$} and \ml{p$_{i,j'}$} are the result of translating the \ml{E$_i$} and \ml{p$_{i,j}$}, respectively. It was shown how to do this translation in Chapter 4, using the \metafn{TD} translation scheme.

This section shows how to transform the definition of \ml{f} into a form which uses \ml{case}-expressions, removing all use of pattern-matching lambda abstractions. The transformation applies to the entire body of the \ml{\tl{}u$_1 \ldots$ \tl{}u$_n$} abstraction, except that we generalize slightly to allow an arbitrary expression instead of \ml{ERROR}.

For the sake of simplicity, we assume that constant patterns have been replaced by conditional equations, as described in Section 4.2.1.

\subsection{The Function match}

Our goal, then, is to transform an expression of the form
\begin{mlnumbered}
    \begin{tabular}{rl}
        &((\tlb{p$_{1,1}$} $\ldots$\tlb{p$_{1,n}$} E$_{1}$) u$_1 \ldots$ u$_n$)\\
    \fatbar{} &$\ldots$\\
    \fatbar{} &((\tlb{p$_{m,1}$} $\ldots$\tlb{p$_{m,n}$} E$_{m}$) u$_1 \ldots$ u$_n$)\\
    \fatbar{} &E
    \end{tabular}
\end{mlnumbered}
into an equivalent expression which uses \ml{case}-expressions rather than pattern-matching lambda abstractions.

The transformation is a bit complicated, and so we will use some new notation to describe it. Specifically, we will use a function \metafn{match}, which takes as its arguments the various parts of the input expression, namely the \ml{p$_{i,j}$}, \ml{E$_i$}, and \ml{u$_j$}, and produces as its output the transformed expression. The function \metafn{match} is similar to the \metafn{TD} and \metafn{TE} translation schemes introduced in Chapter 3, except that both its input and its result are enriched lambda calculus expressions. Furthermore, the double square bracket syntax becomes somewhat cumbersome, so we use a syntax like Miranda instead.

Here, then, is the call to \metafn{match} which we will use to compile the expression
(5.1) given above:
\begin{mlalign}
    \metafn{meta} &[u$_1$, $\ldots$, u$_n$]\\
    &[( [p$_{1,1}$, $\ldots$, p$_{1,n}$], E$_1$ )\\
    & $\ldots$\\
    &( [p$_{m,1}$, $\ldots$, p$_{m,n}$], E$_m$ )]\\
    &E
\end{mlalign}
This call should return an expression equivalent to the expression (5.1), and we take (5.1) as the \textit{definition} of \metafn{match} from a semantic point of view. A call of \metafn{match} takes three arguments: a list of variables, a list of equations, and a default expression. Each equation is a pair, consisting of a list of patterns (representing the left-hand side of the equation) and an expression (representing the right-hand side). Notice that the list of variables and each list of patterns have the same length.

We will also sometimes write calls of \metafn{match} in the form
\begin{mlcoded}
    \metafn{match} us qs E
\end{mlcoded}
Here \ml{us} is the list of argument variables (of length \ml{n}), and \ml{qs} is a list of equations (of length \ml{m}). Each equation \ml{q$_i$} in \ml{qs} has the form \ml{(ps$_i$, E$_i$)}, where \ml{ps$_i$} is the list of patterns on the left-hand side (of length \ml{n}) and \ml{E$_i$} is the expression on the right-hand side.

As a running example, we will use the following Miranda function:
\begin{letalign}
    demo f [\,] ys &$=$ A f ys\\
    demo f (x:xs) [\,] &$=$ B f x xs\\
    demo f (x:xs) (y:ys) &$=$ C f x xs y ys
\end{letalign}
This function is similar in structure to \ml{mappairs}, but it has been changed slightly in order to simplify and clarify the following examples. The right-hand sides use three unspecified expressions \ml{A}, \ml{B}, and \ml{C}.

Translating this into the enriched lambda calculus using \metafn{TD} gives:
\begin{letalign}
        \ml{demo}& \\
    $=$ \tlb{u$_1$}\tlb{u$_2$}\tlb{u$_3$}\!\! &((\tlb{f}\tlb{NIL}\tlb{ys}A f ys) u$_1$ u$_2$ u$_3$) \\
    \hfill\fatbar{} &((\tlb{f}\tlb{(CONS\! x\! xs)}\tlb{NIL}B\! f\! x\! xs) u$_1$ u$_2$ u$_3$) \\
    \hfill\fatbar{} &((\tlb{f}\tlb{(CONS\! x\! xs)}\tlb{(CONS\! y\! ys)}C\! f\! x\! xs\! y\! ys)\! u$_1$\! u$_2$\! u$_3$) \\
    \hfill\fatbar{} &ERROR
\end{letalign}
where \ml{u$_1$}, \ml{u$_2$}, \ml{u$_3$} are new variable names which do not occur free in \ml{A}, \ml{B}, or \ml{C}.
Now, we transform the definition of \ml{demo}, by replacing its body with a call of \metafn{match}:
\begin{mlcoded}
    \setlength{\tabcolsep}{0.25em}
    \begin{tabular}{lllll}
        \ml{demo}&&&& \\
        $=$ \tlb{u$_1$}\tlb{u$_2$}\tlb{u$_3$}\metafn{match} &[u$_1$, u$_2$, u$_3$] &&&\\
        &[\,(\,[f,\! NIL, &ys &],\! (A\! f\! ys) &),\! \\
        &[\,(\,[f,\! CONS\! x\! xs,\! &NIL &],\! (B\! f\! x\! xs)&), \\
        &[\,(\,[f,\! CONS\! x\! xs,\! &CONS\! y\! ys\! &],\! (C\! f\! x\! xs\! y\! ys)&)] \\
        &ERROR &&&
    \end{tabular}
\end{mlcoded}

The following sections give rules to transform any call of \metafn{match} to an
equivalent \ml{case}-expression. We begin with rules for simple cases and proceed to more general cases.

\subsection{The Variable Rule}

In the example above, we have the following call on match:
\begin{mlcoded}
    \setlength{\tabcolsep}{0.25em}
    \begin{tabular}{lllll}
        \metafn{match} &[u$_1$, u$_2$, u$_3$] &&&\\
        &[\,(\,[f,\! NIL, &ys &],\! (A\! f\! ys) &),\! \\
        &[\,(\,[f,\! CONS\! x\! xs,\! &NIL &],\! (B\! f\! x\! xs)&), \\
        &[\,(\,[f,\! CONS\! x\! xs,\! &CONS\! y\! ys\! &],\! (C\! f\! x\! xs\! y\! ys)&)] \\
        &ERROR &&&
    \end{tabular}
\end{mlcoded}
In this case, the list of patterns in every equation begins with a variable. This may be reduced to the equivalent call:
\begin{mlcoded}
    \setlength{\tabcolsep}{0.25em}
    \begin{tabular}{lllll}
        \metafn{match} &[u$_2$, u$_3$] &&&\\
        &[\,(\,[NIL, &ys &],\! (A\! u$_1$\! ys) &),\! \\
        &[\,(\,[CONS\! x\! xs,\! &NIL &],\! (B\! u$_1$\! x\! xs)&), \\
        &[\,(\,[CONS\! x\! xs,\! &CONS\! y\! ys\! &],\! (C\! u$_1$\! x\! xs\! y\! ys)&)] \\
        &ERROR &&&
    \end{tabular}
\end{mlcoded}
This is derived by removing the first variable, \ml{u$_1$}, and in each equation removing the corresponding formal variable, \ml{f}, and replacing \ml{f} by \ml{u}, in the right-hand side of each equation.

The same method works whenever each equation begins with a variable, even if each equation begins with a different variable. For example,
\begin{mlcoded}
        \setlength{\tabcolsep}{0.25em}
    \begin{tabular}{lll}
        \metafn{match} &[u$_2$, u$_3$]& \\
    &[ ( [x, NIL], &(B x) ),\\
    &\phantom{[ }( [y, CONS x xs], &(C y x xs) ) ]\\
    &ERROR
    \end{tabular}
\end{mlcoded}
reduces to the call,
\begin{mlcoded}
        \setlength{\tabcolsep}{0.25em}
    \begin{tabular}{lll}
        \metafn{match} &[u$_3$] &\\
    &[ ( [NIL], &(B u$_2$) ),\\
    &\phantom{[ }( [CONS x xs], &(C u$_2$ x xs) ) ]\\
    &ERROR
    \end{tabular}
\end{mlcoded}
(This particular example arises when compiling the definition of \ml{nodups}.)

In general, if every equation begins with a variable pattern, then the call of match will have the form:
\begin{letalign}
    \metafn{match} (u:us) &[ ( (v$_1$:ps$_1$), E$_1$ ), \\
    &\phantom{[ }\ldots \\
    &\phantom{[ }( (V$_m$:ps$_m$), E$_m$ ) ] \\
    &E
\end{letalign}
This can be reduced to the equivalent call:
\begin{letalign}
    \metafn{match} &us \\
    &[ ( ps$_1$, E$_1$[u$/$v$_1$] ), \\
    &\phantom{[ }( ps$_m$, E$_m$[u$/$v$_m$] ) ] \\
    &E
\end{letalign}
where, as usual, \ml{E[M$/$x]} means `\ml{E} with \ml{M} substituted for \ml{x}'. In order to avoid too many subscripts, a Miranda-like notation has been used here; for example, we write \ml{(u:us)} instead of \ml{[u$_1$,$\ldots$, u$_n$]}. The general case corresponds to the first example above, where \ml{u} is \ml{u$_1$}, \ml{us} is \ml{[u$_2$, u$_3$]}, \ml{v$_1$} is \ml{f}, \ml{ps$_1$} is \ml{[NIL, ys]}, and so on.

It is not hard to show that the rule is correct, that is, that the two \metafn{match} expressions are equivalent. This follows from the definition of \metafn{match} and the semantics of pattern-matching.

\subsection{The Constructor Rule}
The above step has left us with the following call of \metafn{match}:
\begin{mlcoded}
    \setlength{\tabcolsep}{0.25em}
    \begin{tabular}{llll}
        \metafn{match} &[u$_2$, u$_3$]  && \\
    &[ ( [NIL, ys&], (A u$_1$ ys) &), \\
    &\phantom{[ }( [CONS x xs, NIL&], (B u$_1$ x xs) &), \\
    &\phantom{[ }( [CONS x xs, CONS y ys&], (C u$_1$ x xs y ys) &) ] \\
    &ERROR
    \end{tabular}
\end{mlcoded}
In this case, the list of patterns in every equation begins with a constructor. This call is equivalent to the following \ml{case}-expression:
\begin{mlcoded}
    \setlength{\tabcolsep}{0.25em}
    \begin{tabular}{llllll}
        case u$_2$ of &&&&& \\
    \quad NIL &$\Rightarrow$ &\metafn{match} &[u$_3$] && \\
    &&&[\,(\,[ys], &(A u$_1$ ys) &)\,] \\
    &&&ERROR &&\\
    \quad CONS\! u$_4$\! u$_5$ &$\Rightarrow$ &\metafn{match} &[u$_4$, u$_5$, u$_3$] &&\\
    &&&[\,(\,[x, xs, NIL], &(B u$_1$ x xs) &), \\
    &&&\phantom{[\,}(\,[x, xs, CONS y ys], &(C u$_1$ x xs y ys) &)\,] \\
    &&&ERROR
    \end{tabular}
\end{mlcoded}
This call is derived by grouping together all equations that begin with the same constructor. Within each group, new variables are introduced corresponding to each field of the constructor. Thus \ml{NIL}, which has no fields, requires no new variables, while \ml{CONS}, which has two fields, introduces the variables \ml{u$_4$} and \ml{u$_5$}. These new variables are matched against the corresponding subpatterns of the original patterns.

It may be useful here to look at a second example. In compiling the definition of a function like \ml{nodups}, one would encounter the following call of \metafn{match}:
\begin{mlcoded}
    \setlength{\tabcolsep}{0.25em}
    \begin{tabular}{llll}
        \metafn{match} &[u$_1$]  && \\
        &[ ( [NIL&], A &), \\
        &\phantom{[ }( [CONS x NIL&], (B x) &), \\
        &\phantom{[ }( [CONS y (CONS x xs)&], (C y x xs) &) ] \\
        &ERROR
    \end{tabular}
\end{mlcoded}
This can be reduced to the equivalent expression:
\begin{mlcoded}
    \setlength{\tabcolsep}{0.25em}
    \begin{tabular}{llllll}
        case u$_2$ of &&&&& \\
        \quad NIL &$\Rightarrow$ &\metafn{match} &[\,] && \\
        &&&[\,(\,[\,], &A &)\,] \\
        &&&ERROR &&\\
        \quad CONS\! u$_2$\! u$_3$ &$\Rightarrow$ &\metafn{match} &[u$_2$, u$_3$] &&\\
        &&&[\,(\,[x, NIL], &(B x) &), \\
        &&&\phantom{[\,}(\,[y, CONS x xs], &(C y x xs) &)\,] \\
        &&&ERROR
    \end{tabular}
\end{mlcoded}
Again, \ml{NIL} introduces no new variables (leaving a call of \metafn{match} with an empty list of variables), and \ml{CONS} introduces two new variables, \ml{u$_2$} and \ml{u$_3$}.

More generally, it may be the case that not all equations beginning with the same constructor appear next to each other. For example, one might have a call of \metafn{match} such as:
\begin{mlcoded}
    \setlength{\tabcolsep}{0.25em}
    \begin{tabular}{llll}
        \metafn{match} &[u$_1$]  && \\
        &\phantom{[ }( [CONS x NIL&], (B x) &), \\
        &[ ( [NIL&], A &), \\
        &\phantom{[ }( [CONS y (CONS x xs)&], (C y x xs) &) ] \\
        &ERROR
    \end{tabular}
\end{mlcoded}
It is always safe to exchange two equations that begin with a different constructor, so we may rearrange the above to the equivalent call:
\begin{mlcoded}
    \setlength{\tabcolsep}{0.25em}
    \begin{tabular}{llll}
        \metafn{match} &[u$_1$]  && \\
        &[ ( [NIL&], A &), \\
        &\phantom{[ }( [CONS x NIL&], (B x) &), \\
        &\phantom{[ }( [CONS y (CONS x xs)&], (C y x xs) &) ] \\
        &ERROR
    \end{tabular}
\end{mlcoded}
which may be transformed as before.

It may also be the case that not all constructors appear in the original list of equations. For example, a function definition such as:
\begin{letalign}
        last [x] &$=$ x \\
        last (y:(x:xs)) &$=$ last (x:xs)
\end{letalign}
will result in the following call of \metafn{match}:
\begin{mlcoded}
    \setlength{\tabcolsep}{0.25em}
    \begin{tabular}{llll}
        \metafn{match} &[u$_1$]  && \\
        &[ ( [CONS x NIL&], x &), \\
        &\phantom{[ }( [CONS y (CONS x xs)&], (last (CONS x xs)) &) ] \\
        &ERROR
    \end{tabular}
\end{mlcoded}
This can be reduced to the equivalent expression:
\begin{mlcoded}
    \setlength{\tabcolsep}{0.25em}
    \begin{tabular}{llllll}
        case u$_1$ of &&&&& \\
        \quad NIL &$\Rightarrow$ &\metafn{match} &[\,] [\,] ERROR && \\
        \quad CONS\! u$_2$\! u$_3$ &$\Rightarrow$ &\metafn{match} &[u$_2$, u$_3$] &&\\
        &&&[\,(\,[x, NIL], &x &), \\
        &&&\phantom{[\,}(\,[y, CONS x xs], &(last (CONS x xs)) &)\,] \\
        &&&ERROR
    \end{tabular}
\end{mlcoded}
The \ml{case}-expression must still contain a clause for the missing constructor, and the call of \metafn{match} in this clause will have an empty list of equations. (From the definition of \metafn{match}, we know that \ml{(\metafn{match} [\,] [\,] ERROR)} is equivalent to \ml{ERROR}.)

We now discuss the general rule for reducing a call of \metafn{match} where every equation begins with a constructor pattern. Say that the constructors are from a type which has constructors \ml{c$_1$,$\ldots$, c$_1$}. Then the equations can be rearranged into groups of equations \ml{qs$_1$,$\ldots$, qs$_1$}, such that every equation in group \ml{qs$_i$}, begins with constructor \ml{c$_i$}. (If there is some constructor \ml{c} that begins no equation, like \ml{NIL} in the last example above, then the corresponding group \ml{qs$_i$} will be empty.) The call of \metafn{match} will then have the form:
\begin{mlcoded}
    \metafn{match} (u:us) (qs$_1$ $++ \ldots ++$ qs$_k$) E
\end{mlcoded}
where each \ml{qs$_i$}, has the form:
\begin{mlalign}
    [ &( ((c$_i$ ps$'_{i,1}$):ps$_{i,1}$), E$_{i,1}$ )\\
    & $\ldots$\\
    &( ((c$_i$ ps$'_{i,m_i}$):ps$_{i,m_i}$), E$_{i,m_i}$ ) ]
\end{mlalign}
($++$ is list append.) In this expression we have abbreviated the constructor pattern \ml{(c p$_1 \ldots$ p$_r$)} to the form \ml{(c ps)}, where \ml{ps} stands for the list of patterns \ml{[p$_1$, p$_2$,$\ldots$, p$_r$]}. This call to \metafn{match} is reduced to the \ml{case}-expression:
\begin{mlalign}
    case u of&\\
    \qquad c$_1$ us$'_1$ &$\Rightarrow$ \metafn{match} (us$'_1$, $++$ us) qs$'_1$ E\\
    \qquad$\ldots$ &\\
    \qquad c$_k$ us$'_k$ &$\Rightarrow$ \metafn{match} (us$'_k$, $++$ us) qs$'_k$ E\\
\end{mlalign}
where each \ml{qs$'_i$} has the form:
\begin{mlalign}
    [ &( (ps$'_{i,1}$ $++$ ps$_{i,1}$), E$_{i,1}$ ),\\
    & $\ldots$\\
    &( (ps$'_{i,m_i}$ $++$ ps$_{i,m_i}$), E$_{i,m_i}$ ) ]
\end{mlalign}
Here each \ml{us$'_i$} is a list of new variables, containing one variable for each field in \ml{c$_i$}.

For instance, in the example at the beginning of this section, \ml{qs$_2$} is
\begin{mlcoded}
    \setlength{\tabcolsep}{0.1em}
    \begin{tabular}{rlll}
        [ &( [CONS x xs, NIL&], (B u$_1$ x xs) &)\\
    &( [CONS x xs, CONS y ys&], (C u$_1$ x xs y ys) &)\,]
    \end{tabular}
\end{mlcoded}
and \ml{c$_2$} is \ml{CONS}, \ml{ps$'_{2,1}$} is \ml{[x, xs]}, \ml{ps$_{i,1}$} is \ml{[NIL]}, \ml{E$_{2,1}$} is \ml{(B u$_1$ x xs)}, \ml{ps$'_{2,2}$} is \ml{[x, xs]}, \ml{ps$_{2,2}$} is \ml{[CONS y ys]}, and \ml{E$_{2,2}$} is \ml{(C u$_1$ x xs y ys)}. The corresponding \ml{qs$'_2$} is
\begin{mlcoded}
    \setlength{\tabcolsep}{0.1em}
    \begin{tabular}{rlll}
        [ &( [x xs, NIL&], (B u$_1$ x xs) &)\\
        &( [x xs, CONS y ys&], (C u$_1$ x xs y ys) &)\,]
    \end{tabular}
\end{mlcoded}
The corresponding list of new variables, \ml{us$'_2$}, is \ml{[u$_4$,u$_5$]}.

This notation is, of necessity, rather clumsy. The reader will be pleased to discover, in Section 5.3, that this transformation can be written as a functional program which is more concise and (with experience) easier to read.

Again, the correctness of this rule can be proved using the definition of \metafn{match} and the semantics of pattern-matching.

\subsection{The Empty Rule}

After repeated application of the rules above, one eventually arrives at a call of \metafn{match} where the variable list is empty, such as the following:
\begin{mlalign}
    \metafn{match} &[\,] \\
    &[ ( [\,], (A u$_1$ u$_3$) ) ] \\
    &ERROR
\end{mlalign}
This reduces to:
\begin{mlcoded}
    (A u$_1$ u$_3$)
\end{mlcoded}
The correctness of this follows immediately from the definition of \metafn{match}, since \ml{A} cannot return \ml{FAIL}.

In general, the call of \metafn{match} may involve zero, one or more equations. Zero equations may result if the constructor rule is applied and some constructor of the type appears in no equations, as in \ml{last} above. More than one equation can result if some of the original equations overlap.

Thus, the general form of a call of \metafn{match} with an empty variable list is:
\begin{mlcoded}
    \setlength{\tabcolsep}{0.1em}
    \begin{tabular}{rll}
        match &[\,]& \\
    &[ &( [\,], E$_1$ ), \\
    && $\ldots$\\
    &&( [\,], E$_m$ ) ]\\
    &E
    \end{tabular}
\end{mlcoded}
where \ml{m $\geq$ 0}. From the definition of \metafn{match}, this reduces to
\begin{mlcoded}
    E$_1$ \fatbar{} $\ldots$ \fatbar{} E$_m$ \fatbar{} E
\end{mlcoded}
Further, we can often guarantee that none of \ml{E$_1$},$\ldots$, \ml{E$_m$} can be equal to \ml{FAIL}. In this case, the above \metafn{match} expression reduces to \ml{E$_1$} if \ml{m > 0} and to \ml{E} if \ml{m $=$ 0}. Section 5.4.2 discusses this optimization further.

\subsection{An Example}

The rules given so far are sufficient to translate the definitions of \ml{mappairs} and \ml{nodups} to the corresponding \ml{case}-expressions given in the introduction. Notice that the variable names used in the introduction were chosen for readability. In practice, the translation algorithm will usually pick new names. The reader may wish to verify that the rules given above are indeed sufficient to translate the definition
\begin{letalign}
    mappairs f [\,] ys & $=$ [\,]\\
    mappairs f (x:xs) [\,]  & $=$ [\,]\\
    mappairs f (x:xs) (y:ys) & $=$ f x y : mappairs f xs ys
\end{letalign}
to the equivalent:
\begin{mlcoded}
    \setlength{\tabcolsep}{0.25em}
    \begin{tabular}{llllll}
        mappairs &&&&&\\
        $=$ \tlb{u$_1$}\tlb{u$_2$}\tlb{u$_3$} &&&&&\\
        \quad case u$_2$ of &&&&& \\
        \quad NIL &$\Rightarrow$ &NIL &&& \\
        \quad CONS\! u$_4$\! u$_5$ &$\Rightarrow$ &case u$_3$ of &&&\\
        &&NIL &$\Rightarrow$ &NIL & \\
        &&CONS\! u$_6$\! u$_7$ &$\Rightarrow$ & CONS\! &(u$_1$\! u$_4$\! u$_6$) \\
        &&&&&(mappairs u$_1$\! u$_5$\! u$_7$) \\
    \end{tabular}
\end{mlcoded}
The reader may also wish to check that the function \ml{nodups} transforms to the \ml{case}-expression given in the introduction.

\subsection{The Mixture Rule}
The above rules are sufficient for compiling most function definitions into \ml{case}-expressions. However, there is still one case which has not been covered. This arises when not all equations begin with a variable, and not all equations begin with a constructor; that is, when there is a mixture of both kinds of equation. For example, here is an alternative definition of \ml{demo} (similar in structure to the alternative definition of \ml{mappairs}):
\begin{letalign}
    demo$'$ f [\,] ys &$=$ A f ys \\
    demo$'$ f xs [\,] &$=$ B f xs \\
    demo$'$ f (x:xs) (y:ys) &$=$ C f x xs y ys
\end{letalign}
Converting this to a \metafn{match} expression and applying the variable rule to eliminate \ml{f} results in the following:
\begin{mlcoded}
    \setlength{\tabcolsep}{0.25em}
    \begin{tabular}{lllll}
        \metafn{match} &[u$_2$, u$_3$]  &&& \\
        &[ ( [NIL, &ys&], (A u$_1$ ys) &), \\
        &\phantom{[ }( [xs, &NIL&], (B u$_1$ xs) &), \\
        &\phantom{[ }( [CONS x xs, &CONS y ys&], (C u$_1$ x xs y ys) &) ] \\
        &ERROR
    \end{tabular}
\end{mlcoded}
Neither the variable rule nor the constructor rule applies to this expression, because some equations begin with constructors and others with variables.

This is where the third argument to the \metafn{match} function is useful. The above expression is equivalent to:
\begin{mlcoded}
    match [U2, U3]
    ([NIL, ys], (\tlb{(A u: ys)})) \\
    ( match [U2, U3]
    [([xs, NIL], (\tlb{(B u: xs)}))] \\
    ( match [U2, U3]
    [([CONS x xs, CONS y ys], (\tlb{(C uf x xs y ys)}))] \\
    ERROR ))
\end{mlcoded}

\begin{mlcoded}
    \setlength{\tabcolsep}{0.25em}
    \begin{tabular}{lllll}
        \metafn{match} &[u$_2$, u$_3$]  &&& \\
        &[ ( [NIL, &ys&], (A u$_1$ ys) &), \\
        &\phantom{[ }( [xs, &NIL&], (B u$_1$ xs) &), \\
        &\phantom{[ }( [CONS x xs, &CONS y ys&], (C u$_1$ x xs y ys) &) ] \\
        &ERROR
    \end{tabular}
\end{mlcoded}

That is, the equations are broken into groups; first an equation beginning with a constructor, then one beginning with a variable, and then one beginning with a constructor again. If the equation in the first call of \metafn{match} fails to match the arguments then the value of the second call of \metafn{match} is returned. Similarly, if the equation in the second call does not match then the third call is returned, and if the equation in the third call does not match then \ml{ERROR} is returned.

The reader may verify that reducing the three calls of \metafn{match} using the variable, constructor and base case rules results in the following definition of \ml{mappairs$'$}:
\begin{mlcoded}
        \setlength{\tabcolsep}{0.25em}
demo$'$\\
$=$ \tlb{u$_1$}\tlb{u$_2$}\tlb{u$_3$} \\
\hspace{3em}
\begin{minipage}{1.5in}
    \begin{letalign}
case u$_2$ of & \\
NIL & $\Rightarrow$ (A u$_1$ u$_3$)\\
CONS u$_4$ u$_5$ & $\Rightarrow$ \\
\end{letalign}
\end{minipage}\\
\phantom{.\hspace{3em}.}
\begin{minipage}{1.5in}
    \begin{letalign}
        case u$_3$ of & \\
        NIL & $\Rightarrow$ (B u$_1$ u$_2$)\\
        CONS u$_6$ u$_7$ & $\Rightarrow$ \\
    \end{letalign}
\end{minipage}\\
\phantom{.\hspace{8em}.}
\begin{minipage}{1.5in}
    \begin{letalign}
        case u$_2$ of & \\
        NIL & $\Rightarrow$ ERROR\\
        CONS u$_4$ u$_5$ & $\Rightarrow$ \\
    \end{letalign}
\end{minipage}\\
\phantom{.\hspace{13em}.}
\begin{minipage}{1.5in}
    \begin{letalign}
        case u$_3$ of & \\
        NIL & $\Rightarrow$ ERROR\\
        CONS u$_6$ u$_7$ & $\Rightarrow$ (C u$_1$ u$_4$ u$_5$ u$_6$ u$_7$)\\
    \end{letalign}
\end{minipage}
\end{mlcoded}
This involves four \ml{case}-expressions. When the second and third arguments are both non-empty lists then each list is examined twice, as compared with once for the definition of demo. This confirms the claim made in the introduction that `optimizing' the definition of \ml{mappairs} by transforming it into \ml{mappairs$'$} can actually result in worse code.

It may be possible to devise a compilation algorithm that would produce better code for this case. This could be done by simplifying a \ml{case}-expression that appears inside another \ml{case}-expression for the same variable. This sort of optimization is straightforward, although it requires considerably more book-keeping. In this case, \ml{mappairs$'$} would compile to the same \ml{case}-expression as \ml{mappairs}, although the compilation process would be rather more complicated.


In general, a call of \metafn{match} where some equations begin with variables and some with constructors may be transformed as follows. Say we are given a call of \metafn{match} of the form
\begin{mlcoded}
    \metafn{match} us qs E
\end{mlcoded}
The equation list \ml{qs} may be partitioned into $k$ lists \ml{qs$_1$,$\ldots$, qs$_k$} such that
\begin{mlcoded}
    qs $=$ qs$_1 ++ \ldots ++$ qs$_k$
\end{mlcoded}
The partition should be chosen so that each \ml{qs$_i$} either has every equation beginning with a variable or every equation beginning with a constructor. (In the example above, each \ml{qs$_i$} had length 1, but in general this need not be the case.) Then the call of \metafn{match} can be reduced to:
\begin{mlcoded}
    \metafn{match} us qs$_1$ (\metafn{match} us qs$_2$ ($\ldots$ (\metafn{match} us qs$_k$ E)$\ldots$))
\end{mlcoded}
It is easy to use the definition of \metafn{match} to show that this rule is correct.

\subsection{Completeness}
With the addition of the mixture rule, it is now possible to reduce any possible call of \metafn{match} to a \ml{case}-expression. This can be seen by a simple analysis. Given a call \ml{(\metafn{match} us qs E)} then \ml{us} will be either empty, so the empty rule applies, or non-empty. If \ml{us} is non-empty then each equation must have a non-empty pattern list, which must begin with either a variable or a constructor. If all equations begin with a variable then the variable rule applies; if all begin with a constructor then the constructor rule applies; and if some begin with variables and some with constructors then the mixture rule applies.

Further, define the `size' of an equation list as the sum of the sizes of all the patterns in the equation list. It can be seen that all four of the rules result in calls of \metafn{match} with smaller equation lists. This guarantees that the algorithm must eventually terminate.

\section{The Pattern-matching Compiler in Miranda}
This section presents the transformation algorithm as a functional program in Miranda.

\subsection{Patterns}
First, it is necessary to give a data type for representing patterns.
\begin{letalign}
    pattern &\hastype$=$ VAR variable\\
    &\quad \quad $\mid$ CON constructor [pattern]\\
    variable &$==$ [char]\\
    constructor &$==$ [char]
\end{letalign}
For example, \ml{(x:xs)} is represented by \ml{(CON ``CONS'' [VAR ``x'', VAR ``xs''])}.

We need two functions on constructor names. The function \ml{arity} given a constructor returns its arity, and the function \ml{constructors} given a constructor returns a list of all constructors of its type:
\begin{letalign}
    arity &\hastype{} constructor $\rightarrow$ num \\
    constructors &\hastype{} constructor $\rightarrow$ [constructor]
\end{letalign}
For example, \ml{(arity ``NIL'')} returns \ml{0}, and \ml{(arity ``CONS'')} returns \ml{2}. Both \ml{(constructors ``NIL'')} and \ml{(constructors ``CONS'')} return the list \ml{[``NIL'', ``CONS'']}.

\subsection{Expressions}
Next, we need a data type for representing expressions:
\begin{letalign}
    expression &\hastype$=$ CASE variable [clause] \\
    &\quad \quad $\mid$ FATBAR expression expression \\
    &\quad \quad $\mid$ $\ldots$ \\
    clause &\hastype$=$ CLAUSE constructor [variable] expression
\end{letalign}
For example, the \ml{case}-expression:
\begin{letalign}
    case xs of &\\
    \quad NIL & $\Rightarrow$ E$_1$\\
    \quad CONS y ys & $\Rightarrow$ E$_2$\\
\end{letalign}
would be represented by
{    \setlength{\tabcolsep}{0.25em}
\begin{letalign}
    CASE & ``xs'' \\
 \hfill       [ & CLAUSE ``NIL'' [\,] E$_1'$, \\
         & CLAUSE ``CONS'' [``y'', ``ys''] E$'_2$]
\end{letalign}
}
\noindent where \ml{E$_1'$}, \ml{E$_2'$} are the representations of the expressions \ml{E$_1$}, \ml{E$_2$}. Similarly,
\begin{mlcoded}
    E$_1$ \fatbar{} E$_2$
\end{mlcoded}
would be represented by
\begin{mlcoded}
    FATBAR E$_1'$ E$_2'$
\end{mlcoded}

The \ml{$\ldots$} in the definition of the type expression stands for other constructors used to represent other expressions, such as variables, applications and lambda abstractions. We do not need to know anything about these other expressions, except that there is a substitution function defined for them.
\begin{mlcoded}
    subst \hastype{} expression $\rightarrow$ variable $\rightarrow$ variable $\rightarrow$ expression
\end{mlcoded}
For example, if \ml{E} represents the expression \ml{(f x y)}, then \ml{(subst E\, ``\_u1'' ``x'')} represents the expression \ml{(f\, \_u1 y)}.


\subsection{Equations}

An equation is a list of patterns paired with an expression:
\begin{mlcoded}
    equation $==$ ([pattern], expression)
\end{mlcoded}
We will use the letter \ml{q} to denote equations, or else write \ml{(ps,e)}.

We need functions to determine if an equation begins with a variable or a constructor. If it begins with a constructor, we also need a function to return that constructor.
\begin{mlcoded}
\setlength{\tabcolsep}{0.25em}
\begin{tabular}{lll}
        isVar &&\hastype{} equation $\rightarrow$ bool \\
    isVar (VAR v &: ps, e) &$=$ True \\
    isVar (CON c ps$'$ &: ps, e) &$=$ False\\

    isCon &&\hastype{} equation $\rightarrow$ bool \\
    isCon q &&$=$ $\sim$ (isVar q)\\

    getCon &&\hastype{} equation $\rightarrow$ constructor \\
    getCon (CON c ps$'$ &: ps, e) &$=$ c\\
\end{tabular}
\end{mlcoded}

\subsection{Variable Names}

We need some way of generating the new variable names, \ml{u1}, \ml{u2}, and so on. To do this we introduce a function \ml{makeVar} that, given a number, returns a variable name.
\begin{letalign}
    makeVar &\hastype{} num $\rightarrow$ variable \\
    makeVar k &$=$ ``\_u'' $++$ show k
\end{letalign}
For example, \ml{(makeVar 3)} returns `\ml{\_u3}'. Here we preface each new variable name with `\ml{\_}' to avoid it being confused with any variable already in the program.

\subsection{The Functions \ml{partition} and \ml{foldr}}

The implementation of the mixture rule uses a function called \ml{partition}. The call \ml{(partition f xs)} returns a list \ml{[xs$_1$,$\ldots$, xs$_n$]} such that \ml{xs $=$ xs$_1$ $++$$\ldots$ $++$ xs$_n$}, and such that \ml{f x $=$ f x$'$} for any elements \ml{x} and \ml{x$'$} in \ml{xs$_i$}, \ml{i} from \ml{1} to \ml{n}, and such that \ml{f x $\neq$ f x$'$} for any elements \ml{x} in \ml{xs$_i$} and \ml{x$'$} in \ml{xs$_{i+1}$}, \ml{i} from \ml{1} to \ml{n-1}. For example,

\begin{mlcoded}
    \ml{partition odd [1,3,2,4,1] $=$ [ [1,3], [2,4], [1] ]}
\end{mlcoded}

The function \ml{partition} is defined as follows:
\begin{letalign}
    partition &\hastype{} ($*$ $\rightarrow$ $**$) $\rightarrow$ [$*$] $\rightarrow$ [ [$*$] ] \\
    partition f [\,] &$=$ [\,] \\
    partition f [x] &$=$ [ [x] ] \\
    partition f (x:x$'$:xs) &$=$ tack x (partition f (x$'$:xs)), f x $=$ f x$'$ \\
     &$=$ [x] : partition f (x$'$:xs),\qquad otherwise\\
     \\
    tack x xss &$=$ (x : hd xss) : t xss
\end{letalign}

Incidentally, the following definition of tack is not equivalent to the above
\begin{mlcoded}
    tack x (xs:xss) $=$ (x : xs) : xss
\end{mlcoded}
The difference between the two is closely related to the question of strict and
lazy pattern-matching, mentioned in Section 4.3.5 in connection with the
function \ml{firsts}.

The pattern-matching compiler also uses the standard function \ml{foldr}. The
function \ml{foldr} is defined so that
\begin{mlcoded}
    foldr f a [x$_1$, x$_2$, $\cdots$, x$_n$] $=$ f x$_1$ (f x$_2$ ( $\cdots$ (f x$_n$ a) $\cdots$ ))
\end{mlcoded}
For example, \ml{(foldr ($+$) 0 xs)} returns the sum of the list of numbers \ml{xs}. The
function \ml{foldr} is defined by:
\begin{letalign}
    foldr &\hastype{} $(* \rightarrow ** \rightarrow **) \rightarrow ** \rightarrow [*] \rightarrow **$\\
    foldr f a [\,] &$=$ a\\
    foldr f a (x:xs) &$=$ f x (foldr f a xs)
\end{letalign}

\subsection{The Function \metafn{match}}

We are now ready to define the function \metafn{match}. Calls of \metafn{match} have the form
\ml{(\metafn{match} k us qs def)}. Here, as in Section 5.2, \ml{us} represents a list of variables,
\ml{qs} represents a list of equations and \ml{def} is a default expression. The argument
\ml{k} is added to help in generating new variable names; it should be chosen so
that for every \ml{i>k}, \ml{(makeVar i)} is a new variable not in \ml{us}, \ml{qs} or \ml{def}.

For example, the initial call to \metafn{match} to compile the definitions of \ml{mappairs} would be:
\begin{mlcoded}
    \setlength{\tabcolsep}{0.25em}
    \begin{tabular}{lllll}
        \metafn{match} &3 &&&\\
    &[``\_u1'', ``\_u2'', &``\_u3''] && \\
    &[ ( [VAR ``f'', &CON ``NIL'' &[\,], &\\
    && VAR ``ys''&&], E$_1$ ), \\
    &\phantom{[ }( [VAR ``f'', &CON ``CONS'' &[VAR ``x'', VAR ``xs''], &\\
    && CON ``NIL'' &[\,] &], E$_2$ ), \\
    &\phantom{[ }( [VAR ``f'', &CON ``CONS'' &[VAR ``x'', VAR ``xs''], &\\
    && CON ``CONS'' &[VAR ``y'', VAR ``ys''] &], E$_3$ ) ] \\
     &error
    \end{tabular}
\end{mlcoded}
where \ml{E$_1$}, \ml{E$_2$} and \ml{E$_3$} represent the three expressions on the right-hand sides of the equation, and \ml{error} represents the expression \ml{ERROR}.

The definition of \metafn{match} can now be derived in a fairly straightforward way from the description given in Section 5.2. The type of \metafn{match} is:
\begin{mlcoded}
    \metafn{match} \hastype{} num $\rightarrow$ [variable] $\rightarrow$ [equation] $\rightarrow$ expression $\rightarrow$ expression
\end{mlcoded}
The equations for the top-level of \metafn{match} come from the empty rule and the mixture rule.
\begin{letalign}
    \metafn{match} k [\,] &qs def $=$ def\\
    \metafn{match} k (u:us) &qs def $=$ foldr FATBAR def [e $\mid$ ([\,], e) $\leftarrow$ qs ]\\
    \hfill $=$ foldr &(matchVarCon k (u:us)) def (partition isVar qs)
\end{letalign}

The function \ml{matchVarCon} is given a list of equations that either all begin with a variable or all begin with a constructor. It calls \ml{matchVar} or \ml{matchCon}, as appropriate.
\begin{letalign}
    matchVarCon k us qs def \\
    \quad $=$ matchVar k us qs def, &\qquad if isVar (hd qs) \\
    \quad $=$ matchCon k us qs def, &\qquad if isCon (hd qs)
\end{letalign}

The function \ml{matchVar} implements the variable rule.
\begin{mlcoded}
    matchVar k (u:us) qs def \\
    \quad $=$ match k us [(ps, subst e u v) $\mid$ (VAR v : ps, e) $\leftarrow$ qs] def
\end{mlcoded}

The functions \ml{matchCon} and \ml{matchClause} implement the constructor rule. The call \ml{(choose c qs)} returns all equations that begin with constructor \ml{c}.
\begin{mlcoded}
    matchCon k (u:us) qs def \\
    \phantom{.\quad.} $=$ CASE u [matchClause c k (u:us) (choose c qs) def $\mid$ c $\leftarrow$ cs] \\
    \phantom{.\quad\quad.} where\\
    \phantom{.\quad\quad.} cs = constructors (getCon (hd qs))
\end{mlcoded}
\begin{letalign}
    matchClause c k (u:us) qs &def \\
    \phantom{.\quad.} $=$ CLAUSE c us$'$ (match &(k$'$ + k) \\
    &(us$'+\!+$us)\\
    &[(ps$'+\!+$ps, e)\! $\mid$\! (CON c ps$'$\! :\! ps, e)\,$\leftarrow$\,qs] \\
    &def)
\end{letalign}
\vspace{-\baselineskip}
\begin{letalign}
    \phantom{.\quad\quad.} where\\
    \phantom{.\quad\quad.} k$'$ &$=$ arity c \\
    \phantom{.\quad\quad.} us$'$ &$=$ [makeVar (i+k) $\mid$ i $\leftarrow$ [1..k$'$] ]
\end{letalign}
\begin{mlcoded}
choose c qs $=$ [q $\mid$ q $\leftarrow$ qs; getCon q $=$ c]
\end{mlcoded}

This completes the Miranda program for the pattern-matching compiler.

\section{Optimizations}
This section discusses some optimizations to the pattern-matching compiler. Section 5.4.1 describes an optimization which gives greater efficiency when
compiling overlapping equations. This involves further uses of \fatbar{} and \ml{FAIL}, and
Section 5.4.2 describes how these may often be eliminated.

\subsection{Case-expressions with Default Clauses}
If overlapping equations are allowed, then sometimes the pattern-matching compiler described above may transform a small set of equations into a \ml{case}-expression that is much larger. For example, consider the function
\begin{mlcoded}
\begin{tabular}{ll}
    unwieldy [\,] [\,] &$=$ A \\
    unwieldy xs ys &$=$ B xs ys
\end{tabular}
\end{mlcoded}
The pattern-matching compiler transforms this into:
\begin{mlcoded}
\setlength{\tabcolsep}{0.2em}
\begin{tabular}{lllll}
    unwieldy $=$ \tlb{xs}\tlb{ys}&case xs of &&& \\
    &NIL &$\Rightarrow$ case& ys of & \\
    &&&NIL &$\Rightarrow$ A \\
    &&&CONS y$'$ ys$'$ &$\Rightarrow$ B xs ys\\
    &CONS x' xs' &$\Rightarrow$ B xs& ys
\end{tabular}
\end{mlcoded}
Here the expression \ml{(B xs ys)} appears twice. If \ml{(B xs ys)} were replaced by
a very large expression, the increase in size caused by the compilation process
could be very significant.

The problem can be avoided by modifying the rules given in Section 5.2 so
that right-hand sides are never duplicated during the compilation process. In
fact, only one rule can cause right-hand sides to be duplicated, the constructor
rule. This rule is modified as follows.

Recall that the constructor rule transforms a call of \metafn{match} of the form:
\begin{mlcoded}
    \metafn{match} (u:us) (qs$_1$ $++$$ \ldots$ $++$ qs$_k$) E
\end{mlcoded}
to a \ml{case}-expression of the form:
\begin{letalign}
    case u of\\
    \phantom{.\quad.}c$_1$ us$'_1$ &$\Rightarrow$ \metafn{match} (us$'_1$ $++$ us) qs$'_1$ E \\
    \phantom{.\quad.}$\cdots$\\
    \phantom{.\quad.}c$_k$ us$'_k$ &$\Rightarrow$ \metafn{match} (us$'_k$ $++$ us) qs$'_k$ E \\
\end{letalign}
where \ml{qs$_1$, $\ldots$, qs$_k$} and \ml{qs$'_1$, $\ldots$, qs$'_k$} are as described in Section 5.2.3.

Normally \ml{E} will be \ml{ERROR}, but if the mixture rule is used then \ml{E} may itself
be a \metafn{match} expression containing right-hand sides; it is in this case that
duplication may occur. The modified rule prevents this by using \fatbar{} and \ml{FAIL} to
avoid duplicating \ml{E}.


This is done by replacing the \ml{case}-expression above with the equivalent expression:
\begin{letalign}
    case u of\\
    \phantom{.\quad.}c$_1$ us$'_1$ &$\Rightarrow$ \metafn{match} (us$'_1$ $++$ us) qs$'_1$ E \\
    \phantom{.\quad.}$\cdots$\\
    \phantom{.\quad.}c$_k$ us$'_k$ &$\Rightarrow$ \metafn{match} (us$'_k$ $++$ us) qs$'_k$ E \\
    \fatbar{} E
\end{letalign}
If we call the old \ml{case}-expression \ml{C}, then the new expression is \ml{(C$'$ \fatbar{} E)},
where \ml{C$'$} is formed by replacing each \ml{E} in \ml{C} by \ml{FAIL}. It is clear that the new
expression is equivalent to the old expression and, as desired, \ml{E} is not
duplicated by the new rule.

For example, using the new rule, the definition of \ml{unwieldy} will now  transform to:
\begin{mlcoded}
\setlength{\tabcolsep}{0.15em}
\begin{tabular}{llllllr}
unwieldy &$=$ \tlb{xs&}\tlb{ys} &&&&\\
    &(case &xs of&&&& \\
        &&NIL  & $\Rightarrow$ (case& ys of &&\\
        &&&&NIL & $\Rightarrow$ A & \\
        &&&&CONS y$'$ ys$'$  &$\Rightarrow$ FAIL) &(a) \\
    &&& \phantom{$\Rightarrow$ }\fatbar{} FAIL&& &(b) \\
    &&CONS x$'$ xs$'$ &$\Rightarrow$ FAIL)&&& \\
    &\fatbar{} B xs &ys &&&&(c) \\
\end{tabular}
\end{mlcoded}
This expression is a little larger than the previous version of \ml{unwieldy}, but now \ml{(B xs ys)} appears only once. If \ml{(B xs ys)} stands for a large expression, then this new expression may be much smaller than the previous one.

As an example of how this sort of expression is evaluated, consider the call \ml{(unwieldy NIL (CONS 1 NIL))}. This is evaluated as follows. First, the outer \ml{case}-expression is evaluated. Since \ml{xs} is \ml{NIL}, this causes the inner \ml{case} to be evaluated. Since \ml{ys} is \ml{(CONS 1 NIL)}, the inner \ml{case}-expression returns \ml{FAIL} see line (a). So the expression after the inner \ml{\fatbar{}} is returned, which is also \ml{FAIL} see line (b). Thus, the outer \ml{case}-expression returns \ml{FAIL}. So the expression after the outer \ml{\fatbar{}} is returned; see line (c). This is \ml{(B NIL (CONS 1 NIL))}, which is the value returned by the call of \ml{unwieldy}.

\subsection{Optimizing Expressions Containing \ml{\fatbar{}} and \ml{FAIL}}
It is often the case that all occurrences of \ml{FAIL}, and its companion, \ml{\fatbar{}}, can be eliminated. Most of these optimizations depend on reasoning that \ml{FAIL} can never be returned by an expression, because in this case an occurrence of \ml{\fatbar{}} can be eliminated.

Suppose that \ml{FAIL} is returned by an expression \ml{E}. Then it is necessary (though not sufficient) that one of the following conditions must hold:
\begin{numbered}
    \item \ml{FAIL} is mentioned explicitly in \ml{E};
    \item \ml{E} contains a pattern-matching lambda abstraction, whose application may fail;
    \item \ml{FAIL} is the value of one of the free variables of \ml{E}.
\end{numbered}

If the pattern-matching compiler described in this chapter is applied throughout, then no pattern-matching lambda abstractions will remain in the transformed program, and hence (ii) cannot occur. Since the programmer presumably cannot write \ml{FAIL} explicitly in his program, it is not hard (although perhaps tedious) to verify that (iii) cannot occur either.

These observations focus our attention on all the places where \ml{FAIL} can be introduced explicitly by the compiler. There are only two such places:
\begin{numbered}
	\item In the translation of conditional equations (Section 4.2.6). Fortunately, we can easily transform conditional equations to avoid the use of \ml{\fatbar{}} and \ml{FAIL}, and we show how to do so below.
	\item In the variant of the pattern-matching compiler described in the last
	section, where the introduction of \ml{\fatbar{}} and \ml{FAIL}, seems unavoidable. This
	problem motivates the discussion in Section 5.5, in which we describe a
	restricted class of function definitions that can always be compiled
	without using \ml{\fatbar{}} and \ml{FAIL}.
\end{numbered}

\subsubsection{Rules for transforming \ml{\fatbar{}} and \ml{FAIL}}
We now give some rules for transforming expressions involving \ml{\fatbar{}} and FAIL to a
simpler form. In all cases their correctness follows directly from the semantics
of \ml{\fatbar{}}.

First, we may eliminate \ml{\fatbar{}} if FAIL cannot occur on the left:
\begin{mlcoded}
E$_1$ \fatbar{} E$_2$ $\equiv$ E$_1$\\
{\normalfont provided that} E$_1$ {\normalfont cannot return} FAIL.
\end{mlcoded}
For example, this rule is used to derive the optimized version of the empty
rule in Section 5.2.4.

Second, we may eliminate \ml{\fatbar{}} if \ml{FAIL} definitely occurs on the right or left:
\begin{mlcoded}
E \ml{\fatbar{}} FAIL $\equiv$ E {\normalfont and}
FAIL \ml{\fatbar{}} E $\equiv$ E
\end{mlcoded}
For example, these rules can be used to simplify the final definition of unwieldy in
Section 5.4.1.

Third, there is the following useful transformation involving \ml{IF}:
\begin{mlcoded}
(IF E$_1$ E$_2$ E$_3$) \fatbar{} E $\equiv$ IF E$_1$ E$_2$ (E$_3$ \fatbar{} E)\\
{\normalfont provided that neither} E$_1$ {\normalfont nor} E$_2$ {\normalfont can return} FAIL.
\end{mlcoded}
This rule will be useful in simplifying conditional equations, which we now
attend to.

\subsubsection{Eliminating \ml{\fatbar{}} and \ml{FAIL} from conditional equations}

The empty rule for \ml{match}, which was described in Section 5.2.4, resulted in an
expression of the form
\begin{mlcoded}
E$_1$ \fatbar{} \ldots{} \fatbar{} E$_m$ \fatbar{} E
\end{mlcoded}

Now, the \ml{E$_i$} are just the right-hand sides of the original equations. If a
right-hand side consisted of a set of guarded alternatives without a final
`otherwise' case, then it will have been translated to the form:
\begin{mlcoded}
IF G$_1$ A$_1$ (IF \ldots{}\,(IF G$_g$ A$_g$ FAIL) \ldots )
\end{mlcoded}
where $g$ is the number of alternatives (see Section 4.2.6). If there was a final
`otherwise' case (that is, a final alternative with no guard, so that the right-
hand side never fails), then it would have been translated to the form:
\begin{mlcoded}
	IF G$_1$ A$_1$ (IF \ldots{}\,(IF G$_{g-1}$ A$_{g-1}$ A$_{g}$) \ldots )
\end{mlcoded}
Notice that \ml{G$_i$} and \ml{A$_r$} cannot be equal to \ml{FAIL}, because they are only the
transformed versions of expressions written by the programmer.

If the right-hand side is of the first form, we can use the third rule of the
previous section repeatedly, followed by the second, to give:
\begin{mlcoded}
(IF G$_1$ A$_1$ (IF \ldots{}\,(IF G$_g$ A$_g$ FAIL) \ldots )) \fatbar{} E\\
\phantom{.}\qquad \qquad \qquad \qquad $\equiv$\\
IF G$_1$ A$_1$ (IF \ldots{}\,(IF G$_g$ A$_g$ E) \ldots )
\end{mlcoded}
If the right-hand side is of the second form, it cannot return \ml{FAIL}, and so we
can use the first rule of the previous section.

Application of these three rules will eliminate all occurrences of \ml{\fatbar{}} and \ml{FAIL}
in the expression generated by the empty rule, and incidentally thereby give a
worthwhile improvement in efficiency.

\subsubsection{Clever compilation}

Using these rules, many of the instances of \ml{\fatbar{}} and \ml{FAIL} remaining in a function
definition can be eliminated. Later we will consider compiling an expression
into low-level machine code. When we do this, we will see that it is possible to
compile the remaining expressions involving \ml{\fatbar{}} and \ml{FAIL} in a surprisingly
efficient way, so that \ml{\fatbar{}} requires no code at all, and the \ml{FAIL} simply compiles to
a jump instruction. This is discussed in Section 20.4.

\section{Uniform Definitions}

This section introduces a restricted class of function definitions, called
\textit{uniform definitions}. There are two motivations for studying this class. First,
uniform definitions avoid certain problems with reasoning about function
definitions that involve pattern-matching. Second, uniform definitions are
easier to compile, and are guaranteed to avoid certain kinds of inefficient
code.

We begin by discussing some problems with reasoning about function
definitions containing pattern-matching. Consider again the alternate
definition of \ml{mappairs}:
\begin{letalign}
mappairs$'$ f [] ys &= [] \\
mappairs$'$ f xs [] &= [] \\
mappairs$'$ f (x:xs) (y:ys) &= f x y : mappairs$'$ f xs ys
\end{letalign}
Now, consider evaluation of the expression:
\begin{mlcoded}
mappairs$'$ (+) bottom []
\end{mlcoded}
where the evaluation of \ml{bottom} would fail to terminate (for example, \ml{bottom} 
could be defined by the degenerate equation \ml{bottom = bottom}). Matching
against the first equation binds \ml{f} to \ml{(+)} and then attempts to match \ml{[]} against
bottom. In order to perform this match it is necessary to evaluate bottom, and
this of course causes the entire expression to fail to terminate.

On the other hand, consider evaluation of:
\begin{mlcoded}
mappairs$'$ (+) [] bottom
\end{mlcoded}
Now matching against the first equation binds \ml{f} to \ml{(+)}, matching \ml{[]} against \ml{[]}
succeeds, and then \ml{ys} is bound to \ml{bottom} (without evaluating \ml{bottom}). So the
expression returns \ml{[]} instead of failing to terminate. This means that the
definition of \ml{mappairs$'$} is not as symmetric as it appears.

Further, if the first two equations of \ml{mappairs$'$} were written in the opposite
order, the two expressions above would change their meaning: now the first
would return \ml{[]} and the second would fail to terminate. So even though the
first and second equations have the same right-hand side, the order in which
they are written is important.

The original definition of \ml{mappairs} has none of these problems:
\begin{letalign}
mappairs f [] ys &= []\\
mappairs f (x:xs) [] &= [] \\
mappairs f (x:xs) (y:ys) &= f x y : mappairs f xs ys
\end{letalign}
Now the asymmetry between
\ml{(mappairs (+) [] bottom)} and
\ml{(mappairs (+) bottom [])} is apparent from the equations. Further, changing
the order of the equations does not change the meaning of the function.

In general, one might expect that whenever the equations do not overlap,
the order in which they are written does not matter. In fact, this is not true.
Consider the definition:
\begin{mlcoded}
\begin{tabular}{llll}
diagonal x &True &False &= 1 \\
diagonal False &y &True &= 2 \\
diagonal True &False &z &= 3
\end{tabular}
\end{mlcoded}
The three equations of this definition are non-overlapping, that is, at most one
equation can apply. However, by this definition, the evaluation of:
\begin{mlcoded}
diagonal bottom True False
\end{mlcoded}
would return \ml{1}. On the other hand, if the order of equations in the definition
were reversed, so the third equation came first, then the above expression
would fail to terminate. So even though the equations do not overlap, the
order in which they are written is important.

Clearly, it would be useful to have a test that guarantees that the order of
the equations does not matter. We now define the class of \textit{uniform} definitions,
which have this property. The definition of `uniformity' is designed so that it is
easy to test whether a definition is uniform while applying the pattern-matching compiler to it.

\simpletitledbox{DEFINITION}{
A set of equations is {\it uniform} if one of the following three conditions holds:
\vspace{-\baselineskip}
\begin{numbered}
	\item either, all equations begin with a variable pattern, and applying the
	variable rule (of Section 5.2.2) yields a new set of equations that is
	also uniform;
	\item or, all equations begin with a constructor pattern, and applying the
	constructor rule (of Section 5.2.3) yields new sets of equations that
	are all also uniform;
	\item or, all equations have an empty list of patterns, so the empty rule (of
	Section 5.2.4) applies, and there is at most one equation in the set.
\end{numbered}
}

That is, a set of equations is uniform if it can be compiled without using the
mixture rule (of Section 5.2.6), and if the empty rule is only applied to sets
containing zero or one equations. (It is easy for the reader to check that when
the empty rule is applied to more than one equation, the order is relevant.)

Such equation sets are called `uniform' because all equations must begin the
same way, either with a variable pattern or a constructor pattern, whereas the
mixture rule applies when some equations begin with variable patterns and
some with constructor patterns.

It is not difficult to prove the following:

\simpletitledbox{THEOREM}{
If a definition is uniform, changing the order of the equations does not
change the meaning of the definition.
}

The proof is a straightforward induction, and is similar in structure to the
proof of correctness of the pattern-matching compiler that was outlined
(along with its definition) in Section 5.2.

This shows that being uniform is a sufficient condition for the order of the
equations not to matter. It is not a necessary condition, as is shown by the
function \ml{dummy}:
\begin{letalign}
dummy [] &= 1\\
dummy xs &= 1, xs = []
\end{letalign}
Clearly, \ml{dummy} is not uniform, but the order of the equations does not matter.
However, the following result shows that being uniform is indeed necessary if
one considers only the left-hand sides:

\simpletitledbox{THEOREM}{
If the left-hand sides of a definition are such that the order of the equations
does not matter (regardless of the right-hand sides or condition parts of
the equations), the definition is uniform.
}

For example, the order of the equations would matter in \ml{dummy} if the \ml{1} in the
second equation were changed to a \ml{2}. Again, the proof of the theorem is a
straightforward induction. These two theorems give us a simpler way of
characterizing uniform equations, without referring to the pattern-matching
compiler. Namely, a definition is uniform if and only if its left-hand sides are
such that the order of the equations does not matter.

It is also possible to show that every uniform definition is non-overlapping.
The converse is not true: the function diagonal is non-overlapping but is not
uniform. Researchers have often referred to `lack of overlapping' as an
important property, but perhaps they should refer to `uniformity' instead,
since this is the property that guarantees that the order of equations does not
matter.

Uniform equations are related to \textit{strongly left-sequential equations} as
defined by Hoffman and O'Donnell [1983], which are in turn related to
\textit{sequential equations} as defined by Huet and Levy [1979].

Notice that although uniform equations are independent of `top-to-bottom'
order, they still have a `left-to-right' bias. For example, although the
following definition is uniform:
\begin{mlcoded}
	\begin{tabular}{llll}
		xor &False &x &= x \\
		xor &True &False &= True \\
		xor &True &True &= False
	\end{tabular}
\end{mlcoded}
the same definition with the arguments interchanged is not:
\begin{mlcoded}
	\begin{tabular}{llll}
		xor$'$ &x &False &= x \\
		xor$'$ &False &True &= True \\
		xor$'$ &True &True &= False
	\end{tabular}
\end{mlcoded}
Of course, we can always get around this bias by using extra definitions to
rearrange the arguments. For example, we can define
\begin{mlcoded}
xor$''$ x y = xor y x
\end{mlcoded}
and then \ml{xor$"$} is equivalent to \ml{xor$'$}, and both \ml{xor$"$} and \ml{xor} have uniform
definitions.

The existence of left-to-right bias is due to the semantics of pattern-
matching that we have chosen. A different definition of pattern-matching that
avoids left-to-right bias is possible; see Huet and Levy [1979].

There is a second reason why uniform equations are important: they
are easier to implement. The problems with implementing non-uniform
definitions have been referred to implicitly in previous sections. In summary,
they are as follows:
\begin{numbered}
\item The resulting case-expressions may examine some variables more than once (see Section 5.2.6).
\item The compiler must use a modified constructor rule to avoid duplicating
the right-hand side of equations (see Section 5.4.1).
\item The resulting expressions may contain \fatbar{} and \ml{FAIL}. Implementing such
expressions efficiently requires additional simplification rules and/or a
special way of implementing \ml{FAIL} using jump instructions (see Section
5.4.2).
\end{numbered}
The result is that the pattern-matching compiler must be significantly more
complicated if it is to deal with non-uniform expressions. Further, the first
point above means that it may be difficult to know how efficient the code
compiled for a non-uniform definition will be.

An issue related to uniformity is the way conditionals are handled. In
languages such as SASL, conditional expressions and \ml{where} expressions may
appear anywhere in an expression, and the semantics of each is defined
independently. In Miranda, conditions and where clauses are not separate
expressions, but rather must be associated with the right-hand side of
definitions. This increases the power of Miranda, in some ways, but only when
non-uniform definitions are used. Hence, a restriction to uniform equations
would also allow this part of the language to be simplified.

On the other hand, it should be pointed out that non-uniform definitions
are sometimes very convenient. For example, the following definition
reverses lists of length two, and leaves all other lists the same:
\begin{letalign}
reverseTwo [x,y] &= [y,x]\\
reverseTwo xs &= xs
\end{letalign}
The most straightforward way of rewriting this as a uniform definition is much
more long-winded:
\begin{letalign}
reverseTwo [] &= []\\
reverseTwo [x] &= [x]\\
reverseTwo [x,y] &= [y,x]\\
reverseTwo (x:y:z:ws) &= x:y:z:ws
\end{letalign}

In this case, it is easy to see another way of rewriting \ml{reverseTwo}, but, in
general, rewriting may not be so easy.

Functional language designers have long debated whether or not
definitions with overlapping equations should be allowed in functional
languages. As has been shown, it may be more appropriate to debate the
merits of uniform -- as opposed to non-overlapping -- equations. Several
arguments in favor of restricting definitions to uniform equations have been
raised here; but it is also true that non-uniform definitions are on occasion
quite convenient. No doubt the debate will continue to be a lively one.

{\footnotesize
\vspace{\baselineskip}
	
{\centering%
	
* \hspace{2cm} * \hspace{2cm} *

}%

\vspace{0.5\baselineskip}

\noindent Acknowledgements: For help with my work on chapters 4, 5 and 7 I would like to
thank the following. Simon Finn made valuable comments, both detailed and deep.
Simon Peyton Jones and I are listed as co-authors only on Chapter 4, but he has
provided as much thoughtful help as a co-author on the other chapters as well. For help
of a very different nature, I thank Catherine Lyons.

This work was performed while on a research fellowship supported by ICL.
}

\section*{References}

\begin{references}
	\item Augustsson, L. 1985. Compiling pattern matching. In \textit{Conference on Functional Programming Languages and Computer Architecture}, Nancy, pp. 368-81.
	\item Jouannaud (editor), LNCS 201. Springer Verlag.
	\item Barrett, G., and Wadler, P. 1986. \textit{Derivation of a Pattern-matching Compiler.} Programming Research Group, Oxford.
	\item Burstall, R.M., MacQueen, D.B., and Sanella, D.T. 1980. HOPE: an experimental applicative language. In \textit{Proceedings of the ACM Lisp Conference}, August.
	\item Hoffmann, C.M., and O'Donnell, M.J. 1983. Implementation of an interpreter for abstract equations. In \textit{10th ACM Symposium on Principles of Programming Languages}, pp. 111-21. ACM.
	\item Huet, G., and Levy, J.J. 1979. \textit{Computations in Non-ambiguous Linear Term Rewriting Systems.} INRIA technical report 359.
\end{references}

\chapter[Transforming the Enriched Lambda Calculus][Transforming the Enriched Lambda Calculus]{Transforming the Enriched\\Lambda Calculus}
\vspace{3cm}

Having now defined the semantics of pattern-matching, we are in a position to
show how to transform all the constructs of the enriched lambda calculus into
the ordinary lambda calculus.

Section 6.1 shows how to transform pattern-matching lambda abstractions
into the ordinary lambda calculus, while Section 6.2 deals with let- and
letrec-expressions; Sections 6.3 and 6.4 deal with \ml{case}-expressions and the \fatbar{}
operator.

\section{Transforming Pattern-matching Lambda Abstractions}

In order to translate Miranda function definitions involving pattern-matching
into the enriched lambda calculus, we had to introduce pattern-matching
lambda abstractions as a new construct. In this section we will show how they
can be transformed into the ordinary lambda calculus. For each form of \ml{(\tlb{p}E)}
we will give an equivalent form that does not use pattern-matching lambda
abstractions.

In the case when the pattern \ml{p} is a variable there is nothing to do, because
no pattern-matching is involved. The remaining cases are when the pattern is
a constant, a product-constructor pattern or a sum-constructor pattern. These
are dealt with in the following three subsections.

\subsection{Constant Patterns}

This section shows how to transform a pattern-matching lambda abstraction
\ml{(\tlb{k}E)}, with a constant pattern \ml{k}, into the ordinary lambda calculus. First of
all, we recall the semantics of \ml{(\tlb{k}E)} from Section 4.3.2:

\begin{letalign}
\metafnbb{Eval}{\tlb{k}E} a = \metafnbb{Eval}{E} \qquad & \text{if} a $=$ \metafnbb{Eval}{k}\\
\metafnbb{Eval}{\tlb{k}E} a = FAIL &\text{if} a $\neq$ \metafnbb{Eval}{k}  \text{and} a $\neq \bot$  \\
\metafnbb{Eval}{\tlb{k}E} $\bot = \bot$
\end{letalign}
Operationally, \ml{(\tlb{k}E)} tests whether its argument is equal to \ml{k}; if so, it returns
\ml{E}, if not it returns \ml{FAIL}. This simple test can be carried out by the built-in \ml{IF}
function, using the following transformation:
\begin{mlcoded}
(\tlb{k}E) $\equiv$ (\tlb{v}   IF (= k v) E FAIL)
\end{mlcoded}
where \ml{v} is a new variable which does not occur free in \ml{E}. It should be clear
(and can be proved, using the semantics of  \ml{(\tlb{k}E)} and the semantics of \ml{IF} and
\ml{=}) that these two lambda abstractions have the same meaning, and hence are
equivalent. Notice the way in which we introduce a new \ml{\tl{v}} abstraction, so that
we can name the argument directly in its body.

As an example, consider the Miranda definition
\begin{mlcoded}
flip 0 = 1\\
flip 1 = 0
\end{mlcoded}
This will be translated to
\begin{mlcoded}
	\begin{tabular}{lll}
		flip = \tlb{x} &( &((\tlb{0} 1) x) \\
		& \fatbar{}  &((\tlb{1} 0) x) \\
		& \fatbar{} &ERROR)
	\end{tabular}
\end{mlcoded}
Now, transforming out the pattern-matching lambda abstractions gives
\begin{mlcoded}
	\begin{tabular}{lll}
		flip = \tlb{x} &( &((\tlb{v} IF (= 0 v) 1 FAIL) x) \\
		& \fatbar{}  &((\tlb{v} IF (= 1 v) 0 FAIL) x) \\
		& \fatbar{} &ERROR)
	\end{tabular}
\end{mlcoded}
It is now easy to verify that
\begin{mlcoded}
\begin{tabular}{llllll}
flip 0 & $\to$ & $\cdots$ & $\to$ & 1 \\
flip 1 & $\to$ & $\cdots$ & $\to$ & 0 \\
flip 2 & $\to$ & $\cdots$ & $\to$ & ERROR \\
\end{tabular}
\end{mlcoded}

\subsection{Product-constructor Patterns}

Next we consider the case of \ml{(\tlb{p}E)}, where \ml{p} is the product pattern
\ml{(t p$_1$ \ldots p$_r$)}, and \ml{t} is a product constructor of arity $r$. As before, we recall its
semantics (Section 4.3.4):
\begin{letalign}
	\metafnbb{Eval}{\tlb{(t p$_1$ $\ldots$ p$_r$)}E} a = \metafnbb{Eval}{ \tl{p$_1$} $\ldots$ \tlb{p$_r$}E } & (SEL-t-1 a)\\
	& $\cdots$\\
	& (SEL-t-r a)
\end{letalign}

To implement this semantics, we invent a new function
\ml{UNPACK-PRODUCT-t} for each product constructor \ml{t}, and use it in this
transformation:
\begin{mlcoded}
(\tlb{(t p$_1$ $\ldots$ p$_r$)}E) $\equiv$ UNPACK-PRODUCT-t (\tl{p$_1$}$\ldots$\tlb{p$_r$}E)
\end{mlcoded}
The idea is that \ml{UNPACK-PRODUCT-t} takes two arguments, a function and a
structured object, and applies the function to the lazily selected components
of the object. It is defined by the following semantic equation:
\begin{mlcoded}
UNPACK-PRODUCT-t f a $=$ f (SEL-t-1 a) $\cdots$ (SEL-t-r a)
\end{mlcoded}
It can easily be shown that the transformation is valid, by comparing the 
semantics of the expression before and after the transformation.

The right-hand side of the transformation still has pattern-matching lambda
abstractions in it, but they are smaller than the one we began with, and
repeated use of the rules for transforming pattern-matching lambda 
abstractions will eliminate them.

As an example, consider the function \ml{addPair}, which adds together the
elements of a pair:
\begin{mlcoded}
	addPair = \tlb{PAIR x y} + x y
\end{mlcoded}
This will be transformed to
\begin{mlcoded}
	addPair = UNPACK-PRODUCT-PAIR (\tlb{x}\tlb{y}+ x y)
\end{mlcoded}
We can check that it gives the right results by reducing \ml{(addPair (PAIR 3 4))}:
\begin{mlcoded}
	addPair (PAIR 3 4)\\
	\phantom{--}= UNPACK-PRODUCT-PAIR (\tlb{x}\tlb{y}+ x y) (PAIR 3 4) \\
	$\rightarrow$ (\tlb{x}\tlb{y}+ x y) (SEL-PAIR-1 (PAIR 3 4)) (SEL-PAIR-2 (PAIR 3 4)) \\ 
	$\rightarrow$ (\tlb{y}+ (SEL-PAIR-1 (PAIR 3 4)) y) (SEL-PAIR-2 (PAIR 3 4)) \\
	$\rightarrow$ + (SEL-PAIR-1 (PAIR 3 4)) (SEL-PAIR-2 (PAIR 3 4)) \\
	$\rightarrow$ + 3 (SEL-PAIR-2 (PAIR 3 4)) \\  
	$\rightarrow$ + 3 4 \\
	$\rightarrow$ 7
\end{mlcoded}

\subsection{6.1.3 Sum-constructor Patterns}

Finally, consider the case of \ml{(\tlb{p}E)}, where \ml{p} is a sum pattern \ml{(s p$_1$ $\cdots$ p$_r$)}, and
\ml{s} is a sum constructor of arity \ml{r}. The semantics of such lambda abstractions
was derived in Section 4.3.3:
\begin{letalign}
	\metafnbb{Eval}{\tlb{(s p$_1$\,\ldots\,p$_r$)}E} (s a$_1$\,\ldots\,a$_r$) 
	&= \metafnbb{Eval}{\tl{p$_1$}\ldots\tlb{p$_r$}E} a$_1$\,\ldots\,a$_r$ \\
	\metafnbb{Eval}{\tlb{(s p$_1$\,\ldots\,p$_r$)}E} (s$'$ a$_1$\,\ldots\,a$_{r'}$) 
	&= FAIL \qquad {\normalfont if} s$\,\neq\,$s$'$ \\
	\metafnbb{Eval}{\tlb{(s p$_1$\,\ldots\,p$_r$)}E} $\bot$
	&= $\bot$ \\
\end{letalign}


We can make a very similar transformation to the product-constructor case,
leaving all the hard work to a new function \ml{UNPACK-SUM-s}:
\begin{mlcoded}
	(\tlb{(s p$_1$\,\ldots\,p$_r$)}E) $=$ UNPACK-SUM-s (\tl{p$_1$}\,\ldots\,\tlb{p$_r$}E)
\end{mlcoded}

The function \ml{UNPACK-SUM-s} takes two arguments, a function (in this case
\ml{(\tl{p$_1$} \ldots \tlb{p$_r$}. E)}), and a structured object. It checks whether the object is built
with constructor \ml{s}: if not, \ml{FAIL} is returned; if so, \ml{UNPACK-SUM-s} takes the
object apart and applies the function (its first argument) to its components.
\ml{UNPACK-SUM-s} is specified by the following semantic equations:
\begin{letalign}
	UNPACK-SUM-s f (s a$_1$\,\ldots\,a$_r$) &= f\,a$_1$\,\ldots\,a$_r$ \\
	UNPACK-SUM-s f (s$'$ a$_1$\,\ldots\,a$_{r'}$) &= FAIL \qquad {\normalfont if} s\,$\neq$\,s$'$\\
	UNPACK-SUM-s f $\bot$ &= $\bot$
\end{letalign}

As an example, recall the Miranda definition of \ml{reflect}:
\begin{letalign}
	reflect (LEAF n) &= LEAF n\\
	reflect (BRANCH t$_1$ t$_2$) &= BRANCH (reflect t$_2$) (reflect t$_1$)
\end{letalign}
This is translated to:
\begin{mlalign}
	reflect = \tlb{t}&( ((\tlb{(LEAF\ n)}LEAF\ n) t)\\
	&\fatbar{} ((\tlb{(BRANCH t$_1$\ t$_2$)}BRANCH (reflect\ t$_2$) (reflect\ t$_1$)) t)\\
	&\fatbar{}  ERROR)
\end{mlalign}
Now, applying the transformation gives:
\begin{mlalign}
	reflect\\
	= \tl{t}. &( (UNPACK-SUM-LEAF (\tl{n}. LEAF\ n) t)\\
	&\fatbar{}  (UNPACK-SUM-BRANCH (\tlb{t$_1$}\tlb{t$_2$}BRANCH (reflect\ t$_2$) (reflect\ t$_1$)) t)\\
	&\fatbar{}  ERROR)
\end{mlalign}

\subsection{Reducing the Number of Built-in Functions}

The trouble with the transformations of the previous section is that they
introduce several functions associated with each constructor. In this section,
we discuss the implementation of these functions.

A structured object will be represented by the implementation as an
aggregate, consisting of the component fields together with a \textit{structure tag},
which distinguishes objects built by different constructors from each other
(see Section \ml{10.3.1}). It is this tag which can be used by \ml{UNPACK-SUM-s} to
identify the constructor used.

In a type-checked system, it is only necessary to distinguish objects from
other objects of the same type, so the structure tag can be a small integer in the
range \ml{1} \ldots \ml{n} (where \ml{n} is the number of constructors in the type). This means
that, instead of requiring an \ml{UNPACK-SUM-s} function for each constructor \ml{s},
it is only necessary to have a single family of functions \ml{UNPACK-SUM-d-r$_s$},
where \ml{d} is the integer structure tag which is recognized by \ml{UNPACK-SUM-d-r$_s$},
and \ml{r$_s$} is the arity of \ml{s}. In a similar way, the sum constructor functions can be
replaced with a family of functions \ml{PACK-SUM-d-r$_s$}, which take \ml{r$_s$} arguments
and construct an aggregate with \ml{r$_s$} fields and structure tag \ml{d}.

We can perform an analogous set of replacements for the functions
associated with product types. \ml{UNPACK-PRODUCT-t} can be replaced with
\ml{UNPACK-PRODUCT-r$_t$}, where \ml{r} is the arity of \ml{t} (there is no need for a structure
tag here, since \ml{UNPACK-PRODUCT} does not examine it). Similarly, the
product-constructor functions can be replaced with \ml{PACK-PRODUCT-r$_t$}, and
the selector functions \ml{SEL-t-i} can be replaced with \ml{SEL-r$_{t}$-i}. It is sensible to
keep \ml{PACK-SUM} and \ml{PACK-PRODUCT} distinct because, having no structure
tag, objects of product type may have a different representation from objects
of sum type.

To summarize:
\vspace{0.5\baselineskip}

{\setlength{\tabcolsep}{1pt}
\begin{tabular}{rl}
	\ml{s} (a sum-constructor function) &is replaced by \ml{PACK-SUM-d-r$_s$}\\
	\ml{UNPACK-SUM-s} &is replaced by \ml{UNPACK-SUM-d-r$_s$}\\
	\ml{t} (a product-constructor function) &is replaced by \ml{PACK-PRODUCT-r$_t$}\\
	\ml{UNPACK-PRODUCT-t} &is replaced by \ml{UNPACK-PRODUCT-r$_t$}\\
	\ml{SEL-t-i} &is replaced by \ml{SEL-r$_{t}$-i}
\end{tabular}

\begin{tabular}{rl}
where \ml{r$_s$} &$=$ arity of \ml{s},\\
\ml{d} &$=$ structure tag of \ml{s},\\
\ml{r$_t$} &$=$ arity of \ml{t}.\\
\end{tabular}
}
\vspace{0.5\baselineskip}

For example, assuming that we implement lists with structure tag $1$ for \ml{NIL}
and $2$ for \ml{CONS}, then the following replacements would take place:
\vspace{0.5\baselineskip}

{\setlength{\tabcolsep}{1pt}
	\begin{tabular}{rl}
	\ml{NIL} &is replaced by \ml{PACK-SUM-1-0}\\
	\ml{CONS} &is replaced by \ml{PACK-SUM-2-2}\\
	\ml{UNPACK-SUM-NIL} &is replaced by \ml{UNPACK-SUM-1-0}\\
	\ml{UNPACK-SUM-CONS} &is replaced by \ml{UNPACK-SUM-2-2}
\end{tabular}
}
\vspace{0.5\baselineskip}

Likewise, if the type \ml{tree} is declared as before:
\begin{mlcoded}
	tree ::= LEAF num | BRANCH tree tree
\end{mlcoded}
and \ml{LEAF} and \ml{BRANCH} are assigned structure tags 1 and 2 respectively, the
following replacements would take place:
\vspace{0.5\baselineskip}

{\setlength{\tabcolsep}{1pt}
\begin{tabular}{rl}
	\ml{LEAF} &is replaced by \ml{PACK-SUM-1-1}\\
	\ml{BRANCH} &is replaced by \ml{PACK-SUM-2-2}\\
	\ml{UNPACK-SUM-LEAF} &is replaced by \ml{UNPACK-SUM-1-1}\\
	\ml{UNPACK-SUM-BRANCH} &is replaced by \ml{UNPACK-SUM-2-2}
\end{tabular}
}\vspace{0.5\baselineskip}

Finally, if the type \ml{pair} is declared as before:
\begin{mlcoded}
	pair $*$ $**$ ::= PAIR $*$ $**$
\end{mlcoded}
the following replacements would take place:
\vspace{0.5\baselineskip}

{\setlength{\tabcolsep}{1pt}
	\begin{tabular}{rl}
	\ml{PAIR} &is replaced by \ml{PACK-PRODUCT-2}\\
	\ml{UNPACK-PRODUCT-PAIR} &is replaced by \ml{UNPACK-PRODUCT-2}\\
	\ml{SEL-PAIR-1} &is replaced by \ml{SEL-2-1}\\
	\ml{SEL-PAIR-2} &is replaced by \ml{SEL-2-2}
\end{tabular}
\vspace{0.5\baselineskip}

Since functions with different types may be replaced by the same function
(for example, \ml{CONS} and \ml{BRANCH} are both replaced by \ml{PACK-SUM-2-2}), these
replacements should not be performed until after type-checking. For the
same reason, none of these replacements is possible for a system that
performs run-time type-checking (see Section 10.5).

\subsection{Summary}
Figure 6.1 summarizes the transformations developed in this section, and
Figure 6.2 gives the semantics for the two new families of functions we
introduced in order to perform the transformations.


\boxedfigure{
	\begin{mlcoded}
	\begin{tabular}{lll}
			(\tlb{k}E)	& $\equiv$  (\tlb{v} IF (= k v) E FAIL) \\
			&  \phantom{ $\equiv$  } {\normalfont where} v {\normalfont is a new variable that does not } \\
			&	\phantom{ $\equiv$  } \qquad{\normalfont occur free in} E\\
			(\tlb{(t p$_1$ \ldots p$_{r_t}$)E}) & $\equiv$ (UNPACK-PRODUCT-t (\tlb{p$_1$ \ldots p$_{r_t}$ E})) \\
			(\tlb{(s p$_1$ \ldots p$_{r_s}$) E}) & $\equiv$ (UNPACK-SUM-s (\tlb{p$_1$ \ldots p$_{r_s}$ E}))
	\end{tabular}
	\end{mlcoded}
	
		\begin{tabular}{lll}
			where	& \ml{k}  & is a constant \\
			& \ml{t} & is a product constructor of arity \ml{r$_t$}\\
			& \ml{s} & is a sum constructor of arity \ml{r$_s$}
		\end{tabular}
}{Transforming out pattern-matching lambda abstractions}


\boxedfigure{
	\begin{mlcoded}
		UNPACK-PRODUCT-t f a = f (SEL-t-1 a) \ldots (SEL-t-r$_t$ a) \\
		
		\begin{tabular}{lll}
			UNPACK-SUM-s f (s a$_1$ \ldots a$_{r_s}$) &= f a$_1$ \ldots a$_{r_s}$ \\
			UNPACK-SUM-s f (s' a$_1$ \ldots a$_{r_{s'}}$) &= FAIL  \qquad {\normalfont if} s\,$\neq$\,s$'$\\
			UNPACK-SUM-s f $\bot$ &=$\bot$
		\end{tabular}
	\end{mlcoded}
	
	\begin{tabular}{lll}
		where & \ml{t} & is a product constructor of arity \ml{r$_t$}\\
		& \ml{s} & is a sum constructor of arity \ml{r$_s$}
	\end{tabular}
}{Semantics of \ml{UNPACK-PRODUCT} and \ml{UNPACK-SUM}}


\section{Transforming \ml{let} and \ml{letrec}}

In Section 4.2.9 we introduced a new complication to \ml{let}(\ml{rec})-expressions, by
allowing the left-hand side of definitions to be an arbitrary pattern rather than
a simple variable. In this section we show how to transform these generalized
\ml{let}s and \ml{letrec}s into successively simpler forms, arriving eventually at the
ordinary lambda calculus.

Rather than defining the semantics of \ml{let} and \ml{letrec} directly, as we did for
pattern-matching lambda abstractions, we will regard the transformations
described in this section as a definition of their semantics. To define their
meaning in a more direct way would require more mathematical machinery
than we have available in this book.

We begin by sketching a new problem which is introduced by allowing
arbitrary patterns on the left-hand side of definitions. This leads us to define a
class of patterns, the \textit{irrefutable} patterns, which do not suffer from the
problem. Then, before embarking on the transformations themselves, we
give a 'map' to explain their structure.

\subsection{Conformality Checking and Irrefutable Patterns}

Allowing arbitrary patterns on the left-hand side of a definition introduces a
new and somewhat subtle complication. Consider the expression
\begin{mlcoded}
	let (CONS x xs) = B in E
\end{mlcoded}
Here, the pattern \ml{(CONS x xs)} appears on the left-hand side of the definition.
This raises the nasty possibility that B might evaluate to \ml{NIL} instead of
\ml{(CONS B$_1$ B$_2$)}, in which case the pattern would not match, and some sort of
error should, presumably, be reported. This requires that a \textit{conformality
check} be made, to ensure that B conforms with the specified pattern.
Conformality checking will carry some implementation cost, so we would
like to avoid it whenever possible. It can be avoided in precisely those cases
when the pattern match cannot fail; for example, simple product patterns.
However, there are some nested patterns which cannot fail also, which
motivates the following definition:


\simpletitledbox{DEFINITION}{
A pattern \ml{p} is \textit{irrefutable} if it is
\begin{numbered}
	\item either a variable \ml{v}
	\item or a product pattern of form \ml{(t p$_1$ \ldots p$_{r_t}$)} where \ml{p$_1$}, \ldots, \ml{p$_{r_t}$} are irrefutable patterns.
\end{numbered}
Otherwise the pattern is \textit{refutable}.
}

\noindent In other words, the irrefutable patterns consist of arbitrarily nested product
constructors with variables at the leaves. These patterns cannot fail to match
in a type-checked implementation. Variables and simple product patterns are
just two examples of irrefutable patterns.

However, even a single constant or sum constructor (even if nested inside a
product pattern) makes the pattern refutable, since there is a possibility that it
may not match. We need to perform conformality checking for refutable
definitions only.

\subsection{Overview of let and letrec Transformations}
We are now ready to describe the various transformations to simplify \ml{let}(\ml{rec})-
expressions. While few are complicated, they are quite numerous, so we
begin by offering a `map' to aid in navigation through the rest of the section.

For a start, we establish the following terminology:
\begin{numbered}
	\item The left-hand side of each definition of a simple \ml{let(rec)}-expression must
	be a variable.
	\item The left-hand side of each definition of an irrefutable \ml{let(rec)}-expression
	must be an irrefutable pattern.
	\item The left-hand side of each definition of a general \ml{let(rec)}-expression may
	be any arbitrary pattern.
\end{numbered}
With the aid of this terminology, Figure 6.3 depicts the transformations which
will be described below, giving the appropriate section number in brackets.

For the reasons discussed in Section 3.2.4, there are two possible forms into
which we may wish to transform the program, which differ only in their
treatment of \ml{let} and \ml{letrec}:
\begin{numbered}
	\item We may transform the program into the ordinary lambda calculus; this
	gives the simplest resulting program. In this case, general \ml{let}s are trans-
	formed into the ordinary calculus via irrefutable \ml{let}s and simple \ml{let}s.
	General \ml{letrec}s, on the other hand, are first transformed into irrefutable
	\ml{let}s via irrefutable \ml{letrec}s, and then use the \ml{let} transformations.
	\item We may transform the program into the ordinary lambda calculus
	augmented with simple \ml{let(rec)}-expressions; the resulting program is
	slightly more complicated, but can be implemented more efficiently
	(Section 3.2.4). In this case, general \ml{let}s are transformed only into simple
	\ml{let}s, and general \ml{letrec}s are transformed into simple \ml{letrec}s, via irrefutable
	\ml{letrec}s.
\end{numbered}

\begin{figure}[H]
	\centering
	
	\fbox{%
		\begin{minipage}{0.9\textwidth}
			\small
			\setlength{\parindent}{10pt}
			\setlength{\parskip}{0mm plus 0mm minus 0mm}
			
			\includegraphics[width=0.96\textwidth]{chapters/Fig6_3}
			%                    \includesvg{chapters/Fig3_4}
			%                    \input{chapters/Fig3_4.tex}
			
		\end{minipage}%
		
	}%
	
	\caption{\textsf Map of \ml{let(rec)} transformations}
\end{figure}

\noindent Both possibilities are catered for by the transformations shown in Figure 6.3.

In what follows, when considering \ml{let}-expressions we assume that they
contain only one definition. This gives no loss of generality, since a \ml{let}-expression with multiple definitions is trivially equivalent to a nested set of
single-definition \ml{let}-expressions.

The following sections deal with the transformations depicted in Figure 6.3.

\subsection{Transforming Simple \ml{let}s into the Ordinary Lambda Calculus}

Once we have arrived at an expression in which all \ml{let}-expressions are simple,
it is easy to remove them altogether, using the transformation given in Section
3.2.1:

\plainbox{
\begin{mlcoded}
	let v = B in E $\equiv$ (\tlb{v}E) B
\end{mlcoded}
}
\noindent For example,
\begin{mlcoded}
	let x = 4 in (+ x 6) $\equiv$ (\tlb{x}+ x 6) 4
\end{mlcoded}

\subsection{Transforming Irrefutable \ml{let}s into Simple \ml{let}s}

Consider the case of an irrefutable \ml{let}-expression, of the form
\begin{mlcoded}
	let p = B in E
\end{mlcoded}
where \ml{p} is irrefutable. Since the pattern on the left-hand side of the definition
is irrefutable, it must either be a variable or a product pattern. In the former
case, there is nothing to do, since the \ml{let}-expression is already simple. In the
latter case, the \ml{let}-expression takes the form
\begin{mlcoded}
	let (t p$_1$ \ldots p$_r$) = B in E
\end{mlcoded}
where the \ml{p$_i$} are irrefutable patterns, and \ml{B} and \ml{E} are expressions. We can now
make the following transformation:

\plainbox{
\begin{mlcoded}
	\begin{tabular}{llll}
	let (t p$_1$ \ldots p$_r$) = B in E $\equiv$ \ &let &v =&\ B \\
		&in &(let &p$_1$ = SEL-t-1 v \\
		& & &\ldots \\
		& & &p$_r$ = SEL-t-r v \\
		& & \;\;in &E)
	\end{tabular}
\end{mlcoded}
where \ml{v} is a new variable that does not occur free in \ml{E}.
}

The \ml{p$_i$} are bound to selector functions applied to \ml{v}, which is in turn bound to
\ml{B}. Repeated application of this transformation will eliminate all non-simple
irrefutable \ml{let}-expressions.

To take an example, the expression
\begin{mlcoded}
	let (PAIR x y) = B in E
\end{mlcoded}
would be transformed to
\begin{mlcoded}
	\begin{tabular}{lll}
	let v = B in (&let &x = SEL-PAIR-1 v \\
		& &y = SEL-PAIR-2 v\\
		&in &E)
	\end{tabular}
\end{mlcoded}

Notice that if neither \ml{x} nor \ml{y} is evaluated in \ml{E}, then \ml{B} will not be evaluated
either, so the transformation implements lazy product-matching. Lazy
product-matching is just as much of an advantage here as it was in function
definitions. For example, we could recode the function \ml{firsts} from Section
4.3.5 in the following way:
\begin{mlcoded}
	\begin{tabular}{lll}
	firsts [\,] &= &(0, 0)\\
	firsts (x:xs) &= &(x, ev), \qquad odd x\\
	&= &(od, x), \qquad even x\\
	& &where\\
	& &(od, ev) = firsts xs
	\end{tabular}
\end{mlcoded}
We would expect this definition to behave just like that of Chapter 4, so that if
lazy product-matching is used for function definitions then it should also be
used for \ml{let}(\ml{rec})-expressions.

(Note: an alternative transformation would have been possible in this
section, namely:
\begin{mlcoded}
	let p = B in E = (\tlb{p}E) B
\end{mlcoded}
where \ml{p} is an irrefutable pattern. From a semantic point of view, this is
entirely equivalent to the transformation used above. However, for the
efficiency reasons outlined in Section 3.2.4, we prefer to stay in the world of
\ml{let}-expressions as long as possible; hence our choice.)

\subsection{Transforming Irrefutable \ml{letrec}s into Simple \ml{letrec}s}

The transformation from a letrec involving only irrefutable definitions into a
simple \ml{letrec} is very similar to that for \ml{let}-expressions:

\plainbox{
\begin{mlcoded}
	\begin{tabular}{llll}
	letrec &(t p$_1$ \ldots p$_r$) = B $\equiv$ &letrec &v = B \\
	&<other definitions> & & p$_1$ = SEL-t-1 v \\
	in E & & & $\cdots$\\
	&  & &	p$_r$ = SEL-t-r v \\
	& &	& <other definitions>\\
	& & in E & 
	\end{tabular}
\end{mlcoded}
where \ml{v} is a new variable that does not occur free in \ml{E} or \ml{B}.
}


All the transformed definitions must be in a single \ml{letrec}, to ensure that
variables in the patterns \ml{p$_i$} are in scope in \ml{B}. The `\ml{<other definitions>}' simply
takes into account the fact that the \ml{letrec} may contain multiple definitions, and
this transformation should be applied to each of them separately.

%\backmatter

%\include{appendix}

\end{document}
